\documentclass[11pt]{article}
\usepackage[paperwidth=220mm,paperheight=320mm,left=16mm,right=16mm,top=18mm,bottom=18mm,headheight=16pt,headsep=10pt,footskip=14pt]{geometry}
\usepackage{fix-cm}
\usepackage{fontspec,amsmath,amssymb,graphicx,array,multirow}
\usepackage[unicode,hidelinks]{hyperref}
\usepackage{polyglossia}
\setmainlanguage{latin}
\setotherlanguages{arabic,greek,hebrew,syriac}
\setmainfont{Linux Libertine O}[Ligatures=TeX]
\setsansfont{Linux Biolinum O}[Ligatures=TeX]
\newfontfamily\greekfont{FreeSerif}[Script=Greek]
\newfontfamily\arabicfont{Amiri}[Script=Arabic]
\newfontfamily\hebrewfont{Frank Ruhl Hofshi}[Script=Hebrew]
\newfontfamily\syriacfont{Segoe UI Historic}[Script=Syriac]
\newfontfamily\CapFont{Noto Serif}
\newfontfamily\MoonFont{FreeSerif}
\newfontfamily\EthiopicFont{Ebrima}
\newcommand{\textethiopic}[1]{{\EthiopicFont #1}}
\newfontfamily\NallinoSigns{NallinoSigns.otf}[Path=./fonts/]
\newcommand{\AbjadZero}{{\NallinoSigns\symbol{"E000}}}
\hypersetup{pdftitle={Al-Battani, Opus astronomicum. Nallino, Adnotationes ad tabulas (Pars II, pp. 189-317). S07},pdfauthor={Al-Battani; Carlo Alfonso Nallino (notes)}}
\makeatletter
\def\ps@sourceedition{%
\def\@oddhead{\vbox{\hbox to\textwidth{\fontsize{8}{10}\selectfont C. A. NALLINO — ADNOTATIONES AD TABULAS\hfil\leftmark}\vskip3pt\hrule height.25pt}}%
\let\@evenhead\@oddhead
\def\@oddfoot{\hbox to\textwidth{\fontsize{7}{9}\selectfont\rightmark\hfil\thepage}}%
\let\@evenfoot\@oddfoot}
\makeatother
\pagestyle{sourceedition}
\setlength{\parindent}{0pt}\setlength{\parskip}{0pt}
\newcommand{\BeginSource}[2]{\clearpage\markboth{PAG. #2}{#1}\hypertarget{#1}{}\pdfbookmark[0]{#1 — p. #2}{bk-#1}\fontsize{11.1}{13.8}\selectfont}
\tracinglostchars=2
\newcommand{\Name}[1]{{\addfontfeatures{LetterSpace=4.0}#1}}
\newcommand{\Indent}{\hspace*{1.6em}}
\newcommand{\qfrac}[2]{{}^{#1}\!/_{{#2}}}
\newcommand{\sa}{\textsuperscript{\textit{a}}}\newcommand{\sm}{\textsuperscript{\textit{m}}}\newcommand{\sd}{\textsuperscript{\textit{d}}}
\newcommand{\RawBlock}[1]{\par\vspace{3pt}\noindent#1\par\vspace{3pt}}
\newcommand{\CenterLine}[1]{\makebox[\linewidth][c]{#1}}
\newlength{\FullSourceWidth}\newlength{\GutterOffset}\newif\ifNumbersLeft
\newcommand{\DSource}[2]{\BeginSource{AB01-PDF#1}{#2}%
 \setlength{\FullSourceWidth}{\textwidth}\setlength{\GutterOffset}{0pt}%
 \ifodd#2\relax\NumbersLefttrue\else\NumbersLeftfalse\fi}
\newsavebox{\DLBox}
\newcommand{\DLine}[3]{%
 \par\noindent\hypertarget{#1}{}%
 \ifx\relax#2\relax\else
   \ifNumbersLeft
    \rlap{\kern-\GutterOffset\llap{\fontsize{7}{8}\selectfont #2\hspace{3mm}}}%
   \else
    \rlap{\kern\dimexpr\FullSourceWidth-\GutterOffset\relax\hspace{3mm}{\fontsize{7}{8}\selectfont #2}}%
   \fi
 \fi
 \sbox{\DLBox}{\strut #3}%
 \ifdim\wd\DLBox>\linewidth\typeout{DIPLOMATICWIDTH|#1|\the\wd\DLBox|\the\linewidth}\fi
 \usebox{\DLBox}\strut\par}
\begin{document}
\thispagestyle{empty}
\begin{center}\Large AL-BATTĀNĪ\par\vspace{4mm}\LARGE OPUS ASTRONOMICUM\par
\vspace{5mm}\large Caroli Alphonsi Nallino adnotationes ad tabulas (pars II)\par\vspace{8mm}
\Large S07 --- editio diplomatica in progressu\par\vspace{4mm}\large Paginae impressae 189--216\end{center}
\vspace{12mm}\noindent Singulae lineae paginarum impressarum singulis lineis huius editionis respondent, cum
divisionibus vocabulorum, signis et generibus litterarum (rectis, inclinatis, distantibus). Textus ex imaginibus
paginarum transcriptus est; litterae Graecae et Arabicae, numeri et signa in imaginibus amplificatis lecta sunt.
Haec transcriptio a nullo homine recognita est.\par
\DSource{0638}{189}
\par\vspace{120mm}
\DLine{AB01-PDF0638-body-L001}{}{\CenterLine{\fontsize{24}{28}\selectfont ADNOTATIONES}}
\par\noindent\makebox[\linewidth][c]{\rule[.5ex]{18mm}{.4pt}}\par
\DSource{0639}{190}
\null
\DSource{0640}{191}
\par\vspace{10mm}
\par\noindent\rule{\linewidth}{1.2pt}\par\nointerlineskip\vspace{1pt}\noindent\rule{\linewidth}{.4pt}\par
\par\vspace{38mm}
\DLine{AB01-PDF0640-body-L001}{}{\CenterLine{\fontsize{24}{28}\selectfont ADNOTATIONES}}
\par\noindent\makebox[\linewidth][c]{\rule[.5ex]{18mm}{.4pt}}\par
\par\vspace{26mm}
\hypertarget{AB01-PDF0640-H02}{}
\DLine{AB01-PDF0640-body-L002}{}{\CenterLine{\textit{Ad pag. 1–3.} — \textbf{Tabulae regum peregrinorum.}}}
\par\vspace{4mm}
\hypertarget{AB01-PDF0640-P01}{}
\DLine{AB01-PDF0640-body-L003}{}{\Indent In titulo (1) lectionem codicis servavi; censeo tamen al-Battānī \textarabic{تاريخ} pro \textarabic{بتاريخ} scripsisse,}
\DLine{AB01-PDF0640-body-L004}{}{scilicet « et ab eo [incipit] chronologia Almagesti », sicut etiam apud \Name{al-Bīrūnī}, \textit{Chron.}}
\DLine{AB01-PDF0640-body-L005}{5}{p. 88 text., 101 vers., scriptum invenimus. Re vera in \textit{Almagesto} \Name{Ptolemaeus} semper fere}
\DLine{AB01-PDF0640-body-L006}{}{annis aerae Nabonassaris utitur.}
\hypertarget{AB01-PDF0640-P02}{}
\DLine{AB01-PDF0640-body-L007}{}{\Indent Graece et Franco-Gallice exstant hae tabulae, a Ptolemaeo usque ad Antoninum mor-}
\DLine{AB01-PDF0640-body-L008}{}{tuum, ab ignoto Graeco usque ad Constantinopolim a Turcis captam perductae, in libro \textit{Chro-}}
\DLine{AB01-PDF0640-body-L009}{}{\textit{nologie de} \Name{Ptolémée}: \textit{table chronologique des règnes, apparition des étoiles fixes de}}
\DLine{AB01-PDF0640-body-L010}{10}{C. \Name{Ptolémée}, \Name{Théon} etc. [\textit{publiées et}] \textit{traduites par Halma,} Paris 1819, 2\textsuperscript{e} partie, p. 3–5.}
\DLine{AB01-PDF0640-body-L011}{}{Usque ad Diocletianum regem optime tabulas edidit etiam \Name{Ideler}, \textit{Handbuch der mathe-}}
\DLine{AB01-PDF0640-body-L012}{}{\textit{matischen und technischen Chronologie,} Berlin 1825–26, t. I, p. 111–14; regum Assyriae}
\DLine{AB01-PDF0640-body-L013}{}{(quorum ultimus Nabonadius) illustravit Eb. \Name{Schrader}, \textit{Die keilinschriftliche babylonische}}
\DLine{AB01-PDF0640-body-L014}{}{\textit{Königsliste,} in Sitzungsberichte der k. preussischen Akad. der Wissenschaften zu Berlin, 1887,}
\DLine{AB01-PDF0640-body-L015}{15}{nr. XXXI, p. 579–607.}
\hypertarget{AB01-PDF0640-P03}{}
\DLine{AB01-PDF0640-body-L016}{}{\Indent Apud \Name{Ptolemaeum} Graecumque eius continuatorem ita reges distinguuntur: \textit{Assyrio-}}
\DLine{AB01-PDF0640-body-L017}{}{\textit{rum et Medorum,} \textgreek{Ἀσσυρίων καὶ Μηδῶν βασιλεῖς}, usque ad Nabonadium mortuum; \textit{Persarum,}}
\DLine{AB01-PDF0640-body-L018}{}{\textgreek{Περσῶν βασιλεῖς}, a Cyro ad Alexandrum Macedonem; \textit{anni regum Macedonum,} \textgreek{ἔτη βασιλέων}}
\DLine{AB01-PDF0640-body-L019}{}{\textgreek{Μακεδόνων}, a Philippo ad Cleopatram; \textit{regum Romanorum,} \textgreek{Ῥωμαίων βασιλέων}, ab Augusto ad}
\DLine{AB01-PDF0640-body-L020}{20}{quartum Constantini annum absolutum; \textit{regum Christianorum Romanorum,} \textgreek{Χριστιανῶν Ῥω}-}
\DLine{AB01-PDF0640-body-L021}{}{\textgreek{μαίων βασιλέων}, inde ab initio quinti anni Constantini.}
\hypertarget{AB01-PDF0640-P04}{}
\DLine{AB01-PDF0640-body-L022}{}{\Indent Eadem tabula Albateniana, usque ad Basilium imperatorem porrecta, apud \Name{al-Bīrūnī},}
\DLine{AB01-PDF0640-body-L023}{}{\textit{Chron.,} p. 88–89 et 92–96 text., 101 et 103–105 vers.}
\hypertarget{AB01-PDF0640-P05}{}
\DLine{AB01-PDF0640-body-L024}{}{\Indent Ex aliis fontibus manant tabulae regum ab Augusto ad Heraclium apud \Name{aṭ-Ṭabarī},}
\DLine{AB01-PDF0640-body-L025}{25}{\textit{Annales,} ser. I, t. II, p. 741–44; \Name{Ibn al-Athīr}, \textit{Annales,} Kahirae 1301–02 heg., t. I,}
\par\vspace{6pt}\hrule height.25pt\vspace{5pt}
\noindent\begin{minipage}[t]{.48\textwidth}\vspace{0pt}\fontsize{9}{11.5}\selectfont\setlength{\GutterOffset}{0pt}
\hypertarget{AB01-PDF0640-N01}{}
\DLine{AB01-PDF0640-notes_left-L001}{}{\Indent (1) Textus Arabicus harum tabularum est t. III,}
\DLine{AB01-PDF0640-notes_left-L002}{}{p. 228-31. In versione Hispanica hae tabulae omit-}
\end{minipage}
\hfill\begin{minipage}[t]{.48\textwidth}\vspace{0pt}\fontsize{9}{11.5}\selectfont\setlength{\GutterOffset}{0pt}
\DLine{AB01-PDF0640-notes_right-L001}{}{tuntur, praeter (f. 29,r.) conspectum intervallorum,}
\DLine{AB01-PDF0640-notes_right-L002}{}{quem infra, p. 197, edam.}
\end{minipage}
\DSource{0641}{192}
\DLine{AB01-PDF0641-body-L001}{}{p. 139–40 (ed. Tornberg I, p. 228–230); \Name{Ḥamzah al-Iṣfahānī}, \textit{Annales} ed. Gottwaldt,}
\DLine{AB01-PDF0641-body-L002}{}{p. 67–69; \Name{Abū ’l-fidā’}, \textit{Ta’rīkh,} Constantinopoli 1286 heg., t. I, p. 63–69; — atque tabulae}
\DLine{AB01-PDF0641-body-L003}{}{ab Augusto ad Constantinum Porphyrogenitum (m. 959 Chr.) apud \Name{Ḥamzah} p. 76-79, et \Name{Ibn}}
\DLine{AB01-PDF0641-body-L004}{}{\Name{al-Athīr} I, p. 141–48 (ed. Tornberg I, p. 230-42).}
\hypertarget{AB01-PDF0641-P01}{}
\DLine{AB01-PDF0641-body-L005}{5}{\Indent Albatenii tabulas regum (ut etiam geographicas) ab archetypo Graeco per versionem ali-}
\DLine{AB01-PDF0641-body-L006}{}{quam \textit{Syro-christianam} provenisse existimo. Syriacam originem redolent « \textit{Bukhtanaṣṣar}}
\DLine{AB01-PDF0641-body-L007}{}{« \textit{tertius,} qui est Hierosolymorum expugnator » pro \textgreek{Ναβοκολάσσαρος} (1); \textit{Belshaṣṣar} (2) pro}
\DLine{AB01-PDF0641-body-L008}{}{\textgreek{Νηρικασολάσσαρος} (3); \textit{Darius Medus} pro \textgreek{Ναβονάδιος}; \textit{Akhashwīros} pro \textgreek{Ξέρξης}; \textit{Arnabā} pro}
\DLine{AB01-PDF0641-body-L009}{}{(Ptolemaeo) \textgreek{Λάγος}; corruptela nominis \textit{Alexandri} [Severi] \textit{filii Mammaeae,} de qua p. 194}
\DLine{AB01-PDF0641-body-L010}{10}{agam; demum vox \textit{al-ḥanīf,} Iuliano adposita, qua \textgreek{Παραβάτης} « apostata » vertitur. — Ille}
\DLine{AB01-PDF0641-body-L011}{}{\textit{Syriacus} scriptor sub Leone Isaurico (717–741) floruisse videtur, quia ultimus regum est}
\DLine{AB01-PDF0641-body-L012}{}{Theodosius III, anno 717 Chr. mortuus.}
\hypertarget{AB01-PDF0641-P02}{}
\DLine{AB01-PDF0641-body-L013}{}{\Indent Non desunt parvae a textu Graeco discrepantiae. Regno prioris Artaxerxis 41 annos}
\DLine{AB01-PDF0641-body-L014}{}{\Name{Ptolemaeus} tribuit; \Name{al-Battānī} et \Name{al-Bīrūnī} 43, ita ut summa usque ad Ochum mor-}
\DLine{AB01-PDF0641-body-L015}{15}{tuum duobus augeatur annis. Regna Arogi, Darii filii Arsakh, Alexandri Macedonis, 2 + 4}
\DLine{AB01-PDF0641-body-L016}{}{+ 8 = 14 annos apud \Name{Ptolemaeum} complectuntur; 10 + 6 + 6 = 22 (seu \textarabic{ى} 10 pro \textarabic{ب} 2,}
\DLine{AB01-PDF0641-body-L017}{}{\textarabic{و} 6 pro \textarabic{د} 4, \textarabic{و} 6 pro \textarabic{ح} 8) apud \Name{al-Battānī}; 2 + 6 + 8 = 16 apud \Name{al-Bīrūnī}; qua re}
\DLine{AB01-PDF0641-body-L018}{}{regnorum summa usque ad Alexandrum mortuum 424 in \textit{Canone} \Name{Ptolemaei}, 434 apud}
\DLine{AB01-PDF0641-body-L019}{}{\Name{al-Battānī}, 428 apud \Name{al-Bīrūnī}. Miro modo p. 231 textus (t. II, p. 3 vers.) ac in Hispa-}
\DLine{AB01-PDF0641-body-L020}{20}{nica versione intervallum a Bukhtanaṣṣar ad Alexandrum mortuum \textarabic{تكد} 424 annorum statuitur,}
\DLine{AB01-PDF0641-body-L021}{}{quod in versione \textarabic{تلد} 434 ad fidem tabulae regum legendum putavi. — Error in sequentibus}
\DLine{AB01-PDF0641-body-L022}{}{evanescit, cum a Philippo (Aridaeo) nova incipiat numeratio, scilicet aera ab Alexandro mortuo}
\DLine{AB01-PDF0641-body-L023}{}{appellata (vide t. I, p. 242–43).}
\hypertarget{AB01-PDF0641-P03}{}
\DLine{AB01-PDF0641-body-L024}{}{\Indent Postea omnia bene congruunt cum \textit{Canone Ptolemaico} usque ad Antoninum mortuum;}
\DLine{AB01-PDF0641-body-L025}{25}{post quem \Name{al-Battānī} et \Name{al-Bīrūnī} Commodum cum 32 annis collocant, in unum regna}
\DLine{AB01-PDF0641-body-L026}{}{Marci Antonini (19 annorum) et Commodi (13) colligentes. Item regna Severi (18) ac Antonini}
\DLine{AB01-PDF0641-body-L027}{}{Caracallae (7) in unum conflant; sed annorum summae nulla discrepantia manifeste afficiun-}
\DLine{AB01-PDF0641-body-L028}{}{tur. Rursus a \textit{Canone} differunt in regnis Philippi Decii et Galli (6 + 1 + 3 loco 7 + 1 + 2);}
\DLine{AB01-PDF0641-body-L029}{}{summa vero non immutatur. Diocletiano 21 annos pro 20 tribuunt; parvaeque in sequentibus}
\DLine{AB01-PDF0641-body-L030}{30}{annis discrepantiae insunt, quae tamen paulatim ita vicissim compensantur ut annorum sum-}
\DLine{AB01-PDF0641-body-L031}{}{mae ad Iulianum paganum mortuum apud Graecum et Arabes aequentur. Constantino Pogo-}
\DLine{AB01-PDF0641-body-L032}{}{nato (Leontio apud \Name{al-Battānī}) 17 annos Graecus tribuit loco 16, sed Philippico 2 loco 3;}
\DLine{AB01-PDF0641-body-L033}{}{summa igitur ad Theodosium mortuum 1040 est apud omnes (4).}
\par\vspace{6pt}\hrule height.25pt\vspace{5pt}
\noindent\begin{minipage}[t]{.48\textwidth}\vspace{0pt}\fontsize{9}{11.5}\selectfont\setlength{\GutterOffset}{0pt}
\hypertarget{AB01-PDF0641-N01}{}
\DLine{AB01-PDF0641-notes_left-L001}{}{\Indent (1) \textgreek{Ναβουχοδονόσωρ} apud \Name{Georgium Syn-}}
\DLine{AB01-PDF0641-notes_left-L002}{}{\Name{cellum}, ed. G. Dindorf, I, p. 393.}
\hypertarget{AB01-PDF0641-N02}{}
\DLine{AB01-PDF0641-notes_left-L003}{}{\Indent (2) Hebraice \textit{Bēlshaṣṣar} \texthebrew{בֵּלְשַׁאצַּר} est ultimus}
\DLine{AB01-PDF0641-notes_left-L004}{}{Chaldaeorum rex (\textit{Bēl-shar-uṣur}), Nabonidi filius;}
\DLine{AB01-PDF0641-notes_left-L005}{}{\textit{Bēlṭĕshaṣṣar} \texthebrew{בֵּלְטְשַׁאצַּר} nomen Danielis in re-}
\DLine{AB01-PDF0641-notes_left-L006}{}{gia Babyloniae aula. Ambo \textgreek{Βαλτάσαρ} apud Septua-}
\DLine{AB01-PDF0641-notes_left-L007}{}{ginta, \textsyriac{ܒܠܛܫܨܪ} in Syriaca Bibliorum versione (quod}
\DLine{AB01-PDF0641-notes_left-L008}{}{\Name{Payne Smith}, \textit{Thesaurus Syriacus,} s. v., \textit{Bliṭ-}}
\DLine{AB01-PDF0641-notes_left-L009}{}{\textit{shāṣār} vel \textit{Blīṭshāṣār} legit, \Name{Bar Bahlūl} contra}
\DLine{AB01-PDF0641-notes_left-L010}{}{rectius, ut me \Name{I. Guidi} docet, Belṭshāṣār). Nomen}
\DLine{AB01-PDF0641-notes_left-L011}{}{regis Belshaṣṣar ut \Name{al-Battānī} praebent \Name{Ḥamzah}}
\DLine{AB01-PDF0641-notes_left-L012}{}{p. 94, ac \Name{aṭ-Ṭabarī} ser. I, t. I, p. 216; t. II, p. 652}
\DLine{AB01-PDF0641-notes_left-L013}{}{(ubi editor male \textarabic{بلتشصر} recepit); ex librariorum cor-}
\end{minipage}
\hfill\begin{minipage}[t]{.48\textwidth}\vspace{0pt}\fontsize{9}{11.5}\selectfont\setlength{\GutterOffset}{0pt}
\DLine{AB01-PDF0641-notes_right-L001}{}{ruptela \textarabic{بلتنصر} apud \Name{Ibn Khaldūn}, ed. Būlāq 1284}
\DLine{AB01-PDF0641-notes_right-L002}{}{heg., t. II, p. 117. Apud \Name{Ibn al-Athīr} \textarabic{بلتاصر}; \Name{Euty-}}
\DLine{AB01-PDF0641-notes_right-L003}{}{\Name{chium} (\textit{Contextio gemmarum, sive Eutychii Pa-}}
\DLine{AB01-PDF0641-notes_right-L004}{}{\textit{triarchae Alexandrini Annales ed. Edw. Pocockio,}}
\DLine{AB01-PDF0641-notes_right-L005}{}{Oxoniae 1659, t. I, p. 257) \textarabic{بلتاسر}; sed fortasse in}
\DLine{AB01-PDF0641-notes_right-L006}{}{ambobus \textarabic{ـتـ} est ex \textarabic{ـشـ} corrupta. Ad formam Syriacam,}
\DLine{AB01-PDF0641-notes_right-L007}{}{\Name{al-Bīrūnī} l. c. et \Name{Abū ’l-faraǵ} ed. Ṣālḥānī,}
\DLine{AB01-PDF0641-notes_right-L008}{}{Beyrūt 1890, p. 78, \textarabic{بلطشاصر}.}
\hypertarget{AB01-PDF0641-N03}{}
\DLine{AB01-PDF0641-notes_right-L009}{}{\Indent (3) \textgreek{Νιριγασολάσαρος} et \textgreek{Βαλτάσαρ} apud \Name{Syn-}}
\DLine{AB01-PDF0641-notes_right-L010}{}{\Name{cellum}.}
\hypertarget{AB01-PDF0641-N04}{}
\DLine{AB01-PDF0641-notes_right-L011}{}{\Indent (4) Perperam 1042 in editione Halma; error a}
\DLine{AB01-PDF0641-notes_right-L012}{}{summa Constantii incipit, ubi 994 pro 992 editor}
\DLine{AB01-PDF0641-notes_right-L013}{}{scripsit.}
\end{minipage}
\DSource{0642}{193}
\hypertarget{AB01-PDF0642-P01}{}
\DLine{AB01-PDF0642-body-L001}{}{\Indent At quamquam \textit{\textbf{numerorum}} ordo \textit{Albatenianus} cum Graeco et al-Bīrūnī congruit, in}
\DLine{AB01-PDF0642-body-L002}{}{\textit{\textbf{nominibus}} quibusdam confusio inest in codice Escurialensi, quam auferre nolui. Permutantur}
\DLine{AB01-PDF0642-body-L003}{}{enim apud \Name{al-Battānī} nomina (non numeri) Secundi Interregni et Mesesimordaci, Theodo-}
\DLine{AB01-PDF0642-body-L004}{}{sii magni et Valentis; praeterea, post duos Constantinos qui Heraclio subsequuntur, \Name{Graecus}}
\DLine{AB01-PDF0642-body-L005}{5}{et \Name{al-Bīrūnī} hoc ordine regum nomina praebent: Constantinum Pogonatum pro Leontio, Iu-}
\DLine{AB01-PDF0642-body-L006}{}{stinianum posteriorem pro Tiberio, Leontium pro Iustiniano, Tiberium pro Philippico I, Iustinianum}
\DLine{AB01-PDF0642-body-L007}{}{pro Constantino, Philippicum, Anastasium, Theodosium. Ceterum in codice Escurialensi haec}
\DLine{AB01-PDF0642-body-L008}{}{nomina adeo sunt corrupta ut in eis legendis non raro dubitatio oriatur; qua re affirmare}
\DLine{AB01-PDF0642-body-L009}{}{non ausim omnes has nominum permutationes ab al-Battānī provenire.}
\hypertarget{AB01-PDF0642-P02}{}
\DLine{AB01-PDF0642-body-L010}{10}{\Indent Restat ut de nominibus quibusdam agam.}
\hypertarget{AB01-PDF0642-P03}{}
\DLine{AB01-PDF0642-body-L011}{}{\Indent \textit{Darius Medus,} i. e. Nabonadius, in textu est Darius \textit{al-ādharī,} seu « e provincia Ādhar-}
\DLine{AB01-PDF0642-body-L012}{}{« bayǵān »; etiam in tabula geographica t. II, p. 35, nr. 69, haec regio censetur \textit{Mediae} Pto-}
\DLine{AB01-PDF0642-body-L013}{}{lemaei respondere. In archetypo Syriaco procul dubio \textit{mādāyā} \textsyriac{ܡܕܝܐ} erat, quod \Name{al-Bīrūnī},}
\DLine{AB01-PDF0642-body-L014}{}{\textit{Chron.,} p. 89 text., 101 vers. servavit: « Darius (\textarabic{دريوس}) \textit{al-Mādāy} primus ». Apud \Name{aṭ-Ṭa-}}
\DLine{AB01-PDF0642-body-L015}{15}{\Name{barī}, \textit{Annales,} ser. I, t. II, p. 649 et 652 occurrit Darius (\textarabic{داريوش}) \textit{al-Mādhawī,} quod utro-}
\DLine{AB01-PDF0642-body-L016}{}{bique dicitur nomen relativum a maiore eius quodam Mādhī vocato; apud \Name{Eutychium}}
\DLine{AB01-PDF0642-body-L017}{}{t. I, p. 257, « Darius... \textit{al-Mādī,} scilicet e progenie Mādānī » \textarabic{داريوش بن امرير المادي اعني من ولد ماداني};}
\DLine{AB01-PDF0642-body-L018}{}{apud \Name{Abū ’l-faraǵ} ed. Ṣālḥānī, p. 78 et 80, Darius (\textarabic{داريوش}) \textit{al-Mādī,} et in eiusdem Chro-}
\DLine{AB01-PDF0642-body-L019}{}{nico Syriaco: \textit{Mādāyā.} — Ex alio fonte, Persico nempe, alibi \Name{al-Bīrūnī}, \textit{Chron.,} p. 111}
\DLine{AB01-PDF0642-body-L020}{20}{text., 115 vers., \textit{Dārā al-Māhī} i. e. ex regione Māh, quae est forma Persica recentior ve-}
\DLine{AB01-PDF0642-body-L021}{}{tustissimi nominis \textit{Māda} (Media). — In textu Arabico \textarabic{داريوس}, optime cum Graeco \textgreek{Δαρεῖος} con-}
\DLine{AB01-PDF0642-body-L022}{}{venientem, recepi ut codex habet; sed vix dubito \Name{al-Battānī} in fonte suo \textarabic{داريوش} \textit{Dārīwash}}
\DLine{AB01-PDF0642-body-L023}{}{invenisse, ut fere omnes Arabum historici scriptores habent, iuxta Syriacam et Hebraicam}
\DLine{AB01-PDF0642-body-L024}{}{formam \textit{Dāryāwesh} (in inscriptionibus cuneatis Persicis \textit{Dārayawaush,} Babylonicis \textit{Dāri-}}
\DLine{AB01-PDF0642-body-L025}{25}{\textit{yāmush}).}
\hypertarget{AB01-PDF0642-P04}{}
\DLine{AB01-PDF0642-body-L026}{}{\Indent \textit{Arogus} in textu \textarabic{غيرون} et apud \Name{al-Bīrūnī}, \textit{Chron.,} p. 89 text., 101 vers. \textarabic{مرون}; quae}
\DLine{AB01-PDF0642-body-L027}{}{vix orta videntur ex \textgreek{Ἀρωγός} \textit{Canonis Ptolemaici,} vel ex \textgreek{Ἄρσης} (varr. \textgreek{Ἄρσις, Ναρσῆς, Ὀάρσης})}
\DLine{AB01-PDF0642-body-L028}{}{ut ab aliis rex ille vocatur. Erat filius Artaxersis III (Ochi); qua re ex \textgreek{Ἄρσις Ὤχου} (1) in-}
\DLine{AB01-PDF0642-body-L029}{}{terdum \textit{Arsisochus} conflatus est. Apud \Name{Eutychium} t. I, p. 266, \Name{al-Khāzinī} f. 106,v.,}
\DLine{AB01-PDF0642-body-L030}{30}{\Name{Abū ’l-faraǵ} p. 89, ipsumque \Name{al-Bīrūnī}, \textit{Chron.,} p. 111 text., 115 vers. (sed ex alio}
\DLine{AB01-PDF0642-body-L031}{}{fonte), \textit{Arsīs.}}
\hypertarget{AB01-PDF0642-P05}{}
\DLine{AB01-PDF0642-body-L032}{}{\Indent \textit{Darius} [\textit{filius}] \textit{Arsakh} (\textarabic{ارسخ}) in \textit{Canone Ptolemaico} est \textit{Darius tertius.} Plerumque ab}
\DLine{AB01-PDF0642-body-L033}{}{Arabibus, e Persico fonte, \textit{Dārā filius Dārā} vocatus. \Name{Al-Bīrūnī} tamen (\textit{Chron.} p. 89 text.,}
\DLine{AB01-PDF0642-body-L034}{}{101 vers.) \textit{Darius filius Arsīkh} (\textarabic{ارسيخ}), et in eiusdem inedito \textit{Canone Mas‘ūdico} (teste \Name{Sa-}}
\DLine{AB01-PDF0642-body-L035}{35}{\Name{chau} in notis ad vers. \textit{Chron.} p. 400) \textit{Darius filius Arsaq} \textarabic{ارسق}; \Name{Abū ’l-faraǵ} ed. Ṣāl-}
\DLine{AB01-PDF0642-body-L036}{}{ḥānī p. 91, \textit{Darius filius Arshak} \textarabic{ارشك}, ut in eiusdem \textit{Chronico Syriaco,} ed. Bruns et Kirsch,}
\DLine{AB01-PDF0642-body-L037}{}{Lipsiae 1789, p. 38 text., 36 vers. (\textsyriac{ܕܪܝܘܫ ܒܪ ܐܪܫܟ}). — False omnino \Name{Eutychius} I, 266,}
\DLine{AB01-PDF0642-body-L038}{}{\textit{Darius filius Arsīs}; et Aethiopica historiae Arabis \Name{Elmacini} (al-Makīn) redactio \textit{Dārā filius}}
\DLine{AB01-PDF0642-body-L039}{}{\textit{Arsēs} \textethiopic{ዳራ፡ ወልደ፡ አርሴስ፡} (2). — Apud \Name{al-Mas‘ūdī}, \textit{Prairies} II, p. 100, nomen occurrit}
\DLine{AB01-PDF0642-body-L040}{40}{\textarabic{دارو اليسع}, quod, non improbabili coniectura, \Name{J. Marquart} (3) putat al-Bīrūnicum \textit{Darium}}
\par\vspace{6pt}\hrule height.25pt\vspace{5pt}
\noindent\begin{minipage}[t]{.48\textwidth}\vspace{0pt}\fontsize{9}{11.5}\selectfont\setlength{\GutterOffset}{0pt}
\hypertarget{AB01-PDF0642-N01}{}
\DLine{AB01-PDF0642-notes_left-L001}{}{\Indent (1) Ut apud \Name{Eusebii} \textit{Chronicorum cano-}}
\DLine{AB01-PDF0642-notes_left-L002}{}{\textit{num libri duo edd. A. Maius et Ioh. Zohrabus,} Me-}
\DLine{AB01-PDF0642-notes_left-L003}{}{diolani 1818, p. 107. — Apud \Name{Syncellum} \textgreek{Ναρσῆς}}
\DLine{AB01-PDF0642-notes_left-L004}{}{\textgreek{Ὤχου ἀδελφός}.}
\hypertarget{AB01-PDF0642-N02}{}
\DLine{AB01-PDF0642-notes_left-L005}{}{\Indent (2) Apud L. \Name{Goldschmidt}, \textit{Die abessini-}}
\end{minipage}
\hfill\begin{minipage}[t]{.48\textwidth}\vspace{0pt}\fontsize{9}{11.5}\selectfont\setlength{\GutterOffset}{0pt}
\DLine{AB01-PDF0642-notes_right-L001}{}{\textit{schen Handschriften der Stadtbibliothek zu Frank-}}
\DLine{AB01-PDF0642-notes_right-L002}{}{\textit{furt am Main,} Berlin 1897, nr. 21, p. 75.}
\hypertarget{AB01-PDF0642-N03}{}
\DLine{AB01-PDF0642-notes_right-L003}{}{\Indent (3) \textit{Die Assyriaka des Ktesias,} in Philologus,}
\DLine{AB01-PDF0642-notes_right-L004}{}{VI Supplementband, 2. Hälfte, Göttingen 1893, p. 653,}
\DLine{AB01-PDF0642-notes_right-L005}{}{adn. 634.}
\end{minipage}
\par\vspace{5mm}\hfill{\fontsize{8}{10}\selectfont 25}\hspace{10mm}
\DSource{0643}{194}
\DLine{AB01-PDF0643-body-L001}{}{\textit{filium Arsīkh} esse, et Graeco \textgreek{Δαρεῖος Ἀρσάμου} respondere; \textgreek{Ἀρσάμης} enim (Persice \textit{Arshāma})}
\DLine{AB01-PDF0643-body-L002}{}{erat Darii III avus. Si ita res se habet, error in scriptura Syriaca explicatur, cum ex \textsyriac{ܐܪܣܡ}}
\DLine{AB01-PDF0643-body-L003}{}{\textit{Arsam} facile \textsyriac{ܐܪܣܟ} \textit{Arsak} et \textsyriac{ܐܪܣܚ} \textit{Arṡaḥ} (\textit{Arsakh}) oriantur.}
\hypertarget{AB01-PDF0643-P01}{}
\DLine{AB01-PDF0643-body-L004}{}{\Indent \textit{Philippus pater Dhū ’l-qarnayn,} loco Philippi \textit{Aridaei} qui Alexandri Magni erat fra-}
\DLine{AB01-PDF0643-body-L005}{5}{ter uterinus. Error apud Arabes haud infrequens; cfr. t. I, p. 243. — De nomine Dhū ’l-qar-}
\DLine{AB01-PDF0643-body-L006}{}{nayn t. I, p. \textsc{lxxiv}.}
\hypertarget{AB01-PDF0643-P02}{}
\DLine{AB01-PDF0643-body-L007}{}{\Indent \textit{Alexander conditor} est in \textit{Canone Ptolemaico} \textgreek{Ἀλέξανδρος ἄλλος}. Sicut in linea praece-}
\DLine{AB01-PDF0643-body-L008}{}{denti Philippum Aridaeum cum Macedone confuderat, ita nunc \Name{al-Battānī} vel eius fons}
\DLine{AB01-PDF0643-body-L009}{}{nomine Conditoris (\textit{al-bannā’}) Alexandrum alterum appellat quod, ut recte habet \Name{al-Bī-}}
\DLine{AB01-PDF0643-body-L010}{10}{\Name{rūnī} (1), ad Alexandrum Magnum spectat. Hic enim \textgreek{ὁ κτίστης} « [Alexandriae] conditor » in}
\DLine{AB01-PDF0643-body-L011}{}{\textit{Tabulis Manualibus} \Name{Ptolemaei} (ed. Halma p. 2) ac \Name{Theonis} (ed. Halma t. I, p. 27)}
\DLine{AB01-PDF0643-body-L012}{}{vocatur. — Eadem confusio apud \Name{Kūshiyār al-Ǵīlī} (\Name{Ideler}, \textit{Handb. d. mathem. u. techn.}}
\DLine{AB01-PDF0643-body-L013}{}{\textit{Chronol.,} II, 626), et \Name{Ḥāǵǵī Khalīfah} ed. Flügel III, 470, nr. 6471, qui Theonem dicit}
\DLine{AB01-PDF0643-body-L014}{}{« in opere suo astronomico, Canone appellato et ab observationibus illis deducto, aeram adhi-}
\DLine{AB01-PDF0643-body-L015}{15}{« buisse Philippi Romani, Conditoris, fratris \textgreek{τοῦ} Dhū ’l-qarnayn » (\textarabic{استعمل في زيجه المسمّى بالقانون}}
\DLine{AB01-PDF0643-body-L016}{}{\textarabic{المحصول من الرصد المذكور تاريخ فيلبس الرومي البناء اخي ذي القرنين}). Quem locum \Name{J. Lippert}, \textit{Studien}}
\DLine{AB01-PDF0643-body-L017}{}{\textit{auf dem Gebiete der griechisch-arabischen Uebersetzungslitteratur,} Braunschweig 1894,}
\DLine{AB01-PDF0643-body-L018}{}{p. 41, adn. 2, commemorans, non iniuria affirmat illud « Philippum conditorem » corruptum}
\DLine{AB01-PDF0643-body-L019}{}{esse ex aliquo Graeco textu \textgreek{Φίλιππος ὁ μετ’ Ἀλέξανδρον τὸν κτίστην} « Philippus qui [fuit] post}
\DLine{AB01-PDF0643-body-L020}{20}{Alexandrum conditorem ».}
\hypertarget{AB01-PDF0643-P03}{}
\DLine{AB01-PDF0643-body-L021}{}{\Indent \textit{Ptolemaeus Lepus.} Melius erat « Ptolemaeum \textit{Arnabā} » scribere, ut est in codice Escu-}
\DLine{AB01-PDF0643-body-L022}{}{rialensi (\textarabic{اربيا}, l. \textarabic{ارنبا}). In textu edendo t. III, p. 229, \textarabic{ارنب} \textit{arnab,} i. e. « lepus », emendandum}
\DLine{AB01-PDF0643-body-L023}{}{putavi; vix recte, nam si auctor « leporem » voluisset voci \textit{arnab} articulum \textit{al-} praeposuisset.}
\DLine{AB01-PDF0643-body-L024}{}{Similiter « Ptolemaeus filius \textit{Arnabā} » apud \Name{al-Bīrūnī}, \textit{Chron.,} p. 92 text. Eius pater enim}
\DLine{AB01-PDF0643-body-L025}{25}{Graece \textgreek{Λάγος} « lepus » vocabatur; quod nomen a Syriaco interprete optime \textit{arnĕbā} versum,}
\DLine{AB01-PDF0643-body-L026}{}{ab \Name{al-Battānī} et \Name{al-Bīrūnī} (seu potius ab utriusque fonte Arabico) pro Graeca voce ha-}
\DLine{AB01-PDF0643-body-L027}{}{bitum est. Recte alii Arabes \textit{al-arnab,} ut \Name{Ḥamzah} p. 66 text. « Ptolemaeus filius Leporis »;}
\DLine{AB01-PDF0643-body-L028}{}{\Name{Eutychius} I, 298 « Ptolemaeus qui lepus cognominatur » (\textarabic{بطليموس ويلقّب الارنب}); \Name{al-Khā-}}
\DLine{AB01-PDF0643-body-L029}{}{\Name{zinī} f. 117,r. « Ptolemaeus \textarabic{مثوس} (?) filius Lūghūs Leporis »; \Name{Abū ’l-faraǵ} p. 98 « Ptole-}
\DLine{AB01-PDF0643-body-L030}{30}{« maeus filius Lāghūs id est filius Leporis ». — Interdum \textarabic{الارنب} \textit{al-arnab} corrumpitur \textarabic{الاريب}}
\DLine{AB01-PDF0643-body-L031}{}{\textit{al-arīb} « intelligens, callidus, astutus »; vide adnot. de Goeje ad \Name{al-Mas‘ūdī}, \textit{at-Tanbīh,}}
\DLine{AB01-PDF0643-body-L032}{}{p. 114 (locus libri \textit{Mafātīḥ} ibi adductus nunc in edit. van Vloten est p. 113: \textarabic{بطليموس الاريب}}
\DLine{AB01-PDF0643-body-L033}{}{\textarabic{بن اديب}).}
\hypertarget{AB01-PDF0643-P04}{}
\DLine{AB01-PDF0643-body-L034}{}{\Indent \textit{Ptolemaeus Cleopatra.} Pro viro interdum habetur; vide quae narrat \Name{al-Bīrūnī}, \textit{Chron.,}}
\DLine{AB01-PDF0643-body-L035}{35}{p. 92 text., 103 vers.}
\hypertarget{AB01-PDF0643-P05}{}
\DLine{AB01-PDF0643-body-L036}{}{\Indent \textit{Alexander filius Mammaeae,} \textgreek{Ἀλέξανδρος Μαμαίας}, seu Alexander Severus, est in cod.}
\DLine{AB01-PDF0643-body-L037}{}{Escur. \textit{Aleksandros Brnmā} (\textarabic{برنما}), in codicibus \Name{al-Bīrūnī}, \textit{Chron.} p. 94 text., \textit{Aleksandros}}
\DLine{AB01-PDF0643-body-L038}{}{\textit{Bzymā} (\textarabic{بزيما}), quod originem Syriacam ostendit. Primus enim Arabicus huius tabulae interpres}
\DLine{AB01-PDF0643-body-L039}{}{Syriacum \textsyriac{ܒܰܪ ܡܳܡܳܐ} \textit{bar Māmā} « filium Mammaeae » pro unico vocabulo habuit, Alexandri}
\DLine{AB01-PDF0643-body-L040}{40}{cognomine; qua re \textarabic{برمما} \textit{Barmamā} scripsit, a librariis postea varie corruptum. — Forma Syriaca}
\DLine{AB01-PDF0643-body-L041}{}{genuina nominis \textgreek{Μαμαία} est \textsyriac{ܡܡܐ} \textit{Māmā,} in inscriptionibus Phoeniciis \texthebrew{ממה}, in Palmyrenis \texthebrew{ממי}}
\DLine{AB01-PDF0643-body-L042}{}{(cfr. \Name{Nöldeke}, \textit{Beiträge zur semitischen Sprachwissenschaft,} Strassburg 1904, p. 94); qua}
\par\vspace{6pt}\hrule height.25pt\vspace{5pt}
\noindent\begin{minipage}[t]{.48\textwidth}\vspace{0pt}\fontsize{9}{11.5}\selectfont\setlength{\GutterOffset}{0pt}
\hypertarget{AB01-PDF0643-N01}{}
\DLine{AB01-PDF0643-notes_left-L001}{}{\Indent (1) \textit{Chron.} 92 (omissum in vers. p. 103), et 28}
\DLine{AB01-PDF0643-notes_left-L002}{}{text. (32 vers.). — Unde \Name{al-Maqrīzī}, \textit{al-Mawā‘iẓ}}
\end{minipage}
\hfill\begin{minipage}[t]{.48\textwidth}\vspace{0pt}\fontsize{9}{11.5}\selectfont\setlength{\GutterOffset}{0pt}
\DLine{AB01-PDF0643-notes_right-L001}{}{\textit{wa ’l-i‘tibār,} Būlāq 1270 hegirae, I, 260, l. 5 a f.:}
\DLine{AB01-PDF0643-notes_right-L002}{}{\textarabic{الاسكندر البناء المقدوني} « Alexander conditor Macedon ».}
\end{minipage}
\DSource{0644}{195}
\DLine{AB01-PDF0644-body-L001}{}{re optime \Name{Abū ’l-faraǵ} ed. Ṣālḥānī p. 126, e Syriacis fontibus recepit « \textit{Aleksandrūs,} cuius}
\DLine{AB01-PDF0644-body-L002}{}{« matris nomen erat \textit{Māmā} ». — Traditio autem Arabica quae ex \textit{Graeco} manat, genitivum}
\DLine{AB01-PDF0644-body-L003}{}{Graecum \textgreek{Μαμαίας} servat, quem pro Alexandri cognomine habuit; ita \Name{al-Mas‘ūdī}, \textit{Prai-}}
\DLine{AB01-PDF0644-body-L004}{}{\textit{ries} II, 306, \textit{al-Iskandar Māmayās;} \Name{al-Mas‘ūdī}, \textit{at-Tanbīh} p. 133, « \textit{al-Ikṣandros}}
\DLine{AB01-PDF0644-body-L005}{5}{(\textarabic{الاكصندرس}) et cognominatur \textit{Māmayās} », sicut etiam \Name{Ibn al-Athīr}, ed. Kahirensis, t. I,}
\DLine{AB01-PDF0644-body-L006}{}{p. 143 (perperam 145 impressum) in altero regum catalogo habet.}
\hypertarget{AB01-PDF0644-P01}{}
\DLine{AB01-PDF0644-body-L007}{}{\Indent \textit{Decius.} In \textit{Canone} \textgreek{Δέκιος}, et apud \textit{historicos} Arabum scriptores \textit{Dāqiyūs, Daqyūs, Dī-}}
\DLine{AB01-PDF0644-body-L008}{}{\textit{qiyūs.} Sed \Name{al-Battānī} \textit{Dāqiyānūs,} ut semper hic rex vocatur in \textit{Arabico-Muslimicis} re-}
\DLine{AB01-PDF0644-body-L009}{}{dactionibus fabulosae narrationis de « gente cavernae » (\textit{ahl al-kahf}) i. e. de septem Dormien-}
\DLine{AB01-PDF0644-body-L010}{10}{tibus Ephesi; \textit{Deqyānūs} apud \Name{al-Ḥasan al-Qummī an-Naysābūrī}, \textit{Tafsīr} (in marg.}
\DLine{AB01-PDF0644-body-L011}{}{\textit{tafsīr} aṭ-Ṭabarī) t. XV, p. 114; \Name{ath-Tha‘labī}, \textit{‘Arā’is al-maǵālis,} Kahirae 1297 heg.,}
\DLine{AB01-PDF0644-body-L012}{}{p. 369 sqq.; \Name{az-Zamakhsharī}, \textit{al-Kashshāf,} Kahirae 1307 heg., t. I, p. 567; \Name{al-Bay-}}
\DLine{AB01-PDF0644-body-L013}{}{\Name{ḍāwī} ed. Fleischer, t. I, p. 559 seq.; \Name{ad-Damīrī}, \textit{Ḥayāt al-ḥayawān,} Kahirae 1311 heg.,}
\DLine{AB01-PDF0644-body-L014}{}{t. II, p. 240–45 (s. v. \textarabic{كلب}); [Pseudo-] \Name{Ibn Iyās}, \textit{Badā’i‘ az-zuhūr,} Kahirae 1310 heg.,}
\DLine{AB01-PDF0644-body-L015}{15}{p. 182–85, etc. — \Name{Aṭ-Ṭabarī} in \textit{Annalibus,} ser. I, t. II, p. 778 (de septem Dormientibus}
\DLine{AB01-PDF0644-body-L016}{}{agens) ac magno in Qoranum \textit{Commentario} t. XV, p. 123–26, 128, 134–38, \textit{Daqīnūs} \textarabic{دقينوس}.}
\DLine{AB01-PDF0644-body-L017}{}{Nomen exstat in vico \textit{Shahr-i-Daqyānūs} « Deciani urbs » quo sunt ruinae urbis Ǵiruft (vide}
\DLine{AB01-PDF0644-body-L018}{}{tabulas geographicas nr. 255). — Miro modo \textit{Daqyūs} a \textit{Daqyānūs} distinguit lexicon \textit{Qāmūs}}
\DLine{AB01-PDF0644-body-L019}{}{inscriptum eiusque Commentarius \textit{Tāǵ al-‘arūs,} ed. altera, t. IV, p. 152: \textarabic{(دقيوس بالفتح) اسم}}
\DLine{AB01-PDF0644-body-L020}{20}{\textarabic{(ملك) اعجمي (اتخذ مسجدًا على اصحاب الكهف) زاد الصاغاني (ودقيانوس) اسم (ملك هربوا منه) وقصتهم مذكورة}. — In scri-}
\DLine{AB01-PDF0644-body-L021}{}{ptis Syriacis \textsyriac{ܕܩܝܘܣ} vel \textsyriac{ܕܩܝܣ} \textit{Deqyūs;} unde \textit{Daqyūs} in historia \textit{Arabico-Christiana} septem}
\DLine{AB01-PDF0644-body-L022}{}{Dormientium apud \Name{I. Guidi}, \textit{Testi orientali inediti sopra i Sette Dormienti di Efeso,} in}
\DLine{AB01-PDF0644-body-L023}{}{Memorie dell’Accad. dei Lincei, Classe di scienze morali, ser. III, t. XII, 1884, p. 391, adn. 1;}
\DLine{AB01-PDF0644-body-L024}{}{tamen \textit{Deqyānūs} in alia narratione \textit{Arabico-Christiana} apud \Name{Guidi} p. 301. In \textit{Synaxario}}
\DLine{AB01-PDF0644-body-L025}{25}{Arabico-Coptico Iacobitarum \textit{Dākiyūs} (\Name{Guidi} p. 426), ex fonte Graeco, cum quo Aethiopicum}
\DLine{AB01-PDF0644-body-L026}{}{\textethiopic{ዳክዮስ፡} \textit{Dākyōs} congruit, semper in libris sub Zar’a Yā‘qōb (1434–1468 Chr.) conscriptis}
\DLine{AB01-PDF0644-body-L027}{}{occurrens (apud \Name{Guidi} p. 406–19 et 429), vel cum \textethiopic{ዳኬዎስ፡} \textit{Dākēwōs} Synaxarii Aethiopici}
\DLine{AB01-PDF0644-body-L028}{}{(apud \Name{Guidi} p. 427). — Formam \textit{Dāqiyānūs} vel \textit{Deqyānūs} vel \textit{Daqīnūs} puto e scriptura}
\DLine{AB01-PDF0644-body-L029}{}{Syriaca ortam, ob \textsyriac{ܕܩܝܘܣ} male \textsyriac{ܕܩܝܢܘܣ} lectum (1).}
\hypertarget{AB01-PDF0644-P02}{}
\DLine{AB01-PDF0644-body-L030}{30}{\Indent \textit{Iulianus} « \textit{paganus} » (p. 2 et 3). Arabice \textit{al-ḥanīf,} eodem sensu quo vox Syriaca \textit{ḥanpā;}}
\DLine{AB01-PDF0644-body-L031}{}{in Graeco \textgreek{παραβάτης} « apostata ». De illa verbi \textit{ḥanīf} significatione vide Glossarium De Goeje}
\DLine{AB01-PDF0644-body-L032}{}{ad \Name{al-Mas‘ūdī}, \textit{at-Tanbīh,} p. \textsc{xvii–xviii}.}
\hypertarget{AB01-PDF0644-P03}{}
\DLine{AB01-PDF0644-body-L033}{}{\Indent \textit{Heraclius} (\textit{Hiraqlis} et \textit{Hiraql}) « \textit{sapiens} » (p. 2 et 3). In codice \textarabic{صاحب العرب} \textit{ṣāḥib al-}}
\DLine{AB01-PDF0644-body-L034}{}{\textit{-‘arab} « dominus arabum », quod t. III, p. 230 ac 231 servavi, licet nullo modo explicare}
\DLine{AB01-PDF0644-body-L035}{35}{possem; item versio Hispanica paulum infra adducenda: « Hercules \textit{el de los Alaraues} ». Le-}
\DLine{AB01-PDF0644-body-L036}{}{gendum duco \textarabic{صاحب العرف} \textit{ṣāḥib al-‘irf} « dominus cognitionis, cognitione praeditus »; littera \textarabic{ف}}
\DLine{AB01-PDF0644-body-L037}{}{in scriptura maghrebina punctum inferne habet, qua re facillime cum \textarabic{ب} confunditur. Coniec-}
\DLine{AB01-PDF0644-body-L038}{}{tura mea innititur \Name{al-Bīrūnī}, \textit{Chron.,} p. 96 text., 105 vers., ubi \textit{Hiraqlis al-ḥakīm} « He-}
\DLine{AB01-PDF0644-body-L039}{}{« raclius sapiens ». Cur ita vocatus non bene intelligo; in Graeco \textgreek{Ἡράκλειος} tantum. Fortasse}
\DLine{AB01-PDF0644-body-L040}{40}{quod Heraclius imperator de rebus astronomicis scripsit; fragmentum eius chronologicum me-}
\DLine{AB01-PDF0644-body-L041}{}{moravi t. I, p. 244–45. Anno 618 Chr. luculentam in Tabulas Manuales Ptolemaei expositio-}
\par\vspace{6pt}\hrule height.25pt\vspace{5pt}
\noindent\begin{minipage}[t]{.48\textwidth}\vspace{0pt}\fontsize{9}{11.5}\selectfont\setlength{\GutterOffset}{0pt}
\hypertarget{AB01-PDF0644-N01}{}
\DLine{AB01-PDF0644-notes_left-L001}{}{\Indent (1) Cfr. M. J. \Name{de Goeje} sententiam: « Daaren-}
\DLine{AB01-PDF0644-notes_left-L002}{}{« tegen hebben de Arabische, naar mijne meening}
\DLine{AB01-PDF0644-notes_left-L003}{}{« uit oude Syrische bronnen afkomstige, berichten}
\DLine{AB01-PDF0644-notes_left-L004}{}{« Dakianus » (\textit{De legende der Zevenslapers van Efe-}}
\end{minipage}
\hfill\begin{minipage}[t]{.48\textwidth}\vspace{0pt}\fontsize{9}{11.5}\selectfont\setlength{\GutterOffset}{0pt}
\DLine{AB01-PDF0644-notes_right-L001}{}{\textit{ze,} in Verslagen en Mededeelingen der kon. Akad.}
\DLine{AB01-PDF0644-notes_right-L002}{}{van Wetenschappen, Afdeeling Letterkunde, 4\textsuperscript{e} reeks,}
\DLine{AB01-PDF0644-notes_right-L003}{}{deel III, Amsterdam 1900, p. 12).}
\end{minipage}
\DSource{0645}{196}
\DLine{AB01-PDF0645-body-L001}{}{nem composuit, qua supputationes Solis, Lunae, quinque planetarum docuit, sicut apparet ex}
\DLine{AB01-PDF0645-body-L002}{}{\Name{Iac. Morellii} \textit{Bibliotheca manuscripta Graeca et Latina,} Bassani 1802, p. 205–06 et 209;}
\DLine{AB01-PDF0645-body-L003}{}{cfr. etiam \Name{Stephani} philosophi (priore VIII Chr. saeculi medietate) fragmentum de arte ma-}
\DLine{AB01-PDF0645-body-L004}{}{thematica editum in \textit{Catalogo codicum astrologorum Graecorum,} fasc. II, Brussellis 1900,}
\DLine{AB01-PDF0645-body-L005}{5}{p. 182: \textgreek{Ὅ τε γὰρ Πτολομαῖος} [sic] \textgreek{τοῖς ἀπὸ τοῦ Ναβουχοδόνοσορ ἔτεσιν ἐχρήσατο καὶ μησὶν Αἰγυ}-}
\DLine{AB01-PDF0645-body-L006}{}{\textgreek{πτιακοῖς, ὁ δέ γε Θέων καὶ Ἡράκλειος καὶ ὁ Ἀμμώνιος τοῖς τοῦ Φιλίππου καὶ μησὶν Αἰγυπτιακοῖς}.}
\DLine{AB01-PDF0645-body-L007}{}{— Aḥmad ibn ‘Abd al-ḥalīm \Name{Ibn Taymiyyah} in libro \textit{takhǵīl ahl al-inǵīl} « vituperationis}
\DLine{AB01-PDF0645-body-L008}{}{« asseclarum Evangelii » (1), Heraclium imperatorem affirmat insignem astrologum fuisse, ac}
\DLine{AB01-PDF0645-body-L009}{}{e stellarum motibus adventum « Regis circumcisionis », nempe Muḥammad, cognovisse; sed,}
\DLine{AB01-PDF0645-body-L010}{10}{iuxta vetustiores Arabicas traditiones (2), eum Heraclius in somnio viderat. Narrant alii He-}
\DLine{AB01-PDF0645-body-L011}{}{raclium, epistula Prophetae accepta, ad novam religionem se populumque suum convertere}
\DLine{AB01-PDF0645-body-L012}{}{primo voluisse, sed principum populique minas irritum tale consilium fecisse (3). — Philoso-}
\DLine{AB01-PDF0645-body-L013}{}{phorum concilium apud Heraclium (\textit{Hiraql}) regem, eorumque de musica sententias, ex aucto-}
\DLine{AB01-PDF0645-body-L014}{}{ritate \textarabic{اموس} (Ammonii?) commemorat \Name{al-Kindī} (IX saec. Chr.) in inedito de musica libello,}
\DLine{AB01-PDF0645-body-L015}{15}{tract. II, cap. IV, apud \Name{Steinschneider} ad \Name{Baldi} p. 437 (extr. 12, cum additamentis). —}
\DLine{AB01-PDF0645-body-L016}{}{Demum non impossibilis censenda est aliqua cum mythico \textit{Hiraql} (Heraclio aut Hercule) con-}
\DLine{AB01-PDF0645-body-L017}{}{fusio, quem \Name{Ibn al-Qifṭī} (in adnot. Flügel ad \textit{Fihrist,} t. II, p. 125) dicit « sapientem (\textit{ḥa-}}
\DLine{AB01-PDF0645-body-L018}{}{« \textit{kīm}) Babylonium, unum e septem sapientibus » qui olim septem templis Babylonis praefuis-}
\DLine{AB01-PDF0645-body-L019}{}{sent. Huic Hiraql tribuit \Name{an-Nadīm}, \textit{al-Fihrist,} p. 270 et 285, librum \textit{ad-dawāyir wa}}
\DLine{AB01-PDF0645-body-L020}{20}{\textit{’d-dawālīb} « de rotis [omnimodis] et rotis aquariis » (4); eidem aut alii Hiraql magnum de}
\DLine{AB01-PDF0645-body-L021}{}{\textit{alchymia} librum ipse \Name{an-Nadīm}, p. 353. l. 24, atque p. 354, l. 27, et \Name{Ḥāǵǵī Khalīfah}}
\DLine{AB01-PDF0645-body-L022}{}{ed. Flügel, t. V, p. 171 (cfr. etiam \Name{Steinschneider}, \textit{Die arabischen Uebersetzungen aus}}
\DLine{AB01-PDF0645-body-L023}{}{\textit{dem Griechischen,} in Zeitschr. d. deutsch. morgenl. Gesellsch., t. L, 1896, p. 362, ubi alchy-}
\DLine{AB01-PDF0645-body-L024}{}{miae scriptor dicitur, quo fundamento nescio, ipse Heraclius imperator esse). — In Bibliotheca}
\DLine{AB01-PDF0645-body-L025}{25}{Nationali Parisiensi codex Arabicus XIII saeculi Chr. asservatur, quem Catalogus impressus}
\DLine{AB01-PDF0645-body-L026}{}{(nr. 178, 3\textsuperscript{o}) ita describit: « Comment Sibylle (\textarabic{سبلا}) fille de Héraclius (\textarabic{هرقل}), chef des philosophes,}
\DLine{AB01-PDF0645-body-L027}{}{« expliqua le songe que cent hommes eurent, la même nuit, dans la ville de Rome, sous le}
\DLine{AB01-PDF0645-body-L028}{}{« règne d’Auguste César ». — Mythicus Hercules est cuius imago in codice quodam Graeco}
\DLine{AB01-PDF0645-body-L029}{}{Bononiensi occurrit, subscripta \textgreek{Ἀστρονόμος ὁ Ἱρακλῇς} (\textit{Catal. codd. astrologorum Graecorum,}}
\DLine{AB01-PDF0645-body-L030}{30}{fasc. IV, 1903, p. 45); cfr. \Name{Bouché-Leclercq}, \textit{L’astrologie grecque,} p. 577: « C’est bien}
\par\vspace{6pt}\hrule height.25pt\vspace{5pt}
\noindent\begin{minipage}[t]{.48\textwidth}\vspace{0pt}\fontsize{9}{11.5}\selectfont\setlength{\GutterOffset}{0pt}
\hypertarget{AB01-PDF0645-N01}{}
\DLine{AB01-PDF0645-notes_left-L001}{}{\Indent (1) Apud L. \Name{Marraccium}, \textit{Alcorani textus}}
\DLine{AB01-PDF0645-notes_left-L002}{}{\textit{universus,} Patavii 1698, t. I (Prodromus ad refuta-}
\DLine{AB01-PDF0645-notes_left-L003}{}{tionem Alcorani), pars I, p. 44. — Heraclium tan-}
\DLine{AB01-PDF0645-notes_left-L004}{}{quam verum Prophetam Mahometum coluisse affir-}
\DLine{AB01-PDF0645-notes_left-L005}{}{mat idem \Name{Ibn Taymiyyah}, \textit{ar-risālah al-qu-}}
\DLine{AB01-PDF0645-notes_left-L006}{}{\textit{bruṣiyyah,} Kahirae 1319 heg., p. 15. — Obiit auctor}
\DLine{AB01-PDF0645-notes_left-L007}{}{primo diluculo (\textit{saḥar}) diei Lunae 20 dhū ’l-qa‘dah}
\DLine{AB01-PDF0645-notes_left-L008}{}{728, i. e. ultima nocte inter diem Dominicam 25 et}
\DLine{AB01-PDF0645-notes_left-L009}{}{diem Lunae 26 Septembris 1328 (« Dienstag 22. Dū}
\DLine{AB01-PDF0645-notes_left-L010}{}{« ’lqa‘da = 29 Sept. » apud \Name{Brockelmann} II,}
\DLine{AB01-PDF0645-notes_left-L011}{}{p. 101 falsum, ac e \Name{Schreiner}, ZDMG, LII, 1898,}
\DLine{AB01-PDF0645-notes_left-L012}{}{p. 562 repetitum; falsum quoque « primo diluculo}
\DLine{AB01-PDF0645-notes_left-L013}{}{« diei Lunae 10 dhū ’l-qa‘dah » apud \Name{Ibn al-Ālūsī},}
\DLine{AB01-PDF0645-notes_left-L014}{}{\textit{Ǵalā’ al-‘aynayn,} Būlāq 1298 heg., p. 9, unde flu-}
\DLine{AB01-PDF0645-notes_left-L015}{}{xit in biographiam antepositam libro \Name{Ibn Taymiy-}}
\DLine{AB01-PDF0645-notes_left-L016}{}{\Name{yah}, \textit{al-wāsiṭah bayn al-khalq wa ’l-ḥaqq,} Ka-}
\DLine{AB01-PDF0645-notes_left-L017}{}{hirae 1318 heg., p. 8).}
\end{minipage}
\hfill\begin{minipage}[t]{.48\textwidth}\vspace{0pt}\fontsize{9}{11.5}\selectfont\setlength{\GutterOffset}{0pt}
\hypertarget{AB01-PDF0645-N02}{}
\DLine{AB01-PDF0645-notes_right-L001}{}{\Indent (2) \textit{Kitāb al-Aghānī} t. VI, p. 94; \Name{aṭ-Ṭabarī},}
\DLine{AB01-PDF0645-notes_right-L002}{}{\textit{Annales,} ser. I, t. III, p. 1561 sq.}
\hypertarget{AB01-PDF0645-N03}{}
\DLine{AB01-PDF0645-notes_right-L003}{}{\Indent (3) \textit{Aghānī} VI, p. 94–96 (cfr. XV, p. 58, l. 14-15);}
\DLine{AB01-PDF0645-notes_right-L004}{}{\Name{aṭ-Ṭabarī}, l. c., p. 1562-68; \Name{Ibn Sa‘d} apud}
\DLine{AB01-PDF0645-notes_right-L005}{}{\Name{Wellhausen}, \textit{Skizzen,} t. IV (Berlin 1899) p. 2}
\DLine{AB01-PDF0645-notes_right-L006}{}{text., 98 vers.; \Name{al-Wāqidī}, vers. Wellhausen,}
\DLine{AB01-PDF0645-notes_right-L007}{}{p. 329; \Name{Ibn al-Athīr}, Kahirae 1301-02, t. II,}
\DLine{AB01-PDF0645-notes_right-L008}{}{p. 101–02 (ed. Tornberg II, p. 162-63); \Name{al-Bukhārī},}
\DLine{AB01-PDF0645-notes_right-L009}{}{Būlāq 1289 heg., t. I, p. 3-5, II, p. 130-31 (ǵihād),}
\DLine{AB01-PDF0645-notes_right-L010}{}{III, p. 93-94 (tafsīr, ad Qor. III, 57).}
\hypertarget{AB01-PDF0645-N04}{}
\DLine{AB01-PDF0645-notes_right-L011}{}{\Indent (4) \textit{Dawāyir} verti « rotas omnimodas »; cfr.}
\DLine{AB01-PDF0645-notes_right-L012}{}{\textit{Le livre des appareils pneumatiques et des machi-}}
\DLine{AB01-PDF0645-notes_right-L013}{}{\textit{nes hydrauliques de} \Name{Philon de Byzance} \textit{édité}}
\DLine{AB01-PDF0645-notes_right-L014}{}{\textit{d’après les versions arabes et traduit par Carra}}
\DLine{AB01-PDF0645-notes_right-L015}{}{\textit{de Vaux,} in Notices et extraits des mss. de la}
\DLine{AB01-PDF0645-notes_right-L016}{}{Bibliothèque Nationale, t. XXXVIII, Paris 1902, p. 29,}
\DLine{AB01-PDF0645-notes_right-L017}{}{30, 231.}
\end{minipage}
\DSource{0646}{197}
\DLine{AB01-PDF0646-body-L001}{}{« l’Hercules de Tyr que Nonnus appelle \textgreek{ἀστροχίτων} (XL, 408, 579), \textgreek{πρόμος ἄστρων} (ibid. 567).}
\DLine{AB01-PDF0646-body-L002}{}{« Cet Hercule fait un cours d’histoire à Bacchus et lui fait cadeau de sa robe constellée, la}
\DLine{AB01-PDF0646-body-L003}{}{« vêtement légendaire des astrologues ».}
\hypertarget{AB01-PDF0646-P01}{}
\DLine{AB01-PDF0646-body-L004}{}{\Indent Conspectum intervallorum iuxta \textit{Hispanicam} versionem (f. 29,r.) nunc praebebo: « Lo}
\DLine{AB01-PDF0646-body-L005}{5}{« que a entre Nabuchodonosor el primero fata la muerte de Alexandro el primero .CCCC. et}
\DLine{AB01-PDF0646-body-L006}{}{« .XXIIII. (1) annos egipcianos. {\CapFont ⸿} Et depues [sic] regno Phelippo padre de Alexandre el}
\DLine{AB01-PDF0646-body-L007}{}{« Macedonio. et ouo entrel. et entrel començamiento del regnado de Augusto el primero ce-}
\DLine{AB01-PDF0646-body-L008}{}{« sar .CC. et .LXIIII. (2) annos. {\CapFont ⸿} Et desde que regno Cesar Augusto el romano. fata el co-}
\DLine{AB01-PDF0646-body-L009}{}{« mençamiento del regno de Dioclecianos el xristiano .CCC. et .XIII. annos. {\CapFont ⸿} Et desde que}
\DLine{AB01-PDF0646-body-L010}{10}{« regno Dioclecianos fata el començamiento del regno de Bolianoz (3) .LXXVII. annos. {\CapFont ⸿} Et}
\DLine{AB01-PDF0646-body-L011}{}{« despues mataron (4) a Bolianoz. et tornosse el regno a los xristianos assi como fue. (5) Et}
\DLine{AB01-PDF0646-body-L012}{}{« fue del començamiento del regno de Dioclecianos. fata el comienço que regno Hercules el}
\DLine{AB01-PDF0646-body-L013}{}{« de los Alaraues .CCC. et .XXVI. (6) annos. {\CapFont ⸿} Et desi regnaron los Alaraues del comença-}
\DLine{AB01-PDF0646-body-L014}{}{« miento de la alhigera. fata que regno Mauya que fue de fiios de Vmaya. depues [sic] de}
\DLine{AB01-PDF0646-body-L015}{15}{« la muerte de Ahly aben Abitali (7) .XXXIX. annos. et finco el regnado en los fijos de}
\DLine{AB01-PDF0646-body-L016}{}{« Vmaya. fata que se torno a los fijos de Alahbez a complimiento de .C. et .XXVII. annos}
\DLine{AB01-PDF0646-body-L017}{}{« de la alhigera ».}
\par\vspace{3mm}
\hypertarget{AB01-PDF0646-H18}{}
\DLine{AB01-PDF0646-body-L018}{}{\CenterLine{\textit{Ad pag. 4-6.} — \textbf{Tabula khalifarum.}}}
\par\vspace{2mm}
\hypertarget{AB01-PDF0646-P02}{}
\DLine{AB01-PDF0646-body-L019}{}{\Indent Eadem (8), usque ad al-Muṭī‘ mortuum (336 heg.) porrecta apud \Name{al-Mas‘ūdī}, \textit{Prai-}}
\DLine{AB01-PDF0646-body-L020}{20}{\textit{ries,} t. IX, p. 38–49, qui se \Name{al-Battānī} praesertim secutum esse fatetur; eadem etiam, usque}
\DLine{AB01-PDF0646-body-L021}{}{ad al-Muktafī mortuum (334 heg.) perducta, at summis annorum omissis et non sine erro-}
\DLine{AB01-PDF0646-body-L022}{}{ribus, apud \Name{Ḥamzah} p. 156–58 text. Iuvat nonnulla ex \Name{al-Mas‘ūdī} l. c., p. 38–39 referre:}
\DLine{AB01-PDF0646-body-L023}{}{« In anteriore huius libri parte iam peculiare caput constitueramus pro regum chronologia ....}
\DLine{AB01-PDF0646-body-L024}{}{« usque ad hoc tempus, deducta e supputatione atque ex iis quae narrantur in libris biogra-}
\DLine{AB01-PDF0646-body-L025}{25}{« phiarum, historiarum ac eorum qui historiae khalifarum et regum vacaverunt. Sed illo loco}
\DLine{AB01-PDF0646-body-L026}{}{« non exposueramus quae in astronomicis operibus (\textit{fī kutub az-zīǵāt}) narrant astronomi}
\DLine{AB01-PDF0646-body-L027}{}{« (\textit{aṣḥāb an-nuǵūm}) sua temporum ratione innitentes. Hoc igitur capite omnia quae illi in}
\DLine{AB01-PDF0646-body-L028}{}{« astronomicis operibus statuerint ab hegira ad nostrum tempus memorabimus, ut libri utilitas}
\DLine{AB01-PDF0646-body-L029}{}{« augeatur et cognitio discrepantiae inter historicos ac caelestium rerum scriptores deprehen-}
\DLine{AB01-PDF0646-body-L030}{30}{« datur ». Postea, p. 48–49: « Inter hanc et historicorum biographorumque chronologiam di-}
\DLine{AB01-PDF0646-body-L031}{}{« screpantia (\textit{tafāwut}) est ob incrementa (\textit{ziyādāt}) mensium et dierum. In chronologia ab he-}
\DLine{AB01-PDF0646-body-L032}{}{« gira ad hoc tempus statuenda usi sumus iis quae in operibus astronomicis invenimus; nam}
\DLine{AB01-PDF0646-body-L033}{}{« qui huic [rerum caelestium] arti vacant tempora illa accurate servant (\textit{yurā‘ūn}) eorumque}
\DLine{AB01-PDF0646-body-L034}{}{« cognitionem absolutissime definitam (\textit{‘alà ’t-taḥdīd}) adipiscuntur. Quae retulimus, ex opere}
\DLine{AB01-PDF0646-body-L035}{35}{« astronomico Abū ‘Abd Allāh Muḥammad ibn Ǵābir \Name{al-Battānī} (9) aliisque astronomicis}
\DLine{AB01-PDF0646-body-L036}{}{« operibus sumpsimus ». — De tabulis chronologicis vide etiam \Name{al-Mas‘ūdī}, \textit{at-Tanbīh,} p. 329.}
\par\vspace{6pt}\hrule height.25pt\vspace{5pt}
\noindent\begin{minipage}[t]{.48\textwidth}\vspace{0pt}\fontsize{9}{11.5}\selectfont\setlength{\GutterOffset}{0pt}
\hypertarget{AB01-PDF0646-N01}{}
\DLine{AB01-PDF0646-notes_left-L001}{}{\Indent (1) Legit \textarabic{تكد} ut in cod. Escurialensi; vide su-}
\DLine{AB01-PDF0646-notes_left-L002}{}{pra p. 192.}
\hypertarget{AB01-PDF0646-N02}{}
\DLine{AB01-PDF0646-notes_left-L003}{}{\Indent (2) Perperam legit \textarabic{رصد} pro \textarabic{رضد} 294.}
\hypertarget{AB01-PDF0646-N03}{}
\DLine{AB01-PDF0646-notes_left-L004}{}{\Indent (3) Legit \textarabic{بليانس}, ut cod. Escur. hac pagina ha-}
\DLine{AB01-PDF0646-notes_left-L005}{}{bet, pro \textarabic{يليانس} (cod. Esc. f. 155,r. \textarabic{يوليانس}); omisit \textit{al-}}
\DLine{AB01-PDF0646-notes_left-L006}{}{\textit{-ḥanīf} « paganum ».}
\hypertarget{AB01-PDF0646-N04}{}
\DLine{AB01-PDF0646-notes_left-L007}{}{\Indent (4) Legit \textarabic{قتل} « interfectus est » pro \textarabic{ملك} « re-}
\DLine{AB01-PDF0646-notes_left-L008}{}{gnavit ».}
\end{minipage}
\hfill\begin{minipage}[t]{.48\textwidth}\vspace{0pt}\fontsize{9}{11.5}\selectfont\setlength{\GutterOffset}{0pt}
\hypertarget{AB01-PDF0646-N05}{}
\DLine{AB01-PDF0646-notes_right-L001}{}{\Indent (5) Textum corruptum manifeste interpretari}
\DLine{AB01-PDF0646-notes_right-L002}{}{nescivit; postea duorum annorum intervallum omisit.}
\hypertarget{AB01-PDF0646-N06}{}
\DLine{AB01-PDF0646-notes_right-L003}{}{\Indent (6) Legit \textarabic{سكو} ut cod. Escur.}
\hypertarget{AB01-PDF0646-N07}{}
\DLine{AB01-PDF0646-notes_right-L004}{}{\Indent (7) I. e. ‘Alī ibn Abī Ṭālib; deest incisum in}
\DLine{AB01-PDF0646-notes_right-L005}{}{cod. Escur.}
\hypertarget{AB01-PDF0646-N08}{}
\DLine{AB01-PDF0646-notes_right-L006}{}{\Indent (8) Textus Arabicus in t. III, p. 231-34. In ver-}
\DLine{AB01-PDF0646-notes_right-L007}{}{sione Hispanica omissae.}
\hypertarget{AB01-PDF0646-N09}{}
\DLine{AB01-PDF0646-notes_right-L008}{}{\Indent (9) In edit. \textarabic{البناني}, qua re in versione Gallica}
\DLine{AB01-PDF0646-notes_right-L009}{}{\textit{Bennani.}}
\end{minipage}
\DSource{0647}{198}
\hypertarget{AB01-PDF0647-P01}{}
\DLine{AB01-PDF0647-body-L001}{}{\Indent Anni et menses huius tabulae lunares sunt, iuxta Muslimorum calendarium. De hegira}
\DLine{AB01-PDF0647-body-L002}{}{die Iovis constituta in adnot. ad tabulam p. 7 dicam.}
\hypertarget{AB01-PDF0647-P02}{}
\DLine{AB01-PDF0647-body-L003}{}{\Indent « \textit{Khalifarum iustorum} » nomen in inscriptionibus servavi, sicut in codice inveneram.}
\DLine{AB01-PDF0647-body-L004}{}{Sed plerumque nonnisi quatuor priores khalifae ita vocantur; nec impossibile est nomen, ab}
\DLine{AB01-PDF0647-body-L005}{5}{\Name{al-Battānī} quatuor illis dumtaxat praepositum, amanuensium errore etiam reliquis omnibus}
\DLine{AB01-PDF0647-body-L006}{}{tributum esse.}
\hypertarget{AB01-PDF0647-P03}{}
\DLine{AB01-PDF0647-body-L007}{}{\Indent \textit{Banū Hāshim,} in initio p. 5 memorati, tribus vel gens sunt a qua khalifae ‘Abbāsidae}
\DLine{AB01-PDF0647-body-L008}{}{originem ducebant.}
\hypertarget{AB01-PDF0647-P04}{}
\DLine{AB01-PDF0647-body-L009}{}{\Indent Discrepantiae ab \Name{al-Mas‘ūdī} minimi sunt momenti, ut ex earum conspectu patebit:}
\hypertarget{AB01-PDF0647-body-T01}{}\RawBlock{\hspace*{14mm}\begin{tabular}{@{}l@{\hspace{6mm}}r@{\hspace{2.2mm}}r@{\hspace{2.2mm}}r@{\hspace{3mm}}|@{\hspace{3mm}}r@{\hspace{2.2mm}}r@{\hspace{2.2mm}}r@{\hspace{3mm}}||@{\hspace{3mm}}r@{\hspace{2.2mm}}r@{\hspace{2.2mm}}r@{\hspace{3mm}}|@{\hspace{3mm}}r@{\hspace{2.2mm}}r@{\hspace{2.2mm}}r@{}}
& \multicolumn{6}{c||}{\Name{al-Battānī} p. 4} & \multicolumn{6}{c}{\Name{al-Mas‘ūdī} IX, 42}\\[3pt]
‘Omar ibn ‘Abd al-Malik & 2\sa & 5\sm & \textit{13}\sd & 100\sa & 6\sm & \textit{26}\sd & 2\sa & 5\sm & \textit{15}\sd & 100\sa & 6\sm & \textit{28}\sd\\
Yazīd ibn ‘Abd al-Malik & 4 & 0 & 1 & 104 & 6 & \textit{27} & 4 & 0 & 1 & 104 & 6 & \textit{29}\\
Hishām ibn ‘Abd al-Malik & 19 & 8 & \textit{9} & 124 & 3 & 6 & 19 & 8 & \textit{7} & 124 & 3 & 6\\[8pt]
& \multicolumn{6}{c||}{\Name{al-Battānī} p. 5} & \multicolumn{6}{c}{\Name{al-Mas‘ūdī} IX, 46}\\[3pt]
al-Mutawakkil & 14\sa & 9\sm & \textit{9}\sd & 246\sa & 9\sm & \textit{3}\sd & 14\sa & 9\sm & \textit{7}\sd & 246\sa & 9\sm & \textit{1}\sd\\
al-Muntaṣir & 0 & 6 & 0 & 247 & 3 & \textit{3} & 0 & 6 & 0 & 247 & 3 & \textit{1}\\
al-Musta‘īn & 2 & 9 & 3 & 250 & 0 & \textit{6} & 2 & 9 & 3 & 250 & 0 & \textit{4}\\
usque ad al-Mu‘tazz & 0 & 0 & \textit{8} & 250 & 0 & 14 & 0 & 0 & \textit{10} & 250 & 0 & 14\\
\end{tabular}}
\hypertarget{AB01-PDF0647-P05}{}
\DLine{AB01-PDF0647-body-L010}{10}{\Indent Omnes qui post \textit{al-Muktafī} memorantur (Italicis seu inclinatis litteris expressi), quoniam}
\DLine{AB01-PDF0647-body-L011}{}{seriore aetate quam \Name{Albatenii} fuerunt, manifeste ab aliquo lectore additi sunt, qui in an-}
\DLine{AB01-PDF0647-body-L012}{}{norum \textit{al-Muqtadir} summa perperam 18 dies loco 11 posuit et reliquas summas igitur etiam}
\DLine{AB01-PDF0647-body-L013}{}{vitiavit. Nec parvae sunt in his interpolatis discrepantiae ab \Name{al-Mas‘ūdī}.}
\hypertarget{AB01-PDF0647-P06}{}
\DLine{AB01-PDF0647-body-L014}{}{\Indent Sed fortasse ultimus khalifa ab \Name{al-Battānī} memoratus erat \textit{al-Mu‘taḍid;} numeri enim}
\DLine{AB01-PDF0647-body-L015}{15}{sequentis al-Muktafī cum \Name{al-Mas‘ūdī} IX, 47 non congruunt:}
\par\vspace{1mm}
\hypertarget{AB01-PDF0647-H16}{}
\DLine{AB01-PDF0647-body-L016}{}{\CenterLine{\Name{al-Battānī} 6\sa\ 6\sm\ \textit{4}\sd\ \ |\ \ 294\sa\ \textit{9}\sm\ \textit{26}\sd\ \ ||\ \ \Name{al-Mas‘ūdī} 6\sa\ 6\sm\ \textit{20}\sd\ \ |\ \ 294\sa\ \textit{10}\sm\ \textit{12}\sd}}
\par\vspace{1mm}
\hypertarget{AB01-PDF0647-P07}{}
\DLine{AB01-PDF0647-body-L017}{}{\Indent Historici scriptores, sicut \Name{aṭ-Ṭabarī}, ser. III, t. V, p. 2280, \Name{al-Mas‘ūdī}, \textit{at-Tanbīh,}}
\DLine{AB01-PDF0647-body-L018}{}{p. 371 (ubi \textarabic{وستّة اشهر} inserendum), aliique, khalifatum al-Muktafī 6 annorum, 6 mensium, 19}
\DLine{AB01-PDF0647-body-L019}{}{dierum faciunt.}
\hypertarget{AB01-PDF0647-P08}{}
\DLine{AB01-PDF0647-body-L020}{20}{\Indent Rogo demum ut p. 5, ult. l., pro codicis \textarabic{كح} 28 restituantur \textarabic{كج} 23 anni.}
\par\vspace{4mm}
\hypertarget{AB01-PDF0647-H21}{}
\DLine{AB01-PDF0647-body-L021}{}{\CenterLine{\textit{Ad pag. 7.} — \textbf{Signa annorum Arabicorum.}}}
\par\vspace{3mm}
\hypertarget{AB01-PDF0647-P09}{}
\DLine{AB01-PDF0647-body-L022}{}{\Indent Tabulam auctor memorat t. I, p. 71 (cap. XXXII), explicat p. 145 (adp. A).}
\hypertarget{AB01-PDF0647-P10}{}
\DLine{AB01-PDF0647-body-L023}{}{\Indent Duae signorum rationes occurrunt; altera, quam I vocabo, in sinistra paginae parte, al-}
\DLine{AB01-PDF0647-body-L024}{}{tera (II) in dextera. Mihi II, in Hispanica versione (f. 29,r.) omissa, \textit{spuria} videtur; inutilis}
\DLine{AB01-PDF0647-body-L025}{25}{prorsus cum iam I praebeatur, et altera minus commoda, quia ad signum anni dati reperien-}
\DLine{AB01-PDF0647-body-L026}{}{dum necesse est signo illius anni in tabula scripto signum mensis al-muḥarram semper ad-}
\DLine{AB01-PDF0647-body-L027}{}{dere. Signum primi anni hegirae erit nempe 7 + 5 = 5, signum secundi 4 + 5 = 2, etc.}
\DLine{AB01-PDF0647-body-L028}{}{Iuxta rationem I signum anni iam in tabula est, nec additio requiritur nisi ad invenienda signa}
\DLine{AB01-PDF0647-body-L029}{}{mensium qui primo (al-muḥarram) subsequuntur. \Name{Al-Battānī}, l. c., rationem I exponit; II ne}
\DLine{AB01-PDF0647-body-L030}{30}{memorat quidem (1).}
\par\vspace{6pt}\hrule height.25pt\vspace{5pt}{\fontsize{9}{11.5}\selectfont
\hypertarget{AB01-PDF0647-N01}{}
\DLine{AB01-PDF0647-notes-L001}{}{\Indent (1) Tertia ratio ex opere \Name{Maslamah al-Maǵrīṭī}, in tabulis spuriis f. 240,v. et 241,r.}
}
\DSource{0648}{199}
\hypertarget{AB01-PDF0648-P01}{}
\DLine{AB01-PDF0648-body-L001}{}{\Indent Signa annorum collectorum iustis esse unitate minora \Name{al-Battānī} monet; cur ita fa-}
\DLine{AB01-PDF0648-body-L002}{}{ctum exposui t. I, p. 323.}
\hypertarget{AB01-PDF0648-P02}{}
\DLine{AB01-PDF0648-body-L003}{}{\Indent « Anni collecti » (\textit{sinūn maǵmū‘ah}) significant, apud astronomos Arabicos, summam an-}
\DLine{AB01-PDF0648-body-L004}{}{norum ab epocha aerae, non annorum spatium a quocumque temporis momento numeratum.}
\DLine{AB01-PDF0648-body-L005}{5}{« Anni singuli » vel, ut Arabice vocantur, « expansi » (\textit{sinūn mabsūṭah}), sunt anni singilla-}
\DLine{AB01-PDF0648-body-L006}{}{tim sumpti, a quocumque tempore numerati. Nomina annorum \textit{collectorum} et \textit{expansorum}}
\DLine{AB01-PDF0648-body-L007}{}{e scriptis Arabicis in \textit{Tabulas Toletanas, Alphonsinas} aliasque multas aevi medii et recen-}
\DLine{AB01-PDF0648-body-L008}{}{tiores irrepserunt.}
\hypertarget{AB01-PDF0648-P03}{}
\DLine{AB01-PDF0648-body-L009}{}{\Indent Apud Arabes, ut apud Latinos veteres et astronomos omnes, hebdomada a die dominica}
\DLine{AB01-PDF0648-body-L010}{10}{(feria prima) orditur; sabbatum (feria septima) hebdomadam concludit.}
\hypertarget{AB01-PDF0648-P04}{}
\DLine{AB01-PDF0648-body-L011}{}{\Indent \Name{Al-Battānī} epocham hegirae ponit \textit{die Iovis} (feria quinta) 15 Iulii 622 Chr. Ex astro-}
\DLine{AB01-PDF0648-body-L012}{}{nomis Arabicis (1) quos inspexi, soli \Name{Ḥabash} et \Name{al-Khāzinī} (f. 22,v. etc.) usum civilem}
\DLine{AB01-PDF0648-body-L013}{}{Muslimicum sequuntur epochae hegirae die Veneris (feria sexta) 16 Iulii constituendae. Iuvat}
\DLine{AB01-PDF0648-body-L014}{}{quae \Name{Ḥabash} f. 2,v. narrat suis ipsis verbis referre: « Cum sacrarum traditionum narrato-}
\DLine{AB01-PDF0648-body-L015}{15}{« res (2) et viros computandi peritissimos invenissem novilunium (\textarabic{مستهَلّ}) al-muḥarram anni quo}
\DLine{AB01-PDF0648-body-L016}{}{« Arabicus (3) Propheta Muḥammad (propitius est salutemque aeternam praestat ei Deus et}
\DLine{AB01-PDF0648-body-L017}{}{« purae genti eius domus) feria quinta ponere, investigavi verumne hoc esset (\textarabic{يصحّ}) iuxta com-}
\DLine{AB01-PDF0648-body-L018}{}{« putationem. Coniunctionem luminarium ad novilunium al-muḥarram anni primi hegirae re-}
\DLine{AB01-PDF0648-body-L019}{}{« peri fuisse Mekkah in urbe post mediam noctem cuius mane (\textarabic{صبيحة}) feria quarta [i. e. dies}
\DLine{AB01-PDF0648-body-L020}{20}{« Mercurii] fuit, ac supputando inveni apparitionem Lunae novae in Mekkah fuisse nocte diei}
\DLine{AB01-PDF0648-body-L021}{}{« Veneris novilunium al-muḥarram anni hegirae. Igitur nihil invenimus quod feriam quintam}
\DLine{AB01-PDF0648-body-L022}{}{« [diem Iovis] indicaret in coniunctione aut in apparitione; unde nobis patuit gentem illam}
\DLine{AB01-PDF0648-body-L023}{}{« al-muḥarram a nocte diei Iovis iuxta usum communem dumtaxat computasse, toto calculo}
\DLine{AB01-PDF0648-body-L024}{}{« suo mensem dhū ’l-ḥiǵǵah elapsum [statuendo] 29 dierum fuisse (4). Igitur ad apparitionem}
\DLine{AB01-PDF0648-body-L025}{25}{« Lunae novae reversi sumus, et initium al-muḥarram anni primi hegirae constituimus nocte}
\DLine{AB01-PDF0648-body-L026}{}{« diei Veneris in Mekkah, [cuius] mane dies XXVIII mensis afrawardīn-māh anni 44 ab}
\DLine{AB01-PDF0648-body-L027}{}{« Abarwīz (5) rege, et 9 annis ac 338 diebus ante Yazdaǵird regem. Item nocte cuius mane}
\par\vspace{6pt}\hrule height.25pt\vspace{5pt}
\noindent\begin{minipage}[t]{.48\textwidth}\vspace{0pt}\fontsize{9}{11.5}\selectfont\setlength{\GutterOffset}{0pt}
\hypertarget{AB01-PDF0648-N01}{}
\DLine{AB01-PDF0648-notes_left-L001}{}{\Indent (1) Sicut astronomi, ita astrologi diem Iovis}
\DLine{AB01-PDF0648-notes_left-L002}{}{15 Iulii 622 pro epocha aerae habere solent; e qui-}
\DLine{AB01-PDF0648-notes_left-L003}{}{bus \Name{al-Kindī} apud O. \Name{Loth}, \textit{al-Kindî als As-}}
\DLine{AB01-PDF0648-notes_left-L004}{}{\textit{trolog,} in Morgenländische Forschungen, Festschrift}
\DLine{AB01-PDF0648-notes_left-L005}{}{an H. L. Fleischer, Leipzig 1875, p. 294. — Cum 3624}
\DLine{AB01-PDF0648-notes_left-L006}{}{dierum faciant intervallum inter aeras hegirae et}
\DLine{AB01-PDF0648-notes_left-L007}{}{Iezdegerdis, patet in eadem sententia versari \Name{Al-}}
\DLine{AB01-PDF0648-notes_left-L008}{}{\Name{bumasar} [Abū Ma‘shar] \textit{de magnis coniunctio-}}
\DLine{AB01-PDF0648-notes_left-L009}{}{\textit{nibus,} Venetiis 1515, tract. IV in fine, ac tract. VIII,}
\DLine{AB01-PDF0648-notes_left-L010}{}{diff. 2; et \Name{Alchabitium} [al-Qabīṣī], diff. IV,}
\DLine{AB01-PDF0648-notes_left-L011}{}{cap. VI, fol. dd 3,r. edit. 1485, vel f. 17,r. edit. 1521}
\DLine{AB01-PDF0648-notes_left-L012}{}{(titulos ambarum dedi t. I, p. 309, adn. 1).}
\hypertarget{AB01-PDF0648-N02}{}
\DLine{AB01-PDF0648-notes_left-L013}{}{\Indent (2) \textit{Ruwāt al-ḥadīth,} i. e. qui dicta et facta}
\DLine{AB01-PDF0648-notes_left-L014}{}{Prophetae eiusque sociorum diligentissime collecta}
\DLine{AB01-PDF0648-notes_left-L015}{}{tradebant.}
\hypertarget{AB01-PDF0648-N03}{}
\DLine{AB01-PDF0648-notes_left-L016}{}{\Indent (3) Ex ipso \textit{Qorano} apparet Qoranum esse « \textit{li-}}
\DLine{AB01-PDF0648-notes_left-L017}{}{\textit{brum Arabicum} », vel « \textit{claro Arabico sermone}}
\DLine{AB01-PDF0648-notes_left-L018}{}{\textit{scriptum} », omnesque Prophetas qui ante Muḥam-}
\DLine{AB01-PDF0648-notes_left-L019}{}{mad fuerint ad alias quam Arabicam nationes mis-}
\DLine{AB01-PDF0648-notes_left-L020}{}{sos. Unde nomen « \textit{Prophetae Arabici} » quo saepe}
\end{minipage}
\hfill\begin{minipage}[t]{.48\textwidth}\vspace{0pt}\fontsize{9}{11.5}\selectfont\setlength{\GutterOffset}{0pt}
\DLine{AB01-PDF0648-notes_right-L001}{}{Muḥammad vocatur, praesertim in initio ac fine car-}
\DLine{AB01-PDF0648-notes_right-L002}{}{minum interiectorum popularibus iisque longissimis}
\DLine{AB01-PDF0648-notes_right-L003}{}{narrationibus \textit{Sīrat Banī Hilāl} inscriptis ac res ge-}
\DLine{AB01-PDF0648-notes_right-L004}{}{stas gentis Banī Hilāl concelebrantibus (in editt.}
\DLine{AB01-PDF0648-notes_right-L005}{}{Berytensibus, a Christianis factis, illi primi ultimi-}
\DLine{AB01-PDF0648-notes_right-L006}{}{que versus omittuntur).}
\hypertarget{AB01-PDF0648-N04}{}
\DLine{AB01-PDF0648-notes_right-L007}{}{\Indent (4) Triginta dierum fuerat iuxta peculiarem}
\DLine{AB01-PDF0648-notes_right-L008}{}{\Name{Ḥabash} intercalandi rationem, qua ultimus annus}
\DLine{AB01-PDF0648-notes_right-L009}{}{cycli tricennalis pro bisextili habetur; cfr. t. I, p. 323.}
\DLine{AB01-PDF0648-notes_right-L010}{}{Ab hegira novus ciclus orditur.}
\hypertarget{AB01-PDF0648-N05}{}
\DLine{AB01-PDF0648-notes_right-L011}{}{\Indent (5) Chosroes II, cuius regnum a die 27 Iunii}
\DLine{AB01-PDF0648-notes_right-L012}{}{590 Chr. supputatum esse certissime constat; vide}
\DLine{AB01-PDF0648-notes_right-L013}{}{\Name{Tabari}’s \textit{Geschichte der Perser und Araber zur}}
\DLine{AB01-PDF0648-notes_right-L014}{}{\textit{Zeit der Sasaniden übers. von Th. Nöldeke,} Lei-}
\DLine{AB01-PDF0648-notes_right-L015}{}{den 1879, p. 430-31, \Name{Th. Nöldeke}, \textit{Zur persi-}}
\DLine{AB01-PDF0648-notes_right-L016}{}{\textit{schen Chronologie,} in ZDMG., t. L, 1896, p. 141, et}
\DLine{AB01-PDF0648-notes_right-L017}{}{etiam tabulas omnes Sāsānidarum regum apud \Name{al-}}
\DLine{AB01-PDF0648-notes_right-L018}{}{\Name{-Bīrūnī}, \textit{Chron.,} p. 121-31 text., 123-28 vers. Ergo}
\DLine{AB01-PDF0648-notes_right-L019}{}{16 Iulii 622 fuit XXVIII dies mensis afrawardīn anni 32}
\DLine{AB01-PDF0648-notes_right-L020}{}{(non 44 ut in codice Ḥabash) a Chosroe II rege.}
\end{minipage}
\DSource{0649}{200}
\DLine{AB01-PDF0649-body-L001}{}{« fuit XVI dies tammūz [Iulii] anni 933 Alexandri; et a primo die tishrīn prioris (1) ad XV}
\DLine{AB01-PDF0649-body-L002}{}{« tammūz, toto XV computato, 288 dies fuerunt. Item mane XXIII diei mensis kiyahk}
\DLine{AB01-PDF0649-body-L003}{}{« anni 1389 a Bukhtanaṣṣar; et a primo die tūt ad XXII kiyahk, die XXII computato,}
\DLine{AB01-PDF0649-body-L004}{}{« 112 dies fuerunt, qui, additis quinque epagomenis (\textit{al-ayyām al-lawāḥiq}), 117 fiunt. Hos}
\DLine{AB01-PDF0649-body-L005}{5}{« quinque propterea adiecimus quia a [mense Persico] daymāh (2) numeramus, [ex quo] cum}
\DLine{AB01-PDF0649-body-L006}{}{« 12 elapsi sunt annus Coptorum perfectus est ».}
\hypertarget{AB01-PDF0649-P01}{}
\DLine{AB01-PDF0649-body-L007}{}{\Indent Iuxta civilem Muslimorum usum dies ab occasu Solis computatur, ita ut nox ad sequens}
\DLine{AB01-PDF0649-body-L008}{}{mane pertineat. Astronomi diem a meridie supputant (t. I, p. 221); ergo nox diei Veneris iuxta}
\DLine{AB01-PDF0649-body-L009}{}{usum civilem fit apud eos nox diei Iovis. Diei discrepantia in epocha hegirae ita explicatur.}
\DLine{AB01-PDF0649-body-L010}{10}{\Name{Ḥabash} autem aliis calculis fretus est.}
\hypertarget{AB01-PDF0649-P02}{}
\DLine{AB01-PDF0649-body-L011}{}{\Indent Usum mensium lunarium Mekkanis Propheta praecepisse videtur anno 9 hegirae, quo}
\DLine{AB01-PDF0649-body-L012}{}{\textit{Qorani} IX, 36–37 versiculos probabiliter edidit (3); praecepta anno sequenti in ultima sua}
\DLine{AB01-PDF0649-body-L013}{}{contione iteravit (4). Quo calendario \Name{Mekkani} antea uterentur non liquet; doctissimorum}
\DLine{AB01-PDF0649-body-L014}{}{virorum coniecturas vide apud A. \Name{Caussin de Perceval}, \textit{Mémoire sur le calendrier}}
\DLine{AB01-PDF0649-body-L015}{15}{\textit{arabe avant l’Islamisme} (Journal Asiatique, ser. IV, t. 1, 1843, p. 342–79); \Name{Mahmoud}}
\DLine{AB01-PDF0649-body-L016}{}{\Name{Effendi}, \textit{Mémoire sur le calendrier arabe avant l’Islamisme et sur la naissance et l’âge}}
\DLine{AB01-PDF0649-body-L017}{}{\textit{du prophète Mohammad} (Journ. Asiat., V, 11, 1858, p. 109–92); A. \Name{Sprenger}, \textit{Ueber den}}
\DLine{AB01-PDF0649-body-L018}{}{\textit{Kalender der Araber vor Muḥammad} (Zeitschr. d. deutsch. morgenländ. Gesellsch., t. XIII,}
\DLine{AB01-PDF0649-body-L019}{}{1859, p. 134-75); J. \Name{Wellhausen}, \textit{Reste arabischen Heidentums,} ed. altera, Berlin 1897,}
\DLine{AB01-PDF0649-body-L020}{20}{p. 94–101; H. \Name{Winckler}, \textit{Altorientalische Forschungen,} ser. II, t. II (1900), p. 324–50,}
\DLine{AB01-PDF0649-body-L021}{}{374–81. Suam ipsam sententiam postea \Name{Winckler} (5) paulo mutavit, affirmans Prophetae}
\DLine{AB01-PDF0649-body-L022}{}{tempore Mekkanos calendarium \textit{solare} habuisse, a Christianis Arabiae australis, Abrahah duce,}
\DLine{AB01-PDF0649-body-L023}{}{post 525 Chr. allatum; Prophetam autem, post Mekkam anno 8 heg. captam, vetustissimum}
\DLine{AB01-PDF0649-body-L024}{}{calendarium Mekkanorum \textit{lunare} restituisse quod ad cultum Dei lunaris Hubal olim perti-}
\DLine{AB01-PDF0649-body-L025}{25}{nuisset.}
\hypertarget{AB01-PDF0649-P03}{}
\DLine{AB01-PDF0649-body-L026}{}{\Indent Nonnulli, quos \Name{aṭ-Ṭabarī} (\textit{Annales,} ser. I, t. III, p. 1250–56) recenset, epocham aerae}
\DLine{AB01-PDF0649-body-L027}{}{ab ipso Muḥammad institutam dicunt; plurimi, rectius quidem, a khalifa ‘Omar anno 16 (vel}
\DLine{AB01-PDF0649-body-L028}{}{17 vel 18). Epocha est I dies al-muḥarram anni quo Propheta, ab urbe Mekkah secedens,}
\DLine{AB01-PDF0649-body-L029}{}{al-Madīnah pervenit; quod, iuxta Arabicos et nostros astronomos, die Lunae 8 rabī‘ prioris}
\DLine{AB01-PDF0649-body-L030}{30}{iuxta usum civilem (20 Sept. 622) fuit vel die Dominica 7 eiusdem (19 Sept.) iuxta astronomicam}
\DLine{AB01-PDF0649-body-L031}{}{temporum rationem. Vide \Name{Mahmoud Effendi}, l. c., p. 121; \Name{Loth}, \textit{al-Kindî als Astrolog,}}
\DLine{AB01-PDF0649-body-L032}{}{p. 293–94; \Name{al-Bīrūnī}, \textit{Chron.,} p. 330 text., 327–28 vers. Errant Arabes illi qui diem 12}
\DLine{AB01-PDF0649-body-L033}{}{rabī‘ I deligunt.}
\hypertarget{AB01-PDF0649-P04}{}
\DLine{AB01-PDF0649-body-L034}{}{\Indent De codicis erroribus in annis bisextilibus designandis dixi t. I, p. 323, II; omnia ad \Name{Al-}}
\DLine{AB01-PDF0649-body-L035}{35}{\Name{batenii} mentem restitui.}
\par\vspace{6pt}\hrule height.25pt\vspace{5pt}
\noindent\begin{minipage}[t]{.48\textwidth}\vspace{0pt}\fontsize{9}{11.5}\selectfont\setlength{\GutterOffset}{0pt}
\hypertarget{AB01-PDF0649-N01}{}
\DLine{AB01-PDF0649-notes_left-L001}{}{\Indent (1) Seu Octobris; iuxta communem usum (ab}
\DLine{AB01-PDF0649-notes_left-L002}{}{\textit{Albateniano} diversum) annos Seleucidarum a I Octo-}
\DLine{AB01-PDF0649-notes_left-L003}{}{bris computat.}
\hypertarget{AB01-PDF0649-N02}{}
\DLine{AB01-PDF0649-notes_left-L004}{}{\Indent (2) Aera Nabonassaris die I daymāh incepit.}
\hypertarget{AB01-PDF0649-N03}{}
\DLine{AB01-PDF0649-notes_left-L005}{}{\Indent (3) Iuxta \Name{Nöldeke}, \textit{Geschichte des Qorâns,}}
\DLine{AB01-PDF0649-notes_left-L006}{}{Göttingen 1860, p. 166.}
\hypertarget{AB01-PDF0649-N04}{}
\DLine{AB01-PDF0649-notes_left-L007}{}{\Indent (4) \Name{Ibn Hishām}, ed. Wüstenfeld, p. 968–69;}
\DLine{AB01-PDF0649-notes_left-L008}{}{\Name{aṭ-Ṭabarī}, \textit{Annales,} ser. I, t. IV, p. 1753-54; om-}
\DLine{AB01-PDF0649-notes_left-L009}{}{nesque sacrarum traditionum scriptores. De tota ea}
\DLine{AB01-PDF0649-notes_left-L010}{}{contione, varie tradita partimque fortasse ficta, cfr.}
\DLine{AB01-PDF0649-notes_left-L011}{}{\Name{Snouck Hurgronje}, \textit{Het mekkaansche feest,}}
\end{minipage}
\hfill\begin{minipage}[t]{.48\textwidth}\vspace{0pt}\fontsize{9}{11.5}\selectfont\setlength{\GutterOffset}{0pt}
\DLine{AB01-PDF0649-notes_right-L001}{}{Leiden 1880, p. 43–45. — Nullam certam rationem}
\DLine{AB01-PDF0649-notes_right-L002}{}{in chronologia rerum a Propheta gestarum habuisse}
\DLine{AB01-PDF0649-notes_right-L003}{}{vetustissimos historiae scriptores iure monet \Name{Ed.}}
\DLine{AB01-PDF0649-notes_right-L004}{}{\Name{Sachau}, p. \textsc{xxxv} praefationis suae in I partem}
\DLine{AB01-PDF0649-notes_right-L005}{}{tomi III operis \Name{Ibn Saad}, \textit{Biographien Muham-}}
\DLine{AB01-PDF0649-notes_right-L006}{}{\textit{meds, seiner Gefährten und der späteren Träger}}
\DLine{AB01-PDF0649-notes_right-L007}{}{\textit{des Islams hrsg. von E. Sachau,} Leiden 1904 sqq.}
\hypertarget{AB01-PDF0649-N05}{}
\DLine{AB01-PDF0649-notes_right-L008}{}{\Indent (5) \textit{Arabisch-Semitisch-Orientalisch, kultur-}}
\DLine{AB01-PDF0649-notes_right-L009}{}{\textit{geschichtlich-mythologische Untersuchung,} Berlin}
\DLine{AB01-PDF0649-notes_right-L010}{}{1901-02, p. 81–90 (= Mitteilungen der vorderasiati-}
\DLine{AB01-PDF0649-notes_right-L011}{}{schen Gesellschaft, VI Jahrg., 1901, fasc. 4–5).}
\end{minipage}
\DSource{0650}{201}
\par\vspace{8mm}
\hypertarget{AB01-PDF0650-H01}{}
\DLine{AB01-PDF0650-body-L001}{}{\CenterLine{\textit{Ad pag. 8.} — \textbf{Signa mensium annorumque Romanorum.}}}
\par\vspace{4mm}
\hypertarget{AB01-PDF0650-P01}{}
\DLine{AB01-PDF0650-body-L002}{}{\Indent Huius tabulae, t. I, p. 71 (cap. XXXII) memoratae, auctor usum explanavit in adp. B}
\DLine{AB01-PDF0650-body-L003}{}{(t. I, p. 146). Cur signa unitate sint deminuta et cur anni per 28 dividendi exposui t. I, p. 323.}
\hypertarget{AB01-PDF0650-P02}{}
\DLine{AB01-PDF0650-body-L004}{}{\Indent In tabulis similibus apud \Name{al-Bīrūnī}, \textit{Chron.,} p. 195 text., 175 vers., et \Name{Oloug Beg},}
\DLine{AB01-PDF0650-body-L005}{5}{p. 18 text., 17 vers., nulla deminutione signa afficiuntur, et aerae initium prima die Octobris,}
\DLine{AB01-PDF0650-body-L006}{}{ut communis est mos, non Septembris constituitur.}
\hypertarget{AB01-PDF0650-P03}{}
\DLine{AB01-PDF0650-body-L007}{}{\Indent Columna bisextilium praebet sexagesimas partes diei. Cum annus Romanus 365 $\qfrac{1}{4}$ die-}
\DLine{AB01-PDF0650-body-L008}{}{rum seu 365 $\qfrac{15}{60}$ sit, sexagesimae sibi ipsis adiectae diem integrum intercalandum quarto}
\DLine{AB01-PDF0650-body-L009}{}{quoque anno conficiunt.}
\hypertarget{AB01-PDF0650-P04}{}
\DLine{AB01-PDF0650-body-L010}{10}{\Indent Nomina Latina mensium puto ab \textit{interpolatore Arabico-occidentali} addita. Nunquam ab}
\DLine{AB01-PDF0650-body-L011}{}{\Name{al-Battānī} in capitibus de chronologia memorantur; differunt etiam a nominibus apud scri-}
\DLine{AB01-PDF0650-body-L012}{}{ptores \textit{Arabico-orientales} (1) occurrentibus, scilicet: \textit{yānuwāriyūs, fībruwāriyūs, marṭi-}}
\DLine{AB01-PDF0650-body-L013}{}{\textit{yūs, abrīliyūs, māyūs, yūniyūs, yūliyūs, aghūsṭūs, sābṭānbriyūs, oqṭūmbriyūs, nuwām-}}
\DLine{AB01-PDF0650-body-L014}{}{\textit{briyūs, dīqembriyūs,} quae e Graecis proveniunt: \textgreek{Ἰανουάριος, Φεβρουάριος, Μάρτιος, Ἀπρίλλιος},}
\DLine{AB01-PDF0650-body-L015}{15}{\textgreek{Μάϊος, Ἰούνιος, Ἰούλιος, Αὔγουστος, Σεπτέμβριος, Ὀκτώβριος, Νοέμβριος, Δεκέμβριος}. Apud \textit{Arabes}}
\DLine{AB01-PDF0650-body-L016}{}{\textit{orientales} nomina directe e \textit{Latino} aut linguis \textit{Neolatinis} sumpta nonnisi XII saeculo Christi}
\DLine{AB01-PDF0650-body-L017}{}{innotuere, per Cruciatorum bella et per foedera ab Italicis praesertim rebuspublicis cum Ae-}
\DLine{AB01-PDF0650-body-L018}{}{gypti dynastibus inita (2); unde factum est ut hoc nostro tempore menses calendarii Iuliani}
\DLine{AB01-PDF0650-body-L019}{}{ab Arabibus \textit{Aegypti} ita designentur: \textit{yanāyir, fabrāyir, mārs, abrīl, māyū, yūniyah, yū-}}
\DLine{AB01-PDF0650-body-L020}{20}{\textit{liyah, oghusṭus, sebtembar, oktūbar, nōfembar, dīsembar.} — Miro modo \textit{mārt, māyus} et}
\DLine{AB01-PDF0650-body-L021}{}{\textit{oghostūs} (\textarabic{اغستوس}) soli recepti sunt in annum Iulianum quo \textit{Turcae,} post 1675 Chr., in rebus}
\DLine{AB01-PDF0650-body-L022}{}{ad aerarium spectantibus utuntur, dum reliqui menses Syriacis nominibus vocantur.}
\hypertarget{AB01-PDF0650-P05}{}
\DLine{AB01-PDF0650-body-L023}{}{\Indent Nomina huius tabulae spuria, etiam in tabula f. 241,v. (ex opere Arabico-Hispani \Name{Ma-}}
\DLine{AB01-PDF0650-body-L024}{}{\Name{slamah al-Maǵrīṭī} sumpta) occurrentia, haec sunt: \textit{yannayr, fabrayr, mārs, ibrīl, mā-}}
\DLine{AB01-PDF0650-body-L025}{25}{\textit{yuh, yūniyuh, yulyuh} (\textarabic{يليه}), \textit{aghusht, sutambar, aktūbar, nuwanbar, duǵanbar.} Eadem}
\DLine{AB01-PDF0650-body-L026}{}{(paulum mutatis \textit{marsuh, ibrīr, yūliyuh}) appellantur « menses \textit{occidentalium} » (\textarabic{المغاربة}) ab}
\DLine{AB01-PDF0650-body-L027}{}{\Name{al-Bīrūnī}, \textit{Chron.,} p. 50 et 71 text., 59 et 83 vers.; et in astrolabio occurrunt anno 420 heg.}
\DLine{AB01-PDF0650-body-L028}{}{(inc. 19 Ianuar. 1029) Toleti confecto, paulum mutatis \textit{yūliyuh, shutanbar, dukanbar} (3).}
\DLine{AB01-PDF0650-body-L029}{}{— Sub finem VI saec. heg. (XII Chr.), \Name{Ibn Ǵubayr}, scriptor Arabico-Hispanus, iter suum}
\DLine{AB01-PDF0650-body-L030}{30}{per Siciliam describens memorat \textit{yannayr} (apud \Name{Amari}, \textit{Bibl. Arabo-sicula,} text. p. 96).}
\par\vspace{6pt}\hrule height.25pt\vspace{5pt}
\noindent\begin{minipage}[t]{.48\textwidth}\vspace{0pt}\fontsize{9}{11.5}\selectfont\setlength{\GutterOffset}{0pt}
\hypertarget{AB01-PDF0650-N01}{}
\DLine{AB01-PDF0650-notes_left-L001}{}{\Indent (1) Iique paucissimi sunt: \Name{Ḥabash} f. 4,v., qui}
\DLine{AB01-PDF0650-notes_left-L002}{}{ea appellat nomina « mensium Graecorum (\textit{al-yū-}}
\DLine{AB01-PDF0650-notes_left-L003}{}{\textit{nāniyyīn}) »; \Name{al-Bīrūnī}, \textit{Chron.,} p. 50 text. (59}
\DLine{AB01-PDF0650-notes_left-L004}{}{vers.) et 71 text. (83 vers.), « menses Graecorum »}
\DLine{AB01-PDF0650-notes_left-L005}{}{aut « Romanorum »; Chems eddin Abou Abdallah}
\DLine{AB01-PDF0650-notes_left-L006}{}{Mohammed \Name{ed-Dimichqui}, \textit{Cosmographie, texte}}
\DLine{AB01-PDF0650-notes_left-L007}{}{\textit{arabe publié d’après l’édition commencée par C.}}
\DLine{AB01-PDF0650-notes_left-L008}{}{\textit{M. Fraehn par M. F. Mehren,} St. Pétersbourg 1866,}
\DLine{AB01-PDF0650-notes_left-L009}{}{cap. IX, sect. 8, « menses Graecorum et Romano-}
\DLine{AB01-PDF0650-notes_left-L010}{}{rum ». — \textarabic{ذقمبرس} in subscriptione codicis Arabico-}
\DLine{AB01-PDF0650-notes_left-L011}{}{-Christiani, in Syria anno 1188 Alexandri, 877 Chr.,}
\end{minipage}
\hfill\begin{minipage}[t]{.48\textwidth}\vspace{0pt}\fontsize{9}{11.5}\selectfont\setlength{\GutterOffset}{0pt}
\DLine{AB01-PDF0650-notes_right-L001}{}{exarati (vide periodicum \textit{al-Mashriq,} t. VI, 1903,}
\DLine{AB01-PDF0650-notes_right-L002}{}{pag. 1011).}
\hypertarget{AB01-PDF0650-N02}{}
\DLine{AB01-PDF0650-notes_right-L003}{}{\Indent (2) Ut \textit{māy, yūniyuh, aghusht,} in foederibus}
\DLine{AB01-PDF0650-notes_right-L004}{}{XII ac XIII saec. inter Genuenses et Aegyptum, apud}
\DLine{AB01-PDF0650-notes_right-L005}{}{M. \Name{Amari}, \textit{Nuovi ricordi arabici su la storia di}}
\DLine{AB01-PDF0650-notes_right-L006}{}{\textit{Genova,} Genova 1873, p. 5, 10, 16 (= Atti della So-}
\DLine{AB01-PDF0650-notes_right-L007}{}{cietà Ligure di storia patria, t. V).}
\hypertarget{AB01-PDF0650-N03}{}
\DLine{AB01-PDF0650-notes_right-L008}{}{\Indent (3) F. \Name{Woepcke}, \textit{Ueber ein in der k. Bi-}}
\DLine{AB01-PDF0650-notes_right-L009}{}{\textit{bliothek zu Berlin befindliches arabisches Astro-}}
\DLine{AB01-PDF0650-notes_right-L010}{}{\textit{labium,} in Abhandl. der mathem. Klasse der k. Ak.}
\DLine{AB01-PDF0650-notes_right-L011}{}{der Wissensch. zu Berlin, Jahrg. 1858, p. 5.}
\end{minipage}
\par\vspace{5mm}\hfill{\fontsize{8}{10}\selectfont 26}\hspace{10mm}
\DSource{0651}{202}
\DLine{AB01-PDF0651-body-L001}{}{\textit{fabrayr} (p. 101) et \textit{duǵanbar} (p. 77, 94). — In opere \textit{Vocabulista in arabico pubblicato}}
\DLine{AB01-PDF0651-body-L002}{}{\textit{da C. Schiaparelli,} Firenze 1871, in Hispania meridionali saec. XIII confecto, memorantur}
\DLine{AB01-PDF0651-body-L003}{}{\textit{fibrayr, abrīl, māyuh, yūniyuh, yūliyuh, aghusht, shutanbar, nuwanbar, dhuǵanbar} (1). Hae}
\DLine{AB01-PDF0651-body-L004}{}{\textit{Arabico-occidentales} formae a Latino sermone non per virorum doctorum sed per popularem}
\DLine{AB01-PDF0651-body-L005}{5}{traditionem derivatae sunt; apud agricolas enim Africae boreo-occidentalis et Nigritiae annus}
\DLine{AB01-PDF0651-body-L006}{}{Iulianus et nomina Latina mensium ditioni Romanae superfuerunt ac etiam nunc vigent, sicut}
\DLine{AB01-PDF0651-body-L007}{}{J. \Name{Lippert} animadvertit in Mittheilungen des Seminars für orientalische Sprachen zu Berlin,}
\DLine{AB01-PDF0651-body-L008}{}{t. II, pars 2\textsuperscript{a} (Westasiatische Studien), 1899, p. 253 (2). Qua re semper memorantur in libellis}
\DLine{AB01-PDF0651-body-L009}{}{de astronomia ab Arabibus Algeriensibus et Maroccanis in usum scholarum conscriptis; ita}
\DLine{AB01-PDF0651-body-L010}{10}{\Name{Saḥnūn al-Wānsharishī} (de quo t. I, p. 204), p. 14–15, ubi aperte etiam testatur tale}
\DLine{AB01-PDF0651-body-L011}{}{calendarium ab agricolis adhiberi, et \Name{Muḥammad as-Sūsī al-Marghitī} (3), \textit{al-Mumta‘}}
\DLine{AB01-PDF0651-body-L012}{}{\textit{fī sharḥ al-maqna‘} fasc. III, p. 4 sqq., et VI, p. 2, ubi haec « mensium anni \textit{peregrini} (4)}
\DLine{AB01-PDF0651-body-L013}{}{nomina Romana » ita esse efferenda docet: \textit{yannayr, fabrā’ir} seu \textit{fābrā’ir, mārs} vel \textit{māris}}
\DLine{AB01-PDF0651-body-L014}{}{vel \textit{mars, ibrīl, māyuh} vel \textit{māy, yūniyuh} vel \textit{yunyuh, yūliyuh} vel \textit{yulyuh, aghusht, shu-}}
\DLine{AB01-PDF0651-body-L015}{15}{\textit{tanbir, aktūbir} seu \textit{aktubir, nuwanbir, duǵanbir.} In provincia Maroccana Figīg (\textarabic{فقيق} seu \textarabic{فجيج}),}
\DLine{AB01-PDF0651-body-L016}{}{nuperrime tractui Algeriensi adiecta, haec hodie nomina vigent: \textit{yennayr, fubrayr, mars,}}
\DLine{AB01-PDF0651-body-L017}{}{\textit{ibrīr, māyū, yūlyū, yūlyūz, ghusht, shutānbīr, ktūber, nuwānbīr, duwānbīr} (5). In Ma-}
\DLine{AB01-PDF0651-body-L018}{}{roccano imperio boreali, iuxta J. \Name{Lerchundi}, \textit{Vocabulario español-arábigo del dialecto de}}
\DLine{AB01-PDF0651-body-L019}{}{\textit{Marruecos,} Tánger 1892, nomina sunt: \textit{yinnayr, febrāyar} seu \textit{ibrāyr, mars} seu \textit{mārs,}}
\DLine{AB01-PDF0651-body-L020}{20}{\textit{ebrīl, māyū} seu \textit{māyuh, yunyuh, yulyuh, ghosht, shutanbir, ektūber, nuwanbir, duǵanbir.}}
\DLine{AB01-PDF0651-body-L021}{}{Nomina similia Arabum Algeriensium legantur in lexicis eorum dialecti. — Animadvertendum}
\DLine{AB01-PDF0651-body-L022}{}{est miras formas \textit{ibrīlaz} (simul cum \textit{ibrīl}), \textit{māyb} (cum \textit{māyuh}), \textit{yulyuz} (cum \textit{yulyuh}),}
\DLine{AB01-PDF0651-body-L023}{}{\textit{ghushtaǵǵ} (cum \textit{aghusht}), \textit{aktubiraǵǵ} (cum \textit{aktubir}) et \textit{nuwanbirad} (cum \textit{nuwanbir}) apud}
\par\vspace{6pt}\hrule height.25pt\vspace{5pt}
\noindent\begin{minipage}[t]{.48\textwidth}\vspace{0pt}\fontsize{9}{11.5}\selectfont\setlength{\GutterOffset}{0pt}
\hypertarget{AB01-PDF0651-N01}{}
\DLine{AB01-PDF0651-notes_left-L001}{}{\Indent (1) In Siciliae \textit{Chronico Cantabrigensi} (apud}
\DLine{AB01-PDF0651-notes_left-L002}{}{\Name{Amari}, \textit{Biblioteca arabo-sicula,} textus, Lipsia 1857,}
\DLine{AB01-PDF0651-notes_left-L003}{}{p. 165–76), a Christiano Siculo quodam sub finem}
\DLine{AB01-PDF0651-notes_left-L004}{}{X saec. Chr. Arabice scripto, occurrunt \textit{ǵennāriyūh,}}
\DLine{AB01-PDF0651-notes_left-L005}{}{\textit{mārsuh, abrīl, māyuh, yūniyūh, yūliyūh, awsuh,}}
\DLine{AB01-PDF0651-notes_left-L006}{}{\textit{sutanbar, aktūbar, duǵanbar.} Post Panormum a}
\DLine{AB01-PDF0651-notes_left-L007}{}{Rogerio Comite, Normannorum duce, captam (1072}
\DLine{AB01-PDF0651-notes_left-L008}{}{Chr.), \textit{yannār, marṭiyūs} seu \textit{mārsū, abrīl, māyū,}}
\DLine{AB01-PDF0651-notes_left-L009}{}{\textit{awsuh, sutanbar, aqtūbar, duǵanbar} invenio pas-}
\DLine{AB01-PDF0651-notes_left-L010}{}{sim apud \Name{Cusa}, \textit{I diplomi greci ed arabi di Sici-}}
\DLine{AB01-PDF0651-notes_left-L011}{}{\textit{lia pubblicati} ecc., Palermo 1868-82. Item \textit{māyuh}}
\DLine{AB01-PDF0651-notes_left-L012}{}{et \textit{nuwanbar} in inscriptionibus annorum 1149 et}
\DLine{AB01-PDF0651-notes_left-L013}{}{1153, apud \Name{Amari}, \textit{Le epigrafi arabiche di Sicilia,}}
\DLine{AB01-PDF0651-notes_left-L014}{}{parte II (iscrizioni sepolcrali), Palermo 1879-81, p. 87}
\DLine{AB01-PDF0651-notes_left-L015}{}{et 96 (= Documenti pubbl. a cura della Società Si-}
\DLine{AB01-PDF0651-notes_left-L016}{}{ciliana per la storia patria, serie III, vol. I). — De-}
\DLine{AB01-PDF0651-notes_left-L017}{}{mum \textit{abrīl, māyuh, yūniyuh, yūliyūh} \textarabic{يوليوه}, \textit{shu-}}
\DLine{AB01-PDF0651-notes_left-L018}{}{\textit{tanbar, duǵanbar,} in foederibus inter Pisanos et}
\DLine{AB01-PDF0651-notes_left-L019}{}{dynastas Tuneti, Marocci, Hispaniae, apud \Name{Amari},}
\DLine{AB01-PDF0651-notes_left-L020}{}{\textit{I diplomi arabi del R. Archivio Fiorentino,} Fi-}
\DLine{AB01-PDF0651-notes_left-L021}{}{renze 1863–67, t. I, p. 13, 15, 87, 97, 101, 111, 135,}
\DLine{AB01-PDF0651-notes_left-L022}{}{150, 171, 235, et t. II, p. 8; omnia intra annos 1181}
\DLine{AB01-PDF0651-notes_left-L023}{}{et 1445.}
\end{minipage}
\hfill\begin{minipage}[t]{.48\textwidth}\vspace{0pt}\fontsize{9}{11.5}\selectfont\setlength{\GutterOffset}{0pt}
\hypertarget{AB01-PDF0651-N02}{}
\DLine{AB01-PDF0651-notes_right-L001}{}{\Indent (2) Cfr. quoque E. \Name{Doutté}, \textit{Un texte arabe}}
\DLine{AB01-PDF0651-notes_right-L002}{}{\textit{en dialecte oranais,} Mémoires de la Société de lin-}
\DLine{AB01-PDF0651-notes_right-L003}{}{guistique de Paris, t. XII, 1903, p. 349-50.}
\hypertarget{AB01-PDF0651-N03}{}
\DLine{AB01-PDF0651-notes_right-L004}{}{\Indent (3) De quo t. I, p. 205, adn. 1, et p. \textsc{lxxiii}. Li-}
\DLine{AB01-PDF0651-notes_right-L005}{}{bellus etiam nunc assidue perlegitur in scholis Ara-}
\DLine{AB01-PDF0651-notes_right-L006}{}{bicis urbis Fez seu Fās et regionis Algeriensis; v.}
\DLine{AB01-PDF0651-notes_right-L007}{}{\Name{Delphin}, \textit{L’astronomie au Maroc,} in Journal}
\DLine{AB01-PDF0651-notes_right-L008}{}{Asiatique, ser. VIII, t. XVII, 1891, p. 195–96; \Name{Moty-}}
\DLine{AB01-PDF0651-notes_right-L009}{}{\Name{linski}, p. XI libri a me t. I, p. 297, adn. 4, me-}
\DLine{AB01-PDF0651-notes_right-L010}{}{morati.}
\hypertarget{AB01-PDF0651-N04}{}
\DLine{AB01-PDF0651-notes_right-L011}{}{\Indent (4) \textit{Sanah ‘aǵamiyyah} appellat annum sola-}
\DLine{AB01-PDF0651-notes_right-L012}{}{rem Iulianum, die 1 Ianuarii incipientem, quo agri-}
\DLine{AB01-PDF0651-notes_right-L013}{}{colae utuntur.}
\hypertarget{AB01-PDF0651-N05}{}
\DLine{AB01-PDF0651-notes_right-L014}{}{\Indent (5) Teste E. \Name{Doutté}, \textit{Figuig, notes et im-}}
\DLine{AB01-PDF0651-notes_right-L015}{}{\textit{pressions,} in La Géographie, Bulletin de la Société}
\DLine{AB01-PDF0651-notes_right-L016}{}{de Géographie, Paris 1903, p. 198, qui etiam addit:}
\DLine{AB01-PDF0651-notes_right-L017}{}{« L’année musulmane, en effet, n’est employée à Fi-}
\DLine{AB01-PDF0651-notes_right-L018}{}{« guig que par les lettrés, ou lorsqu’on veut dater}
\DLine{AB01-PDF0651-notes_right-L019}{}{« une époque; en tout autre circonstance, on se sert}
\DLine{AB01-PDF0651-notes_right-L020}{}{« non des mois lunaires, mais des mois du calendrier}
\DLine{AB01-PDF0651-notes_right-L021}{}{« julien, dont les véritables noms ont été conservés}
\DLine{AB01-PDF0651-notes_right-L022}{}{« en berbère depuis l’antiquité, et non pas, comme}
\DLine{AB01-PDF0651-notes_right-L023}{}{« on le crois parfois, réintroduits par nous ».}
\end{minipage}
\DSource{0652}{203}
\DLine{AB01-PDF0652-body-L001}{}{\Name{al-Akhḍarī} et \Name{al-Wānsharīshī} occurrentes, e versibus memorialibus quorumdam scrip-}
\DLine{AB01-PDF0652-body-L002}{}{torum male intellectis ortas esse, ut sunt versus apud \Name{al-Marghītī}, fasc. V, p. 6:}
\par\vspace{1mm}
\DLine{AB01-PDF0652-body-L003}{}{\CenterLine{\makebox[56mm][c]{\textarabic{يَنَّيرَ آ فَابْرَائرَدْ وَمَرْصَدِ}}\hspace{8mm}\makebox[56mm][c]{\textarabic{وَمَنْ يُرِدْ مَدْخَلَ شَهْرٍ يُنْشِدِ}}}}
\DLine{AB01-PDF0652-body-L004}{}{\CenterLine{\makebox[56mm][c]{\textarabic{يلْيُزْ أَغُشْتَجَّ شُتَنْبِرَوْ فُهِ}}\hspace{8mm}\makebox[56mm][c]{\textarabic{أَبْرِيلَ زِ مَايَبْ وَهَاء يُنْيُهِ}}}}
\DLine{AB01-PDF0652-body-L005}{5}{\CenterLine{\makebox[56mm][c]{\textarabic{مَجْمُوعُهَا أَدَدْ زَبَهْ زَجَوْ حَدَوْ}}\hspace{8mm}\makebox[56mm][c]{\textarabic{أَكْتبِرَحْ نُوَنْبِرَدْ دُجَنْبِرَوْ}}}}
\par\vspace{1mm}
\DLine{AB01-PDF0652-body-L006}{}{« qui vult initium mensis recitet: \textit{yannayr-ā, fābrā’ir-ad,} et \textit{marṣ-ad,} | \textit{ibrīl-azi, māy-}}
\DLine{AB01-PDF0652-body-L007}{}{« \textit{-ab,} et \textit{hā’ yunyuh, yulyu-z, aghusht-aǵǵa, shutanbir-aw} pronuntia, | \textit{aktubir-aḥ, nu-}}
\DLine{AB01-PDF0652-body-L008}{}{« \textit{wanbir-ad, duǵanbir-aw;} quarum [consonantium] summa est \textit{adad zabaḥ zaǵaw ḥa-}}
\DLine{AB01-PDF0652-body-L009}{}{« \textit{daw} ». Consonantes nominibus adiectae si suo valore numerali accipiantur (\textit{ā} = 1, \textit{d} = 4,}
\DLine{AB01-PDF0652-body-L010}{10}{\textit{z} = 7 etc.) indicant diem hebdomadae quo mensis ordiri debet.}
\hypertarget{AB01-PDF0652-P01}{}
\DLine{AB01-PDF0652-body-L011}{}{\Indent Nomina « mensium Romanorum » quibus \Name{al-Battānī} et astronomi \Name{orientales} (nec}
\DLine{AB01-PDF0652-body-L012}{}{non \Name{Arabes Christiani} Syriae et Mesopotamiae hoc nostro tempore) utuntur, a \textit{Syriacis}}
\DLine{AB01-PDF0652-body-L013}{}{his repetita sunt: \textit{teshrī qdem} (vel \textit{qdīm}), \textit{teshrī ḥrāy} (\textsyriac{ܐܚܪܳܝ}), \textit{kānūn qdem} (vel \textit{qdīm}), \textit{kānūn}}
\DLine{AB01-PDF0652-body-L014}{}{\textit{ḥrāy, shbāṭ, ādār, nīsān, iyār, ḥzīrān, tāmūz, āb, īlūl.}}
\hypertarget{AB01-PDF0652-P02}{}
\DLine{AB01-PDF0652-body-L015}{15}{\Indent « Romani », \textit{ar-Rūm,} ab Arabibus vocantur incolae omnes imperii Byzantini; eodem}
\DLine{AB01-PDF0652-body-L016}{}{sensu \textgreek{Ῥωμαῖοι} apud Graecos tardae aetatis, Procopium, Zonaram, Cedrenum, Theophanem etc.}
\par\vspace{4mm}
\hypertarget{AB01-PDF0652-H17}{}
\DLine{AB01-PDF0652-body-L017}{}{\CenterLine{\textit{Ad pag. 9–18.} — \textbf{Annorum Arabicorum cum Romanis concordantia.}}}
\par\vspace{3mm}
\hypertarget{AB01-PDF0652-P03}{}
\DLine{AB01-PDF0652-body-L018}{}{\Indent Hae tabulae in fine XXXII capitis (t. I, p. 71) et adp. C (t. I, p. 146) memorantur ac ex-}
\DLine{AB01-PDF0652-body-L019}{}{planantur.}
\hypertarget{AB01-PDF0652-P04}{}
\DLine{AB01-PDF0652-body-L020}{20}{\Indent Parvos codicis errores in numeris dierum facillime emendavi, optimis fretus \textit{Vergleichungs-}}
\DLine{AB01-PDF0652-body-L021}{}{\textit{-Tabellen der muhammedanischen und christlichen Zeitrechnung hrsg. von} F. \Name{Wüsten-}}
\DLine{AB01-PDF0652-body-L022}{}{\Name{feld}, Leipzig 1854. Numeri dierum semper sunt unitate minores quam Wüstenfeldianis, quia}
\DLine{AB01-PDF0652-body-L023}{}{\Name{al-Battānī} initium aerae die 15 Iulii, non 16, constituit.}
\hypertarget{AB01-PDF0652-P05}{}
\DLine{AB01-PDF0652-body-L024}{}{\Indent Duas codicis paginas (f. 163,v. ac 164,r.), concordantiam usque ad annum 720 producen-}
\DLine{AB01-PDF0652-body-L025}{25}{tes, omittendas censui; \textit{spuriae} enim procul dubio sunt. Multis mendis scatent, non ex libra-}
\DLine{AB01-PDF0652-body-L026}{}{riorum incuria sed e manifestis computandi erroribus ortis, ut evenit quotiescumque ineptus}
\DLine{AB01-PDF0652-body-L027}{}{interpolator aliquid Albatenio addere conatur. Concordantiam usque ad 600 hegirae, id est}
\DLine{AB01-PDF0652-body-L028}{}{per amplius 300 annos a tabularum epocha duxisse, sat \Name{Albatenio} erat. — In versione}
\DLine{AB01-PDF0652-body-L029}{}{Hispanica, f. 30,r.–35,v., hae tabulae inabsolutae reperiuntur; nam in dextera sectione f. 35,v.}
\DLine{AB01-PDF0652-body-L030}{30}{anni hegirae usque ad 488 scribuntur (7 ultimis lineis vacuis exstantibus), nomina autem primi}
\DLine{AB01-PDF0652-body-L031}{}{diei al-muḥarram et anni Dhū ’l-qarnayn (quorum ultimus 1492) nonnisi usque ad lineam}
\DLine{AB01-PDF0652-body-L032}{}{anni 477 heg., dies demum ac menses Romani usque ad lineam anni 470 heg. tantum. Tabu-}
\DLine{AB01-PDF0652-body-L033}{}{lam amanuensem non absolvisse eo puto, quod hac pagina ubicumque in centenis annorum}
\DLine{AB01-PDF0652-body-L034}{}{hegirae falso 400 pro 500 scripsisse animadvertisset, praetereaque loco omissi signi al-muḥar-}
\DLine{AB01-PDF0652-body-L035}{35}{ram anni 452 (leg. 552) posuisse signum anni 453 (leg. 553), et extra sua loca signa sequentia}
\DLine{AB01-PDF0652-body-L036}{}{constituisse.}
\par\vspace{4mm}
\hypertarget{AB01-PDF0652-H37}{}
\DLine{AB01-PDF0652-body-L037}{}{\CenterLine{\textit{Ad pag. 19-23.} — \textbf{Medii motus Solis et Lunae ad calendarium Arabicum.}}}
\par\vspace{3mm}
\hypertarget{AB01-PDF0652-P06}{}
\DLine{AB01-PDF0652-body-L038}{}{\Indent De his tabulis agitur cap. XXVII et XXXIII (t. I, p. 42 et 71).}
\hypertarget{AB01-PDF0652-P07}{}
\DLine{AB01-PDF0652-body-L039}{}{\Indent In columnarum titulis codex semper \textit{al-masīr al-awsaṭ} « iter medium » habet, quod}
\DLine{AB01-PDF0652-body-L040}{40}{nonnisi p. 20–23 convenit, nam in iis \textit{motus medii,} scilicet gradus quos Sol, Luna etc. dato}
\DSource{0653}{204}
\DLine{AB01-PDF0653-body-L001}{}{temporis intervallo motu medio conficiunt, praebentur. At p. 19 indicantur \textit{longitudines me-}}
\DLine{AB01-PDF0653-body-L002}{}{\textit{diae} quas Sol, Luna etc. ineuntibus annis I, XXXI etc. hegirae tenebant; numeri igitur sunt}
\DLine{AB01-PDF0653-body-L003}{}{\textit{epochae} seu \textit{radices} mediorum motuum, a quibus postea longitudines pro quocumque tem-}
\DLine{AB01-PDF0653-body-L004}{}{poris momento deducuntur. Qua re titulos columnarum p. 19, Arabice sicut in codice exsta-}
\DLine{AB01-PDF0653-body-L005}{5}{bant servatos, paululum in versione mutavi. Cfr. quoque adnotationes ad pag. 24–28, 72–77,}
\DLine{AB01-PDF0653-body-L006}{}{102–103.}
\hypertarget{AB01-PDF0653-P01}{}
\DLine{AB01-PDF0653-body-L007}{}{\Indent Motus \textit{nodi} retrogradus est; sed in tabulis \Name{Albatenianis} ita supputatur ac si esset}
\DLine{AB01-PDF0653-body-L008}{}{directus, qua re ipse auctor t. I, p. 75–76 (cap. XXXVII) iubet motum nodi e tabulis sum-}
\DLine{AB01-PDF0653-body-L009}{}{ptum semper de 360° demere.}
\hypertarget{AB01-PDF0653-P02}{}
\DLine{AB01-PDF0653-body-L010}{10}{\Indent \textbf{\textsf{Pag. 19}}, sub annis 661 et sqq., facile fuit gradus anomaliae Lunae restituere, in codice}
\DLine{AB01-PDF0653-body-L011}{}{omissos.}
\hypertarget{AB01-PDF0653-P03}{}
\DLine{AB01-PDF0653-body-L012}{}{\Indent \textbf{\textsf{Pag. 21}}, cod. nonnisi dhū ’l-ḥiǵǵah annorum communium praebet, ut p. 74 subāṭ (Februa-}
\DLine{AB01-PDF0653-body-L013}{}{rium); nam tabulae Albatenianae pro annis collectis et singulis ita conditae sunt, ut in com-}
\DLine{AB01-PDF0653-body-L014}{}{putando nunquam opus sit bisextiles a communibus annis distinguere. Suadente tamen \Name{Schia-}}
\DLine{AB01-PDF0653-body-L015}{15}{\Name{parelli}, dhū ’l-ḥiǵǵah bisextilium addidi.}
\hypertarget{AB01-PDF0653-P04}{}
\DLine{AB01-PDF0653-body-L016}{}{\Indent \textbf{\textsf{Pag. 22}}, duo codicis errores, qui me effugerunt, emendentur: secundae itineris Solis in}
\DLine{AB01-PDF0653-body-L017}{}{10 diebus (ubi cod. et versio 25″ \textarabic{كه} loco 24″ \textarabic{كد}), et secundae nodi in 21 diebus (ubi 45″ \textarabic{مه}}
\DLine{AB01-PDF0653-body-L018}{}{pro 44″ \textarabic{مد}). — Eadem tabula p. 75, cum tribus discrepantiis ab ipso \Name{al-Battānī} provenien-}
\DLine{AB01-PDF0653-body-L019}{}{tibus, nempe 15″ in motu Solis in novem diebus, 23″ in decem, et 38′ 59″ (pro 39′ 0″) in}
\DLine{AB01-PDF0653-body-L020}{20}{anomalia Lunae in decem diebus.}
\hypertarget{AB01-PDF0653-P05}{}
\DLine{AB01-PDF0653-body-L021}{}{\Indent \textbf{\textsf{Pag. 23}}, duo codicis errores, qui in versionem irrepserunt, emendentur: 25″ \textarabic{كه} cod. pro}
\DLine{AB01-PDF0653-body-L022}{}{24″ \textarabic{كد} in Solis itinere trium horarum, et 13″ \textarabic{يج} (cod. \textarabic{لج} 33″) pro 14″ \textarabic{يد} in anomalia lunari}
\DLine{AB01-PDF0653-body-L023}{}{23 horarum. — Eadem tabula p. 76, ubi in motu lunari quatuordecim horarum 11″ loco 10″,}
\DLine{AB01-PDF0653-body-L024}{}{et in anomalia decem horarum 37″ loco 38″ (exacte 5° 26′ 37″ 28‴ 17\textsuperscript{IV}....); ambae discre-}
\DLine{AB01-PDF0653-body-L025}{25}{pantiae ab \Name{al-Battānī} provenire videntur.}
\par\vspace{8mm}
\hypertarget{AB01-PDF0653-H26}{}
\DLine{AB01-PDF0653-body-L026}{}{\CenterLine{\textit{Ad pag. 24-28.} — \textbf{Motus medii quinque planetarum}}}
\hypertarget{AB01-PDF0653-H27}{}
\DLine{AB01-PDF0653-body-L027}{}{\CenterLine{\textbf{ad calendarium Arabicum.}}}
\par\vspace{3mm}
\hypertarget{AB01-PDF0653-P06}{}
\DLine{AB01-PDF0653-body-L028}{}{\Indent Tabulas memorat auctor t. I, p. 71 et 113 (exeunte XXXIII et ineunte XLV cap.).}
\hypertarget{AB01-PDF0653-P07}{}
\DLine{AB01-PDF0653-body-L029}{}{\Indent Columnae p. 24 Arabice ut p. 25–26 inscriptae; paulo aliter in versione, ut in adn. ad}
\DLine{AB01-PDF0653-body-L030}{30}{p. 19–23 monui. — P. 27 et 28 necessario a p. 106 et 105 non differunt.}
\hypertarget{AB01-PDF0653-P08}{}
\DLine{AB01-PDF0653-body-L031}{}{\Indent In tabulis suis \Name{Ptolemaeus} IX, 4 (ed. Halma, t. II, p. 126–55) anomalias planetarum}
\DLine{AB01-PDF0653-body-L032}{}{superiorum et motus medios longitudinis inferiorum addit. Quae \Name{al-Battānī} omittit, iure}
\DLine{AB01-PDF0653-body-L033}{}{t. I, p. 113 (cap. XLV) animadvertens in systemate Ptolemaico medium motum anomaliae trium}
\DLine{AB01-PDF0653-body-L034}{}{superiorum eumdem esse quem motum medium longitudinis Solis imminutum medio motu ano-}
\DLine{AB01-PDF0653-body-L035}{35}{maliae ipsorum planetarum; atque medium motum longitudinis Veneris ac Mercurii non differre}
\DLine{AB01-PDF0653-body-L036}{}{a medio Solis motu.}
\hypertarget{AB01-PDF0653-P09}{}
\DLine{AB01-PDF0653-body-L037}{}{\Indent De parvis discrepantiis in motuum quantitatibus a \Name{Ptolemaeo} vide t. I, p. 66, et adnota-}
\DLine{AB01-PDF0653-body-L038}{}{tiones meas t. I, p. 239, 242.}
\hypertarget{AB01-PDF0653-P10}{}
\DLine{AB01-PDF0653-body-L039}{}{\Indent Folium 167 codicis, in versione omissum, erat additamentum nihil ad tabulas Albatenia-}
\DLine{AB01-PDF0653-body-L040}{40}{nas attinens.}
\DSource{0654}{205}
\par\vspace{4mm}
\hypertarget{AB01-PDF0654-H01}{}
\DLine{AB01-PDF0654-body-L001}{}{\CenterLine{\textit{Ad pag. 29-32.} — \textbf{Syzygiae mediae ad annos Aegyptios.}}}
\par\vspace{3mm}
\hypertarget{AB01-PDF0654-P01}{}
\DLine{AB01-PDF0654-body-L002}{}{\Indent Tabularum usum partim illustravit auctor cap. XLII (t. I, p. 92–93).}
\hypertarget{AB01-PDF0654-P02}{}
\DLine{AB01-PDF0654-body-L003}{}{\Indent In his tabulis miro modo anni \textit{Aegyptii} (seu \textit{vagi,} sine intercalatione Iuliana) coniuncti}
\DLine{AB01-PDF0654-body-L004}{}{apparent cum aera Dhū ’l-qarnayn, quae alibi semper per annos \textit{fixos Iulianos} computatur.}
\DLine{AB01-PDF0654-body-L005}{5}{Tabulis attente perspectis certissime patet hos annos vagos a Dhū ’l-qarnayn numerari, sed 1}
\DLine{AB01-PDF0654-body-L006}{}{thōth i. e. anni initium semper cum 1 thōth congruere annorum Nabonassaris, a \Name{Ptolemaeo}}
\DLine{AB01-PDF0654-body-L007}{}{in \textit{Almagesto} adhibitorum; qua re 1 thōth primi anni a Dhū ’l-qarnayn fuit 1 thōth 437}
\DLine{AB01-PDF0654-body-L008}{}{Nabonassaris, 9 Novembris 312 ante Chr. n., id est 69 diebus post initium aerae Dhū ’l-qarnayn}
\DLine{AB01-PDF0654-body-L009}{}{iuxta usum Albatenianum. Manifeste hanc tabulam \Name{al-Battānī} in usum eorum condidit qui}
\DLine{AB01-PDF0654-body-L010}{10}{annis \textit{Aegyptiis vagis} Nabonassaris cum \Name{Ptolemaeo} uti vellent; sed cur eos maluerit a Dhū}
\DLine{AB01-PDF0654-body-L011}{}{’l-qarnayn numerare nescio.}
\hypertarget{AB01-PDF0654-P03}{}
\DLine{AB01-PDF0654-body-L012}{}{\Indent Confusio etiam oriri potest e nomine « \textit{mensium Coptorum} » quod p. 31 occurrit. Nam,}
\DLine{AB01-PDF0654-body-L013}{}{iuxta t. I, p. 67, l. 19–20, \textit{annus Copticus} est fixus, eiusque 1 thōth (I, p. 68, l. 23 sqq., et}
\DLine{AB01-PDF0654-body-L014}{}{p. 69, l. 7–21) semper in 29 Augusti incidit; praeterea, sicut in adn. t. I, p. 244–45 exposui,}
\DLine{AB01-PDF0654-body-L015}{15}{Albatenio \textit{aera Coptorum} est aera quaedam Augusti vel Alexandrina, a \Name{Theone} quoque adhi-}
\DLine{AB01-PDF0654-body-L016}{}{bita, annis fixis praedita, cuius initium fuit 29 Aug. 25 ante Chr. n., seu 287 annis, 3 diebus de-}
\DLine{AB01-PDF0654-body-L017}{}{tractis, post aeram Dhū ’l-qarnayn. Hic tamen sunt menses annorum \textit{vagorum,} quos miro modo}
\DLine{AB01-PDF0654-body-L018}{}{etiam t. I, p. 70, l. 30 sqq. \textit{annos Coptorum} vocat, addens eos « annos esse Aegyptios a Dhū}
\DLine{AB01-PDF0654-body-L019}{}{« ’l-qarnayn »; eodemque sensu in praeceptis ad syzygias supputandas (t. I, p. 92–93) accipiuntur.}
\hypertarget{AB01-PDF0654-P04}{}
\DLine{AB01-PDF0654-body-L020}{20}{\Indent Animadvertendum etiam est diem 29 Aug. 25 ante Chr. n., aerae Coptorum seu Augusti}
\DLine{AB01-PDF0654-body-L021}{}{initium, esse primum diem thōth anni 724 Nabonassaris, 300 Aridaei seu Alexandri mortui,}
\DLine{AB01-PDF0654-body-L022}{}{ac initium anni 287 a Dhū ’l-qarnayn iuxta Copticum incipiendi morem (a 29 Aug.). Qua re}
\DLine{AB01-PDF0654-body-L023}{}{t. I, p. 70 et 92–93, ad annos vagos Coptorum, a Dhū ’l-qarnayn numeratos, convertendos}
\DLine{AB01-PDF0654-body-L024}{}{in annos fixos, iubet \Name{al-Battānī} ab annis datis detrahere 287, ac residuum per 4 dividere.}
\DLine{AB01-PDF0654-body-L025}{25}{Re vera initium anni 287 Dhū ’l-qarnayn, \textit{Coptico} ordiendi more, cadit 29 Augusti tam iuxta}
\DLine{AB01-PDF0654-body-L026}{}{annos vagos quam fixos, ut etiam tribus sequentibus annis; quod postea nonnisi quadriennio}
\DLine{AB01-PDF0654-body-L027}{}{rursus eveniet cuius primus annus 287 + 365 × 4 = 1747 Dhū ’l-qarnayn. Quem annum ta-}
\DLine{AB01-PDF0654-body-L028}{}{bulae \Name{Albatenianae} non assequuntur.}
\hypertarget{AB01-PDF0654-P05}{}
\DLine{AB01-PDF0654-body-L029}{}{\Indent Quinque tabularum columnae haec praebent:}
\hypertarget{AB01-PDF0654-P06}{}
\DLine{AB01-PDF0654-body-L030}{30}{\Indent I. Annos aut menses datos, cum quibus est in tabulas intrandum.}
\hypertarget{AB01-PDF0654-P07}{}
\DLine{AB01-PDF0654-body-L031}{}{\Indent II. Dies primi mensis Aegyptii (thōth), et partes dierum sexagesimas, quibus dati anni}
\DLine{AB01-PDF0654-body-L032}{}{prima coniunctio aut oppositio media evenit.}
\hypertarget{AB01-PDF0654-P08}{}
\DLine{AB01-PDF0654-body-L033}{}{\Indent III. Mediam longitudinem Solis, temporis momento in II indicato.}
\hypertarget{AB01-PDF0654-P09}{}
\DLine{AB01-PDF0654-body-L034}{}{\Indent IV. Anomaliam mediam Lunae eodem temporis momento.}
\hypertarget{AB01-PDF0654-P10}{}
\DLine{AB01-PDF0654-body-L035}{35}{\Indent V. Argumentum latitudinis Lunae seu distantiam a nodo boreali.}
\hypertarget{AB01-PDF0654-P11}{}
\DLine{AB01-PDF0654-body-L036}{}{\Indent \Name{Ptolemaeus} (1) in III non mediam longitudinem Solis exhibet, at Solis distantiam (\textgreek{ἀποχή})}
\DLine{AB01-PDF0654-body-L037}{}{ab apogeo, quod perperam in 65° 30′ fixum arbitrabatur; in V argumentum latitudinis Lunae}
\DLine{AB01-PDF0654-body-L038}{}{a limite boreali (= longit. nodi + 90°) numeratum. In reliquis constructio tabularum apud}
\DLine{AB01-PDF0654-body-L039}{}{utrumque aequalis; sed elementa pro annis collectis aliquantulo discrepant ob motum Solis}
\DLine{AB01-PDF0654-body-L040}{40}{in tabulis Albatenianis diversum et ob longitudinis intervallum inter urbem Alexandriam et}
\DLine{AB01-PDF0654-body-L041}{}{ar-Raqqah. Elementa sua \Name{Ptolemaeus} VI, 2 (ed. Halma, t. I, p. 374–78) exposuit, quae vide}
\par\vspace{6pt}\hrule height.25pt\vspace{5pt}
\noindent\begin{minipage}[t]{.48\textwidth}\vspace{0pt}\fontsize{9}{11.5}\selectfont\setlength{\GutterOffset}{0pt}
\hypertarget{AB01-PDF0654-N01}{}
\DLine{AB01-PDF0654-notes_left-L001}{}{\Indent (1) Cuius tabulae in \textit{Almag.} VI, 3 (ed. Halma,}
\DLine{AB01-PDF0654-notes_left-L002}{}{t. I, p. 378-83); earum usus, et conversio syzygia-}
\end{minipage}
\hfill\begin{minipage}[t]{.48\textwidth}\vspace{0pt}\fontsize{9}{11.5}\selectfont\setlength{\GutterOffset}{0pt}
\DLine{AB01-PDF0654-notes_right-L001}{}{rum mediarum in veras postea docetur VI, 4 (Hal-}
\DLine{AB01-PDF0654-notes_right-L002}{}{ma, I, 384-87).}
\end{minipage}
\DSource{0655}{206}
\DLine{AB01-PDF0655-body-L001}{}{etiam apud \Name{Delambre}, \textit{Hist. de l’astron. ancienne,} t. II, p. 223–24; sua \Name{al-Battānī} non}
\DLine{AB01-PDF0655-body-L002}{}{indicat. Ut discrepantiae quantitas cognoscatur, primam coniunctionem anni \textit{vagi} 915 Dhū}
\DLine{AB01-PDF0655-body-L003}{}{’l-qarnayn, 1351 Nabonassaris supputo iuxta \textit{Almagestum.} Ultimus annus tabulae Ptolemaei}
\DLine{AB01-PDF0655-body-L004}{}{est 1101 Nabonassaris;}
\par\vspace{1mm}
\DLine{AB01-PDF0655-body-L005}{5}{\hspace*{30mm}\makebox[74mm][l]{prima coniunctio anni 1101 Nabonassaris}\makebox[7mm][c]{}\makebox[7mm][r]{22}\rlap{\textsuperscript{\textit{d}}}\makebox[7mm][r]{41}\rlap{′}\makebox[7mm][r]{45}\rlap{″}\hspace{2.5mm}thōth}
\DLine{AB01-PDF0655-body-L006}{}{\hspace*{30mm}\makebox[74mm][l]{variatio intervallo 250 annorum vagorum}\makebox[7mm][c]{—}\makebox[7mm][r]{0}\rlap{}\makebox[7mm][r]{27}\rlap{}\makebox[7mm][r]{51}\rlap{}}
\par\nointerlineskip\vspace{1pt}\noindent\hspace*{111mm}\rule{22mm}{.4pt}\par\nointerlineskip\vspace{2pt}
\DLine{AB01-PDF0655-body-L007}{}{\hspace*{30mm}\makebox[74mm][l]{pro 1351 Nabonassaris, 915 Dhū ’l-qarn.}\makebox[7mm][c]{}\makebox[7mm][r]{22}\rlap{\textsuperscript{\textit{d}}}\makebox[7mm][r]{13}\rlap{′}\makebox[7mm][r]{54}\rlap{″}\hspace{2.5mm}thōth}
\par\vspace{1mm}
\hypertarget{AB01-PDF0655-P01}{}
\DLine{AB01-PDF0655-body-L008}{}{\Indent T. I, pag. 42, cap. XXVII, differentiam horariam inter Alexandriam et ar-Raqqah statuit}
\DLine{AB01-PDF0655-body-L009}{}{al-Battānī « circiter $\qfrac{2}{3}$ horae aequinoctialis », seu 10° seu 0\textsuperscript{\textit{d}} 1′ 40″; igitur}
\par\vspace{1mm}
\DLine{AB01-PDF0655-body-L010}{10}{\hspace*{30mm}\makebox[74mm][l]{}\makebox[7mm][c]{}\makebox[7mm][r]{22}\rlap{\textsuperscript{\textit{d}}}\makebox[7mm][r]{13}\rlap{′}\makebox[7mm][r]{54}\rlap{″}}
\DLine{AB01-PDF0655-body-L011}{}{\hspace*{30mm}\makebox[74mm][l]{}\makebox[7mm][c]{+}\makebox[7mm][r]{0}\rlap{}\makebox[7mm][r]{1}\rlap{}\makebox[7mm][r]{40}\rlap{}}
\par\nointerlineskip\vspace{1pt}\noindent\hspace*{111mm}\rule{22mm}{.4pt}\par\nointerlineskip\vspace{2pt}
\DLine{AB01-PDF0655-body-L012}{}{\hspace*{30mm}\makebox[74mm][l]{pro meridiano ar-Raqqah iuxta \Name{Ptolem.}}\makebox[7mm][c]{}\makebox[7mm][r]{22}\rlap{\textsuperscript{\textit{d}}}\makebox[7mm][r]{15}\rlap{′}\makebox[7mm][r]{34}\rlap{″}}
\DLine{AB01-PDF0655-body-L013}{}{\hspace*{30mm}\makebox[74mm][l]{iuxta tabulam \Name{Albatenii}}\makebox[7mm][c]{}\makebox[7mm][r]{22}\rlap{}\makebox[7mm][r]{14}\rlap{}\makebox[7mm][r]{44}\rlap{}}
\par\nointerlineskip\vspace{1pt}\noindent\hspace*{111mm}\rule{22mm}{.4pt}\par\nointerlineskip\vspace{2pt}
\DLine{AB01-PDF0655-body-L014}{}{\hspace*{30mm}\makebox[74mm][r]{differentia a \Name{Ptolemaeo}}\makebox[7mm][c]{—}\makebox[7mm][r]{0}\rlap{\textsuperscript{\textit{d}}}\makebox[7mm][r]{0}\rlap{′}\makebox[7mm][r]{50}\rlap{″}\hspace{2.5mm}= — 0\textsuperscript{\textit{h}} 20\textsuperscript{\textit{m}}.}
\par\vspace{1mm}
\hypertarget{AB01-PDF0655-P02}{}
\DLine{AB01-PDF0655-body-L015}{15}{\Indent Si differentiam longitudinis inter Alexandriam et ar-Raqqah sumimus 12° 45′ = 0\textsuperscript{\textit{h}} 51\textsuperscript{\textit{m}} =}
\DLine{AB01-PDF0655-body-L016}{}{= 0\textsuperscript{\textit{d}} 2′ 7″.5, ut a tabulis geographicis Albatenianis (t. II, nr. 109 ac 150) deducitur, differentia}
\DLine{AB01-PDF0655-body-L017}{}{in coniunctione fit — 0\textsuperscript{\textit{d}} 1′ 17″.5 = — 0\textsuperscript{\textit{h}} 31\textsuperscript{\textit{m}}.}
\hypertarget{AB01-PDF0655-P03}{}
\DLine{AB01-PDF0655-body-L018}{}{\Indent \textbf{\textsf{Pag. 29-30}}, pro \textit{annis collectis.} Periodos 25 annorum delegit, \Name{Ptolemaei} exemplum se-}
\DLine{AB01-PDF0655-body-L019}{}{quens, quia 25 anni vagi (seu 9125 dies) 309 lunationes fere exacte complectuntur; differentia}
\DLine{AB01-PDF0655-body-L020}{20}{ab illis 9125 diebus demenda nonnisi 0\textsuperscript{\textit{d}} 2′ 47″ 5‴ seu 0\textsuperscript{\textit{d}} 1\textsuperscript{\textit{h}} 6\textsuperscript{\textit{m}} 50\textsuperscript{\textit{s}} attingit, qua re per multas}
\DLine{AB01-PDF0655-body-L021}{}{periodos 25 annorum vagorum una tantum coniunctio unaque oppositio mensi thōth primi anni}
\DLine{AB01-PDF0655-body-L022}{}{periodi erit, taliumque coniunctionum aut oppositionum tempora paululo ab alteris different. —}
\DLine{AB01-PDF0655-body-L023}{}{Oppositionum tabula construitur, a diebus coniunctionum dimidiam lunationem seu 14\textsuperscript{\textit{d}} 45′ 55″}
\DLine{AB01-PDF0655-body-L024}{}{demendo, motusque huic lunationi respondentes a motibus coniunctionum.}
\hypertarget{AB01-PDF0655-P04}{}
\DLine{AB01-PDF0655-body-L025}{25}{\Indent \textbf{\textsf{Pag. 31}}. — \textbf{[}Trium tabellarum quae hac pagina continentur, una, \textit{mensium} scilicet, est}
\DLine{AB01-PDF0655-body-L026}{}{necessaria et in \textit{Almagesto} quoque (VI, 3, ed. Halma t. I, p. 383) exhibetur. — \textit{Prima} inutilis;}
\DLine{AB01-PDF0655-body-L027}{}{quis enim per quadrantes mensis Aegyptii seu Coptici, i. e. per 7 $\qfrac{1}{2}$ dies unquam computabit? —}
\DLine{AB01-PDF0655-body-L028}{}{\textit{Ultima,} pro annorum intervallis a quolibet temporis momento numeratis, parvae est utilitatis;}
\DLine{AB01-PDF0655-body-L029}{}{praeterea ultima eius columna manifeste in codice ab imperito lectore suppleta erat, nam motus}
\DLine{AB01-PDF0655-body-L030}{30}{latitudinis Lunae supputabatur quasi 108° 47′ 19″ $\qfrac{1}{2}$ esset intervallo 100 annorum vagorum,}
\DLine{AB01-PDF0655-body-L031}{}{loco 108° 48′ 23″ $\qfrac{1}{2}$ quae e tabulis Albatenianis p. 29–30 eliciuntur. Eam denuo supputavi ut}
\DLine{AB01-PDF0655-body-L032}{}{cum Albatenii elementis congrueret. — Primam ultimamque tabellam \textit{spurias} esse suspicor;}
\DLine{AB01-PDF0655-body-L033}{}{p. 87 tamen similes exstant (1). Tituli columnarum ultimae significant pro tempore futuro nu-}
\DLine{AB01-PDF0655-body-L034}{}{meros dierum esse subtrahendos, motuum autem addendos. — \textsc{Schiaparelli}.\textbf{]} — Ultima linea}
\DLine{AB01-PDF0655-body-L035}{35}{primae tabellae pro « monsis » legendum « mensis ». Tabellae prima et secunda eaedem sunt}
\DLine{AB01-PDF0655-body-L036}{}{quae p. 87, cum duabus parvis discrepantiis: sub $\qfrac{1}{4}$ mensis hic 57″, illic 58″ in diebus; sub}
\DLine{AB01-PDF0655-body-L037}{}{III mense (ayyār) hic 1″, illic 0″ in anomalia Lunae.}
\hypertarget{AB01-PDF0655-P05}{}
\DLine{AB01-PDF0655-body-L038}{}{\Indent \textbf{\textsf{Pag. 32}}, pro \textit{annis singulis.} Ad amussim cum \Name{Ptolemaica} VI, 3 (ed. Halma t. I, p. 382)}
\DLine{AB01-PDF0655-body-L039}{}{congruit; eius supputatio facillima. Excessus 13 lunationem super 365 dies anni vagi est 18\textsuperscript{\textit{d}}}
\DLine{AB01-PDF0655-body-L040}{40}{53′ 51″ 43‴; ergo si A vocemus dies mensis thōth ac diei sexagesimas respondentes primo}
\DLine{AB01-PDF0655-body-L041}{}{anni cuiusdam novilunio (aut plenilunio), anno sequenti primum novilunium (aut plenilunium)}
\par\vspace{6pt}\hrule height.25pt\vspace{5pt}{\fontsize{9}{11.5}\selectfont
\hypertarget{AB01-PDF0655-N01}{}
\DLine{AB01-PDF0655-notes-L001}{}{\Indent (1) Item in Hispanica versione f. 41,v. et 56,v.}
}
\DSource{0656}{207}
\DLine{AB01-PDF0656-body-L001}{}{eveniet tempore A + 18\textsuperscript{\textit{d}} 53′ 52″. Quod, si lunationem integram seu 29\textsuperscript{\textit{d}} 31′ 50″ superaverit,}
\DLine{AB01-PDF0656-body-L002}{}{erit secundum mensis thōth novilunium, aut (30 dies excedente) novilunium mensis sequentis;}
\DLine{AB01-PDF0656-body-L003}{}{tunc ad primam coniunctionem aut oppositionem anni reperiendam oportebit de summa luna-}
\DLine{AB01-PDF0656-body-L004}{}{tionem demere ac de motibus columnarum reliquarum motus lunationi respondentes. — In}
\DLine{AB01-PDF0656-body-L005}{5}{columna dierum emendanda sunt: 40″ pro 50″ sub anno 6, 39″ pro 29″ sub 11, 24″ pro 34″}
\DLine{AB01-PDF0656-body-L006}{}{sub 14, 23″ pro 13″ sub 19.}
\hypertarget{AB01-PDF0656-P01}{}
\DLine{AB01-PDF0656-body-L007}{}{\Indent \textit{Termini eclipsium} (\textgreek{ὅροι ἐκλειπτικοί}). — Eos « in pagina mensium syzygiarum » indica-}
\DLine{AB01-PDF0656-body-L008}{}{visse \Name{al-Battānī} affirmat cap. XLII et XLIV (t. I, p. 96 ac 104), et eo quidem loco in ver-}
\DLine{AB01-PDF0656-body-L009}{}{sione Hispanica (f. 41,v. et 56,v.) termini occurrunt; sed codex Escurialensis eos praebet}
\DLine{AB01-PDF0656-body-L010}{10}{hic sub tabula annorum singulorum et pag. 88. Suos \Name{Ptolemaeus} in tabula syzygiarum pro}
\DLine{AB01-PDF0656-body-L011}{}{annis singulis dedit, eorum supputationem postea docens (VI, 5, ed. Halma t. I, p. 388–96); ob}
\DLine{AB01-PDF0656-body-L012}{}{diversam Solis aequationem diversasque diametros angulares luminarium et umbrae necessario}
\DLine{AB01-PDF0656-body-L013}{}{\textit{Ptolemaici} ab \textit{Albatenianis} differunt. Cum nullo loco \Name{al-Battānī} doceat quo modo terminos}
\DLine{AB01-PDF0656-body-L014}{}{supputaverit, numeros codicis qui ad \textit{eclipses solares} spectant, manifestissimis mendis affectos,}
\DLine{AB01-PDF0656-body-L015}{15}{proxime emendare studueram, numeros deligens e quibus facilius amanuensium errores orti vi-}
\DLine{AB01-PDF0656-body-L016}{}{derentur, et proportiones terminorum Almagesti quantum fieri poterat servans. Correctionum}
\DLine{AB01-PDF0656-body-L017}{}{conspectum praebeo, primo numeros codicis hac pagina 32 (f. 172,r.) contentos indicans, deinde}
\DLine{AB01-PDF0656-body-L018}{}{numeros codicis pag. 88 (f. 194,v.), tertio numeros receptos, quarto errorum explicationem:}
\par\vspace{2mm}
\DLine{AB01-PDF0656-body-L019}{}{\hspace*{34.0mm}\makebox[5.0mm][l]{\smash{\raisebox{-6.9pt}{$\left\{\rule[-10pt]{0pt}{20pt}\right.$}}}\makebox[12.0mm][r]{\textit{154}\rlap{\textit{°}}}\makebox[9.0mm][r]{44\rlap{′}}\hspace{3mm}\rule[-\dp\strutbox]{.4pt}{\baselineskip}\hspace{2mm}\makebox[12.0mm][r]{159\rlap{°}}\makebox[9.0mm][r]{44\rlap{′}}\hspace{3mm}\rule[-\dp\strutbox]{.4pt}{\baselineskip}\hspace{2mm}\makebox[12.0mm][r]{159\rlap{°}}\makebox[9.0mm][r]{44\rlap{′}}\hspace{3mm}\rule[-\dp\strutbox]{.4pt}{\baselineskip}\hspace{2mm}\makebox[60.0mm][l]{\smash{\textarabic{قند} pro \textarabic{قنط}}}}
\DLine{AB01-PDF0656-body-L020}{20}{\hspace*{34.0mm}\makebox[5.0mm][l]{\smash{}}\makebox[12.0mm][r]{\smash{\textit{190}}}\makebox[9.0mm][r]{\smash{16}}\hspace{3mm}\rule[-\dp\strutbox]{.4pt}{\baselineskip}\hspace{2mm}\makebox[12.0mm][r]{\smash{\textit{167}}}\makebox[9.0mm][r]{\smash{16}}\hspace{3mm}\rule[-\dp\strutbox]{.4pt}{\baselineskip}\hspace{2mm}\makebox[12.0mm][r]{\smash{191}}\makebox[9.0mm][r]{\smash{16}}\hspace{3mm}\rule[-\dp\strutbox]{.4pt}{\baselineskip}\hspace{2mm}\makebox[60.0mm][l]{\smash{\textarabic{قض} et \textarabic{قصز} pro \textarabic{قضا}}}}
\DLine{AB01-PDF0656-body-L021}{}{\hspace*{34.0mm}\makebox[5.0mm][l]{\smash{\raisebox{-6.9pt}{$\left\{\rule[-10pt]{0pt}{20pt}\right.$}}}\makebox[12.0mm][r]{\smash{0}}\makebox[9.0mm][r]{\smash{0}}\hspace{3mm}\rule[-\dp\strutbox]{.4pt}{\baselineskip}\hspace{2mm}\makebox[12.0mm][r]{\smash{0}}\makebox[9.0mm][r]{\smash{0}}\hspace{3mm}\rule[-\dp\strutbox]{.4pt}{\baselineskip}\hspace{2mm}\makebox[12.0mm][r]{\smash{0}}\makebox[9.0mm][r]{\smash{0}}\hspace{3mm}\rule[-\dp\strutbox]{.4pt}{\baselineskip}\hspace{2mm}\makebox[60.0mm][l]{\smash{}}}
\DLine{AB01-PDF0656-body-L022}{}{\hspace*{34.0mm}\makebox[5.0mm][l]{\smash{}}\makebox[12.0mm][r]{\smash{20}}\makebox[9.0mm][r]{\smash{16}}\hspace{3mm}\rule[-\dp\strutbox]{.4pt}{\baselineskip}\hspace{2mm}\makebox[12.0mm][r]{\smash{20}}\makebox[9.0mm][r]{\smash{16}}\hspace{3mm}\rule[-\dp\strutbox]{.4pt}{\baselineskip}\hspace{2mm}\makebox[12.0mm][r]{\smash{20}}\makebox[9.0mm][r]{\smash{16}}\hspace{3mm}\rule[-\dp\strutbox]{.4pt}{\baselineskip}\hspace{2mm}\makebox[60.0mm][l]{\smash{}}}
\DLine{AB01-PDF0656-body-L023}{}{\hspace*{34.0mm}\makebox[5.0mm][l]{\smash{\raisebox{-6.9pt}{$\left\{\rule[-10pt]{0pt}{20pt}\right.$}}}\makebox[12.0mm][r]{\smash{348}}\makebox[9.0mm][r]{\smash{\textit{57}}}\hspace{3mm}\rule[-\dp\strutbox]{.4pt}{\baselineskip}\hspace{2mm}\makebox[12.0mm][r]{\smash{\textit{345}}}\makebox[9.0mm][r]{\smash{\textit{53}}}\hspace{3mm}\rule[-\dp\strutbox]{.4pt}{\baselineskip}\hspace{2mm}\makebox[12.0mm][r]{\smash{348}}\makebox[9.0mm][r]{\smash{44}}\hspace{3mm}\rule[-\dp\strutbox]{.4pt}{\baselineskip}\hspace{2mm}\makebox[60.0mm][l]{\smash{\textarabic{نز} et \textarabic{نج} pro \textarabic{مد}; \textarabic{سمه} pro \textarabic{سمح}}}}
\DLine{AB01-PDF0656-body-L024}{}{\hspace*{34.0mm}\makebox[5.0mm][l]{\smash{}}\makebox[12.0mm][r]{\smash{360}}\makebox[9.0mm][r]{\smash{0}}\hspace{3mm}\rule[-\dp\strutbox]{.4pt}{\baselineskip}\hspace{2mm}\makebox[12.0mm][r]{\smash{360}}\makebox[9.0mm][r]{\smash{0}}\hspace{3mm}\rule[-\dp\strutbox]{.4pt}{\baselineskip}\hspace{2mm}\makebox[12.0mm][r]{\smash{360}}\makebox[9.0mm][r]{\smash{0}}\hspace{3mm}\rule[-\dp\strutbox]{.4pt}{\baselineskip}\hspace{2mm}\makebox[60.0mm][l]{\smash{}}}
\par\vspace{2mm}
\hypertarget{AB01-PDF0656-P02}{}
\DLine{AB01-PDF0656-body-L025}{25}{\Indent Si iustae essent emendationes (1), distantiam Lunae a nodo necessariam ut contactus di-}
\DLine{AB01-PDF0656-body-L026}{}{scorum eveniat \Name{al-Battānī} statuisset 20° 16′ in boream, 11° 16′ in austrum. \Name{Ptolemaeus}}
\DLine{AB01-PDF0656-body-L027}{}{20° 41′ et 11° 22′. Sed non impossibile est pro 191° emendationis meae esse 190° codicis ser-}
\DLine{AB01-PDF0656-body-L028}{}{vandum, tunc vero 349° \textarabic{سمط} loco codicis \textarabic{سمه} 345° (qui error in scriptura maghrebina facil-}
\DLine{AB01-PDF0656-body-L029}{}{limus) et \textarabic{سمح} 348°; distantiae a nodo ita fiunt 20° 16′, 10° 26′. — Sed nunc puto veros nu-}
\DLine{AB01-PDF0656-body-L030}{30}{meros \textit{alios esse quam quos in tabulis imprimendis receperam.}}
\hypertarget{AB01-PDF0656-P03}{}
\DLine{AB01-PDF0656-body-L031}{}{\Indent Ut melius dubitationis causa in emendando appareat, rationem qua \Name{Ptolemaeus} ter-}
\DLine{AB01-PDF0656-body-L032}{}{minos \textit{eclipsium solarium} constituit exponam. Sit \textgreek{Σ} summa semidiametrorum Solis et Lunae}
\DLine{AB01-PDF0656-body-L033}{}{in eclipsibus, sive distantia centrorum necessaria ut contactus fiat. Parallaxis Lunae negligi}
\DLine{AB01-PDF0656-body-L034}{}{non potest; maximam parallaxim latitudinis, quae omnibus septem Terrae climatibus inter Me-}
\DLine{AB01-PDF0656-body-L035}{35}{roëm (cuius maximus dies 13\textsuperscript{\textit{h}}) et Boristhenis ostia (ubi maximus dies 16\textsuperscript{\textit{h}}) contentis contingere pos-}
\DLine{AB01-PDF0656-body-L036}{}{sit, Ptolemaeus considerat esse 58′ in austrum, 8′ in boream; tunc maximas parallaxes lunares}
\DLine{AB01-PDF0656-body-L037}{}{longitudinis esse 15′ (in Scorpio Piscibusque) et 30′ (in Leone ac Geminis). Ergo latitudo apparens}
\DLine{AB01-PDF0656-body-L038}{}{Lunae necessaria ad contactum erit summum \textgreek{Σ} + 58′ pro Luna boreali, \textgreek{Σ} + 8′ pro australi.}
\DLine{AB01-PDF0656-body-L039}{}{Ab his latitudinibus, per tabulas aut per regulam}
\par\vspace{1mm}
\hypertarget{AB01-PDF0656-H40}{}
\DLine{AB01-PDF0656-body-L040}{40}{\CenterLine{sin dist. {\MoonFont ☾} a nodo boreali $= \mathit{R}\;\dfrac{\text{sin datae lat. {\MoonFont ☾}}}{\text{sin maximae lat. {\MoonFont ☾}}}$\,,}}
\par\vspace{6pt}\hrule height.25pt\vspace{5pt}
\noindent\begin{minipage}[t]{.48\textwidth}\vspace{0pt}\fontsize{9}{11.5}\selectfont\setlength{\GutterOffset}{0pt}
\hypertarget{AB01-PDF0656-N01}{}
\DLine{AB01-PDF0656-notes_left-L001}{}{\Indent (1) Numeros addo Hispanicae versionis f. 41,v.,}
\DLine{AB01-PDF0656-notes_left-L002}{}{inter parentheses notans discrepantias f. 56,v. Ter-}
\DLine{AB01-PDF0656-notes_left-L003}{}{mini eclipsium Solis: 159° 44′ — 197° 16′ (190° 57′),}
\end{minipage}
\hfill\begin{minipage}[t]{.48\textwidth}\vspace{0pt}\fontsize{9}{11.5}\selectfont\setlength{\GutterOffset}{0pt}
\DLine{AB01-PDF0656-notes_right-L001}{}{0° — 20° 16′, 348° 21′ (348° 57′) — 360°. Termini Lu-}
\DLine{AB01-PDF0656-notes_right-L002}{}{nae: 165° 53′ (165° 13′) — 194° 47′, 0° — 14° 47′, 345°}
\DLine{AB01-PDF0656-notes_right-L003}{}{13′ — 360°.}
\end{minipage}
\DSource{0657}{208}
\DLine{AB01-PDF0657-body-L001}{}{deducuntur distantiae apparentes \textgreek{Δ} et \textgreek{Δ}′ Lunae a nodo, seu loci coniunctionis \textit{apparentis.}}
\DLine{AB01-PDF0657-body-L002}{}{Vera coniunctio sit etiam adhuc futura; parallaxes longitudinis (15′ pro Luna australi, 30′}
\DLine{AB01-PDF0657-body-L003}{}{pro boreali) addantur distantiis \textgreek{Δ} et \textgreek{Δ}′; provenient distantiae D et D′ seu loci veri \textit{verae} co-}
\DLine{AB01-PDF0657-body-L004}{}{niunctionis. Coniunctio media adhuc sit futura; quantitas}
\par\vspace{1mm}
\hypertarget{AB01-PDF0657-H05}{}
\DLine{AB01-PDF0657-body-L005}{5}{\CenterLine{(\textgreek{α})\hspace{14mm}maxima aequatio {\MoonFont ⊙} $+\;\dfrac{\text{maxima aeq. {\MoonFont ⊙} $+$ maxima aeq. simpl. {\MoonFont ☾}}}{12}$\,,}}
\par\vspace{1mm}
\DLine{AB01-PDF0657-body-L006}{}{quam Ptolemaeus statuit esse maximam differentiam longitudinis quae inter locum verum verae}
\DLine{AB01-PDF0657-body-L007}{}{coniunctionis et medium mediae cadere possit, addatur distantiis D et D′; exibunt distantiae}
\DLine{AB01-PDF0657-body-L008}{}{maximae Lunae a nodo in \textit{mediis} coniunctionibus, seu \textit{limites eclipsium solarium.}}
\hypertarget{AB01-PDF0657-P01}{}
\DLine{AB01-PDF0657-body-L009}{}{\Indent Ita \Name{Ptolemaeus}, quem tamen \Name{Regiomontanus} (1) monet erravisse in D ac D′ sup-}
\DLine{AB01-PDF0657-body-L010}{10}{putandis; non enim parallaxes longitudinis erant addendae, sed ipsarum duodecima pars, quia}
\DLine{AB01-PDF0657-body-L011}{}{tempore quo Luna spatium parallaxis longitudinis percurrit, Sol eius duodecimam conficit. —}
\DLine{AB01-PDF0657-body-L012}{}{\Name{Al-Battānī}, cap. XXX (t. I, p. 58), in eclipsibus statuit maximam semidiametrum angularem}
\DLine{AB01-PDF0657-body-L013}{}{Solis 16′ 50″, maximam Lunae 17′ 40″; ergo summa semidiametrorum est 34′ 30″. Si \textit{paral-}}
\DLine{AB01-PDF0657-body-L014}{}{\textit{laxibus Ptolemaicis} usus esset, invenisset}
\par\vspace{1mm}
\DLine{AB01-PDF0657-body-L015}{15}{\hspace*{36.0mm}\makebox[36.0mm][l]{\smash{latitudinem {\MoonFont ☾} austr.}}\makebox[15.0mm][r]{\smash{34′ 30″}}\makebox[7.0mm][c]{\smash{+}}\makebox[14.0mm][r]{\smash{58′ 0″}}\makebox[7.0mm][c]{\smash{=}}\makebox[22.0mm][r]{\smash{1° 32′ 30″}}}
\DLine{AB01-PDF0657-body-L016}{}{\hspace*{36.0mm}\makebox[36.0mm][l]{\smash{latitudinem {\MoonFont ☾} boreal.}}\makebox[15.0mm][r]{\smash{34′ 30″}}\makebox[7.0mm][c]{\smash{+}}\makebox[14.0mm][r]{\smash{8′ 0″}}\makebox[7.0mm][c]{\smash{=}}\makebox[22.0mm][r]{0° 42′ 30″\rlap{,}}}
\par\vspace{1mm}
\DLine{AB01-PDF0657-body-L017}{}{quibus latitudinibus respondent distantiae a nodo 17° 58′ 48″.6 et 8° 9′ 16″ (2). \Name{Albatenio}}
\DLine{AB01-PDF0657-body-L018}{}{maxima aequatio Solis est 1° 59′ 10″, maxima simplex Lunae 5° 1′; ergo (\textgreek{α}) est 2° 34′ 10″.}
\DLine{AB01-PDF0657-body-L019}{}{Unde distantiae provenient:}
\par\vspace{1mm}
\DLine{AB01-PDF0657-body-L020}{20}{\hspace*{34.0mm}\makebox[52.0mm][c]{\smash{iuxta rationem Ptolemaei}}\makebox[10.0mm][l]{\smash{}}\makebox[52.0mm][c]{\smash{iuxta eandem emendatam}}}
\DLine{AB01-PDF0657-body-L021}{}{\hspace*{36.0mm}\makebox[10.0mm][r]{17\rlap{°}}\makebox[9.0mm][r]{59\rlap{′}}\makebox[8.0mm][l]{\smash{}}\makebox[10.0mm][r]{8\rlap{°}}\makebox[9.0mm][r]{9\rlap{′}}\makebox[16.0mm][l]{\smash{}}\makebox[10.0mm][r]{17\rlap{°}}\makebox[9.0mm][r]{59\rlap{′}}\makebox[8.0mm][l]{\smash{}}\makebox[10.0mm][r]{8\rlap{°}}\makebox[9.0mm][r]{9\rlap{′}}}
\DLine{AB01-PDF0657-body-L022}{}{\hspace*{36.0mm}\makebox[10.0mm][r]{\smash{}}\makebox[9.0mm][r]{\smash{15}}\makebox[8.0mm][l]{\smash{}}\makebox[10.0mm][r]{\smash{}}\makebox[9.0mm][r]{\smash{30}}\makebox[16.0mm][l]{\smash{}}\makebox[10.0mm][r]{\smash{}}\makebox[9.0mm][r]{\smash{1}}\makebox[8.0mm][l]{\smash{}}\makebox[10.0mm][r]{\smash{}}\makebox[9.0mm][r]{\smash{3}}}
\DLine{AB01-PDF0657-body-L023}{}{\hspace*{36.0mm}\makebox[10.0mm][r]{\smash{2}}\makebox[9.0mm][r]{\smash{34}}\makebox[8.0mm][l]{\smash{}}\makebox[10.0mm][r]{\smash{2}}\makebox[9.0mm][r]{\smash{34}}\makebox[16.0mm][l]{\smash{}}\makebox[10.0mm][r]{\smash{2}}\makebox[9.0mm][r]{\smash{34}}\makebox[8.0mm][l]{\smash{}}\makebox[10.0mm][r]{\smash{2}}\makebox[9.0mm][r]{\smash{34}}}
\par\nointerlineskip\vspace{1pt}\noindent\hspace*{36.0mm}\rule{10.0mm}{.4pt}\rule{9.0mm}{.4pt}\hspace{8.0mm}\rule{10.0mm}{.4pt}\rule{9.0mm}{.4pt}\hspace{16.0mm}\rule{10.0mm}{.4pt}\rule{9.0mm}{.4pt}\hspace{8.0mm}\rule{10.0mm}{.4pt}\rule{9.0mm}{.4pt}\par\nointerlineskip\vspace{2pt}
\DLine{AB01-PDF0657-body-L024}{}{\hspace*{36.0mm}\makebox[10.0mm][r]{20\rlap{°}}\makebox[9.0mm][r]{48\rlap{′}}\makebox[8.0mm][l]{\smash{}}\makebox[10.0mm][r]{11\rlap{°}}\makebox[9.0mm][r]{13\rlap{′}}\makebox[16.0mm][l]{\smash{}}\makebox[10.0mm][r]{20\rlap{°}}\makebox[9.0mm][r]{34\rlap{′}}\makebox[8.0mm][l]{\smash{}}\makebox[10.0mm][r]{10\rlap{°}}\makebox[9.0mm][r]{46\rlap{′}}}
\par\vspace{1mm}
\hypertarget{AB01-PDF0657-P02}{}
\DLine{AB01-PDF0657-body-L025}{25}{\Indent Sed distantiae utraque ratione supputatae, cum numeris codicis non congruunt. Calculos}
\DLine{AB01-PDF0657-body-L026}{}{iterum suscipiam, ponens \Name{al-Battānī} parvam differentiam in semidiametro Solis repertam}
\DLine{AB01-PDF0657-body-L027}{}{neglexisse et quantitate \textit{Ptolemaica} usum esse (15′ 40″); summa semidiametrorum tunc}
\DLine{AB01-PDF0657-body-L028}{}{15′ 40″ + 17′ 40″ = 33′ 20″. Quibus addantur parallaxes latitudinis 58′ et 8′; provenient}
\DLine{AB01-PDF0657-body-L029}{}{latitudines apparentes 1° 31′ 20″ in boream, 0° 41′ 20″ in austrum, sicut in \textit{Almagesto.} Ab}
\DLine{AB01-PDF0657-body-L030}{30}{iis per tabulas nostras logarithmicas deducuntur distantiae apparentes a nodo 17° 44′ 45″,}
\DLine{AB01-PDF0657-body-L031}{}{7° 55′ 45″; per tabulas \Name{Albatenii} 17° 45′ 42″, 7° 55′ 45″. Sed ponamus \Name{al-Battānī} quan-}
\DLine{AB01-PDF0657-body-L032}{}{titates recepisse quas \Name{Ptolemaeus} rudiore supputatione invenerat, 17° 27′ et 7° 52′. His lati-}
\DLine{AB01-PDF0657-body-L033}{}{tudinibus addantur parallaxes longitudinis 15′ et 30′; utrique summae adiiciatur quantitas (\textgreek{α})}
\DLine{AB01-PDF0657-body-L034}{}{iuxta aequationes \textit{Albatenianas} luminarium seu 2° 34′ (iuxta aequationes Ptolemaei 3°); pro-}
\DLine{AB01-PDF0657-body-L035}{35}{\hspace*{0.0mm}\makebox[54.0mm][l]{\smash{dibunt}}\makebox[16.0mm][r]{\smash{17° 27′}}\makebox[7.0mm][c]{\smash{+}}\makebox[13.0mm][r]{\smash{0° 15′}}\makebox[7.0mm][c]{\smash{+}}\makebox[13.0mm][r]{\smash{2° 34′}}\makebox[7.0mm][c]{\smash{=}}\makebox[18.0mm][r]{\smash{20° 16′}}}
\DLine{AB01-PDF0657-body-L036}{}{\hspace*{0.0mm}\makebox[54.0mm][l]{\smash{}}\makebox[16.0mm][r]{\smash{7° 52′}}\makebox[7.0mm][c]{\smash{+}}\makebox[13.0mm][r]{\smash{0° 30′}}\makebox[7.0mm][c]{\smash{+}}\makebox[13.0mm][r]{\smash{2° 34′}}\makebox[7.0mm][c]{\smash{=}}\makebox[18.0mm][r]{10° 56′\rlap{.}}}
\DLine{AB01-PDF0657-body-L037}{}{Prior numerus ad amussim cum terminis borealibus codicis congruit; alter cum australibus}
\DLine{AB01-PDF0657-body-L038}{}{convenit si \textarabic{نو} 56′ pro \textarabic{يو} 16′ cod., \textarabic{د} 4′ pro \textarabic{نج} 53′ et \textarabic{نز} 57′ cod., \textarabic{سمط} 349° (ut antea iam pro-}
\par\vspace{6pt}\hrule height.25pt\vspace{5pt}
\noindent\begin{minipage}[t]{.48\textwidth}\vspace{0pt}\fontsize{9}{11.5}\selectfont\setlength{\GutterOffset}{0pt}
\hypertarget{AB01-PDF0657-N01}{}
\DLine{AB01-PDF0657-notes_left-L001}{}{\Indent (1) I. \Name{de Monteregio} \textit{et} G. \Name{Purbachii}}
\DLine{AB01-PDF0657-notes_left-L002}{}{\textit{Epitome in Cl. Ptolemaei magnam compositionem,}}
\DLine{AB01-PDF0657-notes_left-L003}{}{Basileae 1543, lib. VI, prop. VII, p. 118.}
\hypertarget{AB01-PDF0657-N02}{}
\DLine{AB01-PDF0657-notes_left-L004}{}{\Indent (2) Ita iuxta tabulas logarithmicas et formu-}
\end{minipage}
\hfill\begin{minipage}[t]{.48\textwidth}\vspace{0pt}\fontsize{9}{11.5}\selectfont\setlength{\GutterOffset}{0pt}
\DLine{AB01-PDF0657-notes_right-L001}{}{lam Albatenianam praecedenti pagina memoratam.}
\DLine{AB01-PDF0657-notes_right-L002}{}{E tabulis \Name{Albatenii} t. II, p. 78, in quibus inter-}
\DLine{AB01-PDF0657-notes_right-L003}{}{polatio rudior fit ob intervallum 1° in argumento}
\DLine{AB01-PDF0657-notes_right-L004}{}{tabularum, 17° 59′ 48″ et 8° 9′ 18″.}
\end{minipage}
\DSource{0658}{209}
\DLine{AB01-PDF0658-body-L001}{}{posueram) pro \textarabic{سمه} 345° et \textarabic{سمح} 348° cod. restituantur. — Rogo igitur ut hic et pag. 88, \textit{termini}}
\DLine{AB01-PDF0658-body-L002}{}{\textit{eclipsium Solis ita emendentur:}}
\par\vspace{1mm}
\DLine{AB01-PDF0658-body-L003}{}{\hspace*{50.0mm}\makebox[5.0mm][l]{\smash{\raisebox{-6.9pt}{$\left\{\rule[-10pt]{0pt}{20pt}\right.$}}}\makebox[11.0mm][r]{159\rlap{°}}\makebox[9.0mm][r]{44\rlap{′}}\makebox[22.0mm][l]{\smash{}}\makebox[5.0mm][l]{\smash{\raisebox{-6.9pt}{$\left\{\rule[-10pt]{0pt}{20pt}\right.$}}}\makebox[11.0mm][r]{0\rlap{°}}\makebox[9.0mm][r]{0\rlap{′}}\makebox[22.0mm][l]{\smash{}}\makebox[5.0mm][l]{\smash{\raisebox{-6.9pt}{$\left\{\rule[-10pt]{0pt}{20pt}\right.$}}}\makebox[11.0mm][r]{\textit{349}\rlap{\textit{°}}}\makebox[9.0mm][r]{\textit{4}\rlap{\textit{′}}}}
\DLine{AB01-PDF0658-body-L004}{}{\hspace*{50.0mm}\makebox[5.0mm][l]{\smash{}}\makebox[11.0mm][r]{\smash{\textit{190}}}\makebox[9.0mm][r]{\smash{\textit{56}}}\makebox[22.0mm][l]{\smash{}}\makebox[5.0mm][l]{\smash{}}\makebox[11.0mm][r]{\smash{20}}\makebox[9.0mm][r]{\smash{16}}\makebox[22.0mm][l]{\smash{}}\makebox[5.0mm][l]{\smash{}}\makebox[11.0mm][r]{\smash{360}}\makebox[9.0mm][r]{\smash{0}}}
\par\vspace{1mm}
\hypertarget{AB01-PDF0658-P01}{}
\DLine{AB01-PDF0658-body-L005}{5}{\Indent \Name{Regiomontanus}, l. c., affirmat \Name{al-Battānī} invenisse 20° 12′ et 10° 40′; quo modo}
\DLine{AB01-PDF0658-body-L006}{}{hos numeros supputaverit nescio. — \Name{Ḥabash}, f. 122,r. et 133,r. terminos Ptolemaei servavit.}
\hypertarget{AB01-PDF0658-P02}{}
\DLine{AB01-PDF0658-body-L007}{}{\Indent \textit{Termini eclipsium lunarium} certissime restitui potuerunt:}
\par\vspace{1mm}
\DLine{AB01-PDF0658-body-L008}{}{\hspace*{36.0mm}\makebox[6.0mm][l]{\smash{\raisebox{-6.9pt}{$\left\{\rule[-10pt]{0pt}{20pt}\right.$}}}\makebox[12.0mm][r]{165\rlap{°}}\makebox[9.0mm][r]{\smash{...}}\hspace{3mm}\rule[-\dp\strutbox]{.4pt}{\baselineskip}\hspace{2mm}\makebox[12.0mm][r]{\textit{355}\rlap{\textit{°}}}\makebox[9.0mm][r]{\textit{53}\rlap{\textit{′}}}\hspace{3mm}\rule[-\dp\strutbox]{.4pt}{\baselineskip}\hspace{2mm}\makebox[12.0mm][r]{165\rlap{°}}\makebox[9.0mm][r]{13\rlap{′}}\hspace{3mm}\rule[-\dp\strutbox]{.4pt}{\baselineskip}\hspace{2mm}\makebox[60.0mm][l]{\smash{\textarabic{سنه} false pro \textarabic{قصه}, et \textarabic{نج} pro \textarabic{يج}}}}
\DLine{AB01-PDF0658-body-L009}{}{\hspace*{36.0mm}\makebox[6.0mm][l]{\smash{}}\makebox[12.0mm][r]{\smash{194}}\makebox[9.0mm][r]{\smash{47}}\hspace{3mm}\rule[-\dp\strutbox]{.4pt}{\baselineskip}\hspace{2mm}\makebox[12.0mm][r]{\smash{\textit{167}}}\makebox[9.0mm][r]{\smash{\textit{44}}}\hspace{3mm}\rule[-\dp\strutbox]{.4pt}{\baselineskip}\hspace{2mm}\makebox[12.0mm][r]{\smash{194}}\makebox[9.0mm][r]{\smash{47}}\hspace{3mm}\rule[-\dp\strutbox]{.4pt}{\baselineskip}\hspace{2mm}\makebox[60.0mm][l]{\smash{\textarabic{قصز} false pro \textarabic{قضد}, et \textarabic{مد} pro \textarabic{مز}}}}
\DLine{AB01-PDF0658-body-L010}{10}{\hspace*{36.0mm}\makebox[6.0mm][l]{\smash{\raisebox{-6.9pt}{$\left\{\rule[-10pt]{0pt}{20pt}\right.$}}}\makebox[12.0mm][r]{\smash{0}}\makebox[9.0mm][r]{\smash{0}}\hspace{3mm}\rule[-\dp\strutbox]{.4pt}{\baselineskip}\hspace{2mm}\makebox[12.0mm][r]{\smash{0}}\makebox[9.0mm][r]{\smash{0}}\hspace{3mm}\rule[-\dp\strutbox]{.4pt}{\baselineskip}\hspace{2mm}\makebox[12.0mm][r]{\smash{0}}\makebox[9.0mm][r]{\smash{0}}\hspace{3mm}\rule[-\dp\strutbox]{.4pt}{\baselineskip}\hspace{2mm}\makebox[60.0mm][l]{\smash{}}}
\DLine{AB01-PDF0658-body-L011}{}{\hspace*{36.0mm}\makebox[6.0mm][l]{\smash{}}\makebox[12.0mm][r]{\smash{14}}\makebox[9.0mm][r]{\smash{47}}\hspace{3mm}\rule[-\dp\strutbox]{.4pt}{\baselineskip}\hspace{2mm}\makebox[12.0mm][r]{\smash{14}}\makebox[9.0mm][r]{\smash{47}}\hspace{3mm}\rule[-\dp\strutbox]{.4pt}{\baselineskip}\hspace{2mm}\makebox[12.0mm][r]{\smash{14}}\makebox[9.0mm][r]{\smash{47}}\hspace{3mm}\rule[-\dp\strutbox]{.4pt}{\baselineskip}\hspace{2mm}\makebox[60.0mm][l]{\smash{}}}
\DLine{AB01-PDF0658-body-L012}{}{\hspace*{36.0mm}\makebox[6.0mm][l]{\smash{\raisebox{-6.9pt}{$\left\{\rule[-10pt]{0pt}{20pt}\right.$}}}\makebox[12.0mm][r]{\smash{345}}\makebox[9.0mm][r]{\smash{\textit{53}}}\hspace{3mm}\rule[-\dp\strutbox]{.4pt}{\baselineskip}\hspace{2mm}\makebox[12.0mm][r]{\smash{345}}\makebox[9.0mm][r]{\smash{\textit{53}}}\hspace{3mm}\rule[-\dp\strutbox]{.4pt}{\baselineskip}\hspace{2mm}\makebox[12.0mm][r]{\smash{345}}\makebox[9.0mm][r]{\smash{13}}\hspace{3mm}\rule[-\dp\strutbox]{.4pt}{\baselineskip}\hspace{2mm}\makebox[60.0mm][l]{\smash{\textarabic{نج} false pro \textarabic{يج}}}}
\DLine{AB01-PDF0658-body-L013}{}{\hspace*{36.0mm}\makebox[6.0mm][l]{\smash{}}\makebox[12.0mm][r]{\smash{360}}\makebox[9.0mm][r]{\smash{0}}\hspace{3mm}\rule[-\dp\strutbox]{.4pt}{\baselineskip}\hspace{2mm}\makebox[12.0mm][r]{\smash{360}}\makebox[9.0mm][r]{\smash{0}}\hspace{3mm}\rule[-\dp\strutbox]{.4pt}{\baselineskip}\hspace{2mm}\makebox[12.0mm][r]{\smash{360}}\makebox[9.0mm][r]{\smash{0}}\hspace{3mm}\rule[-\dp\strutbox]{.4pt}{\baselineskip}\hspace{2mm}\makebox[60.0mm][l]{\smash{}}}
\par\vspace{1mm}
\hypertarget{AB01-PDF0658-P03}{}
\DLine{AB01-PDF0658-body-L014}{}{\Indent Unde elucet \Name{al-Battānī} distantiam maximam Lunae a nodo in \textit{mediis} oppositionibus ad}
\DLine{AB01-PDF0658-body-L015}{15}{contactum cum umbra posuisse 14° 47′. Supputatio facilis. In eclipsibus reperit (t. I, p. 58)}
\DLine{AB01-PDF0658-body-L016}{}{maximam semidiametrum Lunae 17′ 40″, maximam umbrae 46′ 0″; earum summa 1° 3′ 40″,}
\DLine{AB01-PDF0658-body-L017}{}{cui latitudini respondet distantia a nodo 12° 16′ 23″ iuxta tabulas nostras logarithmicas,}
\DLine{AB01-PDF0658-body-L018}{}{12° 16′ 25″ iuxta tabulas \Name{al-Battānī} t. II, p. 78. Procul dubio \Name{Albatenius} rudiorem valo-}
\DLine{AB01-PDF0658-body-L019}{}{rem 12° 13′ recepit, proximum 12° 12′ \textit{Almagesti;} quibus addita quantitate (\textgreek{α}) seu 2° 34′, pro-}
\DLine{AB01-PDF0658-body-L020}{20}{veniunt 14° 47′. — Miro modo \Name{Regiomontanus}, VI, 8, p. 119, ex sua coniectura: « Al-}
\DLine{AB01-PDF0658-body-L021}{}{« bategni tamen dicit terminum esse 14. grad. 15. minut. ». — \Name{Ḥabash}, f. 122, r., quan-}
\DLine{AB01-PDF0658-body-L022}{}{titatem \Name{Ptolemaei} servat, i. e. 15° 12′ a nodo.}
\par\vspace{6mm}
\hypertarget{AB01-PDF0658-H23}{}
\DLine{AB01-PDF0658-body-L023}{}{\CenterLine{\textit{Ad pag. 33–54.} — \textbf{Tabulae geographicae.}}}
\par\vspace{3mm}
\hypertarget{AB01-PDF0658-P04}{}
\DLine{AB01-PDF0658-body-L024}{}{\Indent Memorantur t. I, p. 20, 30, 73, 137 (l. 2), seu cap. VI, XIV, XXXIV, LVI. — Textus}
\DLine{AB01-PDF0658-body-L025}{25}{Arabicus, in quo codicis numeros immutatos servavi, est t. III, p. 234–41.}
\hypertarget{AB01-PDF0658-P05}{}
\DLine{AB01-PDF0658-body-L026}{}{\Indent Arabice et Gallice eas iam ediderat ac illustraverat J. \Name{Lelewel}, \textit{Géographie du moyen}}
\DLine{AB01-PDF0658-body-L027}{}{\textit{âge,} Breslau 1852, t. IV, Épilogue, p. 61–100 (tabulae, cum notis, p. 64–93), iuxta apographum}
\DLine{AB01-PDF0658-body-L028}{}{codicis Escurialensis charta pellucida factum. Sed linguae fere ignarus multisque Arabicis}
\DLine{AB01-PDF0658-body-L029}{}{subsidiis destitutus, non semper potuit clarus vir bonas emendationes et coniecturas proponere,}
\DLine{AB01-PDF0658-body-L030}{30}{nec de tabularum loco in historia geographiae recte iudicare. Duae etiam mappae quas ad}
\DLine{AB01-PDF0658-body-L031}{}{Albatenii mentem construere studuit plurimis locis emendandae sunt.}
\hypertarget{AB01-PDF0658-P06}{}
\DLine{AB01-PDF0658-body-L032}{}{\Indent Easdem tabulas ego Italice edidi et multis notis praemissaque dissertatione illustravi in}
\DLine{AB01-PDF0658-body-L033}{}{meo \textit{Le tabelle geografiche d’al-Battâni,} Torino 1898, seorsum edito ex \textit{Cosmos di} \Name{Guido}}
\DLine{AB01-PDF0658-body-L034}{}{\Name{Cora}, ser. II, t. XII, p. 161–83. Quae tunc scripseram, aucta et passim etiam emendata, hic}
\DLine{AB01-PDF0658-body-L035}{35}{atque in notis ad ipsas tabulas praebeo.}
\hypertarget{AB01-PDF0658-P07}{}
\DLine{AB01-PDF0658-body-L036}{}{\Indent Longitudines ac latitudines praecipuarum urbium in omnibus astronomicis operibus Arabum}
\DLine{AB01-PDF0658-body-L037}{}{indicantur, et in Tabulis manualibus \Name{Theonis Alexandrini}. Desunt in \textit{Almagesto,} quia}
\DLine{AB01-PDF0658-body-L038}{}{\Name{Ptolemaeo}, ut ipse II, 12 (ed. Halma, t. I, p. 148) dicit, in animo erat Geographiam seor-}
\DLine{AB01-PDF0658-body-L039}{}{sum edere.}
\hypertarget{AB01-PDF0658-P08}{}
\DLine{AB01-PDF0658-body-L040}{40}{\Indent Ex ipsarum tabularum inscriptionibus (t. II, p. 33 ac 37) et ex iis quae t. I, p. 19–20 narrat,}
\DLine{AB01-PDF0658-body-L041}{}{patet \Name{al-Battānī} situs urbium e \textit{Libro Figurae Terrae} (\textit{Kitāb ṣūrat al-arḍ}) sumpsisse,}
\par\vspace{5mm}\hfill{\fontsize{8}{10}\selectfont 27}\hspace{10mm}
\DSource{0659}{210}
\DLine{AB01-PDF0659-body-L001}{}{quo longitudines et latitudines aut observationibus astronomicis aut itinerariis « modo veritati}
\DLine{AB01-PDF0659-body-L002}{}{proximo » (p. 20, l. 4 et 9) deductae erant. Addit etiam p. 20: Has tabulas geographicas con-}
\DLine{AB01-PDF0659-body-L003}{}{didimus « ad fidem descriptionis (\textit{rasm}) quam in noto libro \textit{Figurae Terrae} (1) invenimus;}
\DLine{AB01-PDF0659-body-L004}{}{« item media loca provinciarum et regionum, quae 94 sunt, separatim notavimus, ut \Name{Ptole-}}
\DLine{AB01-PDF0659-body-L005}{5}{« \Name{maeus} (2) fecit. Sed hoc libro [Figurae Terrae] errores longitudinum et latitudinum repe-}
\DLine{AB01-PDF0659-body-L006}{}{« riuntur ». In fine etiam XIV capitis (t. I, p. 30) monet: « Latitudo loci quae e tabula latitu-}
\DLine{AB01-PDF0659-body-L007}{}{« dinum urbium deprehenditur, veritati tantum est propinqua, nec ita certa ut latitudo obser-}
\DLine{AB01-PDF0659-body-L008}{}{« vationibus cognita ». Nec multum fidei longitudinibus illis \Name{al-Battānī} tribuisse, ex eo quoque}
\DLine{AB01-PDF0659-body-L009}{}{elucet quod t. I intervalla inter meridianum ar-Raqqah et meridianos Alexandriae atque An-}
\DLine{AB01-PDF0659-body-L010}{10}{tiochiae aliquanto diversa sumpsit quam quae e tabulis eruuntur. Posuit enim}
\par\vspace{1mm}
\DLine{AB01-PDF0659-body-L011}{}{\hspace*{16.0mm}\makebox[42.0mm][l]{\smash{t. I, p. 56, l. 26:}}\makebox[62.0mm][l]{\smash{inter ar-Raqqah et Antiochiam}}\makebox[52.0mm][r]{\smash{\textit{amplius} 10\textsuperscript{\textit{m}} = 2° 30′}}}
\DLine{AB01-PDF0659-body-L012}{}{\hspace*{16.0mm}\makebox[42.0mm][l]{\smash{t. I, p. 57, l. 22–23:}}\makebox[62.0mm][l]{\smash{}}\makebox[52.0mm][r]{\smash{\textit{fere} $\qfrac{1}{4}$ horae = 3° 45′}}}
\DLine{AB01-PDF0659-body-L013}{}{\hspace*{16.0mm}\makebox[42.0mm][l]{\smash{in tabulis nr. 150 et 142:}}\makebox[62.0mm][l]{\smash{}}\makebox[52.0mm][r]{\smash{73° 15′ — 69° 0′ = 4° 15′}}}
\par\vspace{1mm}
\DLine{AB01-PDF0659-body-L014}{}{servata scilicet Antiochiae longitudine Ptolemaica. In nostris mappis intervallum 2° 55′ circiter.}
\DLine{AB01-PDF0659-body-L015}{15}{— Item}
\par\vspace{1mm}
\DLine{AB01-PDF0659-body-L016}{}{\hspace*{17.0mm}\makebox[41.0mm][l]{\smash{t. I, p. 42:}}\makebox[114.0mm][r]{\smash{inter ar-Raqqah et Alexandriam \textit{circiter} $\qfrac{2}{3}$ horae = 10° 0′}}}
\DLine{AB01-PDF0659-body-L017}{}{\hspace*{17.0mm}\makebox[41.0mm][l]{\smash{in tabulis nr. 150 et 109}}\makebox[114.0mm][r]{\smash{73° 15′ — 60° 30′ = 12° 45′}}}
\par\vspace{1mm}
\DLine{AB01-PDF0659-body-L018}{}{servata longitudine Ptolemaica Alexandriae. In mappis nostris intervallum 9° 12′.}
\hypertarget{AB01-PDF0659-P01}{}
\DLine{AB01-PDF0659-body-L019}{}{\Indent Latitudines autem urbium Mekkah (nr. 103), Ḥarrān (nr. 152) et Baghdād (nr. 168) etiam}
\DLine{AB01-PDF0659-body-L020}{20}{t. II, p. 67 adhibuit quales his tabulis geographicis indicantur. Situm ar-Raqqah suis astrono-}
\DLine{AB01-PDF0659-body-L021}{}{micis observationibus determinavit.}
\hypertarget{AB01-PDF0659-P02}{}
\DLine{AB01-PDF0659-body-L022}{}{\Indent Nunc quid ille « \textit{Liber Figurae Terrae,} qui vocatur \textit{Geographia} » fuerit investigandum}
\DLine{AB01-PDF0659-body-L023}{}{est. Collatis sitibus tabularum Albatenii cum iis quos Abū ’l-fidā’ (3) iuxta \Name{al-Khuwārizmī}}
\DLine{AB01-PDF0659-body-L024}{}{passim praebet, \Name{Lelewel} putavit illum esse librum \textit{Figurae Terrae} a \Name{Muḥammad ibn}}
\DLine{AB01-PDF0659-body-L025}{25}{\Name{Mūsà al-Khuwārizmī} iussu khalifae al-Ma’mūn conscriptum (4). Sed melius re perpensa}
\DLine{AB01-PDF0659-body-L026}{}{et integro al-Khuwārizmī libro perlecto qui in codice Bibliothecae Argentoratensis nunc exstat,}
\DLine{AB01-PDF0659-body-L027}{}{sententiam illam esse abiiciendam statim elucet. Nomina Graeca saepe aliter ab al-Khuwārizmī,}
\DLine{AB01-PDF0659-body-L028}{}{aliter ab al-Battānī scribuntur; cfr. prioris \textit{Ghermāniyā, Muwesiyah, Ǵanǵis, Nighīrā, Qada-}}
\DLine{AB01-PDF0659-body-L029}{}{\textit{bathmūs} cum alterius \textit{Ǵehrmāniyā, Muwehsiyā, Ghanǵis, Niǵīrā; Qāṭābathmūs.} Littera \textgreek{γ}}
\DLine{AB01-PDF0659-body-L030}{30}{Graecorum ab al-Khuwārizmī semper \textarabic{غ} \textit{gh} redditur (5), ab al-Battānī modo \textarabic{غ} \textit{gh,} modo \textarabic{ج} \textit{ǵ:}}
\DLine{AB01-PDF0659-body-L031}{}{\textit{Ghāliyā, Belghīqī, Yāzūghūs, Lūghdunīsiyā, Eghifṭos, Ghālāṭiyā, Faflāghūniyā, Ghanǵis}}
\DLine{AB01-PDF0659-body-L032}{}{(nr. 89), \textit{Ǵehrmāniyā, Ṭiǵiṭāniyā, Frūǵiyā, Niǵīrā, Frūǵis} (nr. 116). Loca media 56 pro-}
\DLine{AB01-PDF0659-body-L033}{}{vinciarum al-Khuwārizmī determinat in longitudine et latitudine; al-Battānī 93 provinciarum}
\DLine{AB01-PDF0659-body-L034}{}{(quae sunt \textgreek{ἐπαρχίαι} seu \textgreek{χῶραι} Ptolemaei), etiam in coordinatis geographicis ab al-Khuwārizmī}
\DLine{AB01-PDF0659-body-L035}{35}{discrepans. Praeterea, neglectis octo nominibus (nr. 102, 112, 203, 216, 230, 231, 233, 249)}
\par\vspace{6pt}\hrule height.25pt\vspace{5pt}
\noindent\begin{minipage}[t]{.48\textwidth}\vspace{0pt}\fontsize{9}{11.5}\selectfont\setlength{\GutterOffset}{0pt}
\hypertarget{AB01-PDF0659-N01}{}
\DLine{AB01-PDF0659-notes_left-L001}{}{\Indent (1) Plato addit: « qui vocatur \textit{Geographia} ».}
\hypertarget{AB01-PDF0659-N02}{}
\DLine{AB01-PDF0659-notes_left-L002}{}{\Indent (2) Cfr. t. I, p. 20, adn. 2.}
\hypertarget{AB01-PDF0659-N03}{}
\DLine{AB01-PDF0659-notes_left-L003}{}{\Indent (3) \textit{La Géographie d’}\Name{Aboulféda}, \textit{texte arabe}}
\DLine{AB01-PDF0659-notes_left-L004}{}{\textit{publié par Reinaud et Mac Guckin de Slane,} Pa-}
\DLine{AB01-PDF0659-notes_left-L005}{}{ris 1840 (Gallice versa a Reinaud et Guyard, Pari-}
\DLine{AB01-PDF0659-notes_left-L006}{}{siis 1848–83, 3 tom.); ex opere al-Khuwārizmī, quod}
\DLine{AB01-PDF0659-notes_left-L007}{}{plerumque \textit{rasm} vocat, longitudines et latitudines 92}
\DLine{AB01-PDF0659-notes_left-L008}{}{urbium excerpsit.}
\hypertarget{AB01-PDF0659-N04}{}
\DLine{AB01-PDF0659-notes_left-L009}{}{\Indent (4) Innitebantur \textit{mappis} Ptolemaicis \textit{Graecis}}
\end{minipage}
\hfill\begin{minipage}[t]{.48\textwidth}\vspace{0pt}\fontsize{9}{11.5}\selectfont\setlength{\GutterOffset}{0pt}
\DLine{AB01-PDF0659-notes_right-L001}{}{(non versione Syriaca Geographiae ut saepe dicitur}
\DLine{AB01-PDF0659-notes_right-L002}{}{a nostratibus), quibus multas easque saepe optimas}
\DLine{AB01-PDF0659-notes_right-L003}{}{emendationes intulerat. Vide quae disserui \textit{al-Ḫu-}}
\DLine{AB01-PDF0659-notes_right-L004}{}{\textit{wārizmī e il suo rifacimento della Geografia di}}
\DLine{AB01-PDF0659-notes_right-L005}{}{\textit{Tolomeo,} in Memorie della R. Accademia dei Lin-}
\DLine{AB01-PDF0659-notes_right-L006}{}{cei, Classe di scienze morali, ser. V, t. II, 1894-95,}
\DLine{AB01-PDF0659-notes_right-L007}{}{p. 3-53. — Cfr. t. I, p. 312, adn. 4.}
\hypertarget{AB01-PDF0659-N05}{}
\DLine{AB01-PDF0659-notes_right-L008}{}{\Indent (5) \textit{Ǵanǵis} quia haec est forma Arabica no-}
\DLine{AB01-PDF0659-notes_right-L009}{}{minis fluvii.}
\end{minipage}
\DSource{0660}{211}
\DLine{AB01-PDF0660-body-L001}{}{quorum lectio nimis incerta, 58 urbes occurrunt in tabulis Albatenianis quae apud al-Khu-}
\DLine{AB01-PDF0660-body-L002}{}{wārizmī desiderantur. Ex situbus 172 urbium quarum nomina certa sunt in Albatenii codice, 68}
\DLine{AB01-PDF0660-body-L003}{}{e \textit{Geographia} \Name{Ptolemaei} (1) proveniunt, 57 ex \Name{al-Khuwārizmī} (in longitudinibus additis}
\DLine{AB01-PDF0660-body-L004}{}{10 gradibus (2), demum 47 ex \textit{alio fonte Arabico} sed mappis Ptolemaicis plus minusve feli-}
\DLine{AB01-PDF0660-body-L005}{5}{citer accommodati (3). Ex urbibus ab al-Khuwārizmī non memoratis haec sunt quarum situs}
\DLine{AB01-PDF0660-body-L006}{}{cum Ptolemaicis (4) congruunt: \textit{94, 97, 98, 99, 100, 101, 105, 106, 108, 110,} 111, 114,}
\DLine{AB01-PDF0660-body-L007}{}{134, 139, 140, 171, 173, 179, 191, 200, 204, \textit{205, 212,} 213, \textit{215, 217, 218,} 219, 220, 223,}
\DLine{AB01-PDF0660-body-L008}{}{224, 225, \textit{232, 272;} haec autem quarum longitudines et latitudines ab alio fonte (5) manant:}
\DLine{AB01-PDF0660-body-L009}{}{137, *141, *155, *157, *159, *160, 163, *165, *166, *180, 201, 202, 207, 208, *209, 211,}
\DLine{AB01-PDF0660-body-L010}{10}{221, \textit{222,} *226, *228, 229, *249, 265, *266.}
\hypertarget{AB01-PDF0660-P01}{}
\DLine{AB01-PDF0660-body-L011}{}{\Indent Igitur auctor \textit{Libri Figurae Terrae} prae manibus Geographiam \Name{Ptolemaei} et opus}
\DLine{AB01-PDF0660-body-L012}{}{\Name{al-Khuwārizmī} habuit, quibus non pauca addidit. Ptolemaeum \textit{Syriace} versum (6) aut}
\DLine{AB01-PDF0660-body-L013}{}{\textit{Arabice} e \textit{Syriaco,} cognovit; Syriacam enim originem redolent \textit{eh} pro \textgreek{ε} in \textit{Brehṭanīqī,}}
\DLine{AB01-PDF0660-body-L014}{}{\textit{Ǵehrmāniyā, Mehṭānīsā, Muwehsiyā, Qehsarensiyā, Lehfṭis,} biblicum \textit{Kūsh} pro \textgreek{Αἰθιοπία}}
\DLine{AB01-PDF0660-body-L015}{15}{(nr. 45 et 46) et fortasse \textit{Asdūd} (nr. 213) pro Azdūd. Duo vici \textit{Ḥuwwārīn} (nr. 208) et \textit{Ūrim}}
\DLine{AB01-PDF0660-body-L016}{}{(nr. 204) memorantur, ignoti fere Arabicis geographis at, quia episcoporum sedes, a scripto-}
\DLine{AB01-PDF0660-body-L017}{}{ribus \textit{Syro-Christianis} non raro celebrati. Nec praetereundum est e \textit{Syriaco} archetypo etiam}
\DLine{AB01-PDF0660-body-L018}{}{fluxisse geographica VI capitis (vide adn. t. I, p. 166–77) et chronologica t. II, p. 1–3; for-}
\DLine{AB01-PDF0660-body-L019}{}{tasse etiam ventorum rosam t. II, p. 92. Forte « Liber Figurae Terrae qui vocatur Geogra-}
\DLine{AB01-PDF0660-body-L020}{20}{phia » est recensio Arabica operis Ptolemaei, quam Ḥarrānensis celeberrimus \Name{Thābit ibn}}
\DLine{AB01-PDF0660-body-L021}{}{\Name{Qurrah} (211–288 heg., 826-901 Chr.), Graecae linguae ac Syriacae peritissimus, composuerat.}
\DLine{AB01-PDF0660-body-L022}{}{Verisimile enim est, IX saeculo Chr., Thābit Geographiam Ptolemaei non sine mutationibus}
\DLine{AB01-PDF0660-body-L023}{}{transtulisse ut novis temporibus melius responderet; inde \Name{al-Battānī} eas urbes delegisse,}
\DLine{AB01-PDF0660-body-L024}{}{\textit{Syriae} praesertim et \textit{Mesopotamiae,} quae sibi praecipuae viderentur.}
\hypertarget{AB01-PDF0660-P02}{}
\DLine{AB01-PDF0660-body-L025}{25}{\Indent Loco VI capitis antea allato ac in tabularum titulo (t. II, p. 33) provinciae, quarum puncta}
\DLine{AB01-PDF0660-body-L026}{}{media indicantur, 94 esse dicuntur; totidem enumerantur VIII libro Geographiae \Name{Ptolemaei}.}
\DLine{AB01-PDF0660-body-L027}{}{At nonnisi 93 postea in tabulis invenimus; quod ex librariorum negligentia aut etiam ex aliqua}
\DLine{AB01-PDF0660-body-L028}{}{auctoris confusione ortum esse potest, cuius causa ipse Ptolemaeus. Hic enim libro VIII, § 3}
\DLine{AB01-PDF0660-body-L029}{}{conspectus generalis, affirmat Europam 34 aut 35 regiones, Libyam 12, Asiam 49 complecti,}
\DLine{AB01-PDF0660-body-L030}{30}{quarum summa est 95 aut 96; sed § 4 eas esse 87 dicit! Eodem VIII libro, cap. 29, Europae}
\DLine{AB01-PDF0660-body-L031}{}{34 provincias tribuit, 12 Libyae, 49 Asiae (ergo summa 95); sed postea nomina 46 dumtaxat}
\DLine{AB01-PDF0660-body-L032}{}{memorat in Asia, qua re summa fit 92. Tamen lib. VIII, cap. 29, § 29, Ptolemaeus ipse sum-}
\DLine{AB01-PDF0660-body-L033}{}{mam conficit invenitque 94! Nec aliae parvae discrepantiae desunt si haec omnia cum prio-}
\DLine{AB01-PDF0660-body-L034}{}{ribus septem libris comparentur.}
\hypertarget{AB01-PDF0660-P03}{}
\DLine{AB01-PDF0660-body-L035}{35}{\Indent In tabulis quas al-Battānī e \textit{Libro Figurae Terrae} excerpsit, provinciae \textit{Europae} ad}
\par\vspace{6pt}\hrule height.25pt\vspace{5pt}
\noindent\begin{minipage}[t]{.48\textwidth}\vspace{0pt}\fontsize{9}{11.5}\selectfont\setlength{\GutterOffset}{0pt}
\hypertarget{AB01-PDF0660-N01}{}
\DLine{AB01-PDF0660-notes_left-L001}{}{\Indent (1) Nr. 94-101, 105-111, 113, 114, 116-128, 130,}
\DLine{AB01-PDF0660-notes_left-L002}{}{132-136, 139, 140, 142, 144, 145, 149, 158 (?), 171,}
\DLine{AB01-PDF0660-notes_left-L003}{}{173, 179, 181, 182, 191, 199, 200, 204, 205 (?), 206 (?),}
\DLine{AB01-PDF0660-notes_left-L004}{}{212-215, 217–220, 223-225, 232, 272, 273.}
\hypertarget{AB01-PDF0660-N02}{}
\DLine{AB01-PDF0660-notes_left-L005}{}{\Indent (2) Nr. 115, 143, 152, 164, 168, 169 (?), 170,}
\DLine{AB01-PDF0660-notes_left-L006}{}{174-177, 183 (?), 184–187, 188 (?), 190, 194-196, 198,}
\DLine{AB01-PDF0660-notes_left-L007}{}{210, 212, 227, 234-248, 250, 252, 254–261, 262 (?),}
\DLine{AB01-PDF0660-notes_left-L008}{}{263, 264, 267-269, 271.}
\hypertarget{AB01-PDF0660-N03}{}
\DLine{AB01-PDF0660-notes_left-L009}{}{\Indent (3) Nr. 103, 104, 129, 131, 137, 138, 146, 147,}
\DLine{AB01-PDF0660-notes_left-L010}{}{148 (?), 150 (ex Albatenii observationibus), 151, 153-}
\DLine{AB01-PDF0660-notes_left-L011}{}{157, 159–163, 165–167, 172, 178, 180, 189, 192, 193,}
\end{minipage}
\hfill\begin{minipage}[t]{.48\textwidth}\vspace{0pt}\fontsize{9}{11.5}\selectfont\setlength{\GutterOffset}{0pt}
\DLine{AB01-PDF0660-notes_right-L001}{}{197, 201, 202, 207, 208, 209, 211 (?), 221 (?), 222 (?),}
\DLine{AB01-PDF0660-notes_right-L002}{}{226, 228, 229 (?), 251 (?), 253, 265, 266, 270.}
\hypertarget{AB01-PDF0660-N04}{}
\DLine{AB01-PDF0660-notes_right-L003}{}{\Indent (4) Litteris inclinatis indico nomina Graeca}
\DLine{AB01-PDF0660-notes_right-L004}{}{Arabibus fere omnibus ignota.}
\hypertarget{AB01-PDF0660-N05}{}
\DLine{AB01-PDF0660-notes_right-L005}{}{\Indent (5) Stellula * designo urbes quibus nihil in}
\DLine{AB01-PDF0660-notes_right-L006}{}{Geographia Ptolemaei respondet.}
\hypertarget{AB01-PDF0660-N06}{}
\DLine{AB01-PDF0660-notes_right-L007}{}{\Indent (6) Versio Syriaca etiam X saec. Christi ex-}
\DLine{AB01-PDF0660-notes_right-L008}{}{stabat, teste \textit{Kitāb al-Fihrist} p. 268. — Prima ver-}
\DLine{AB01-PDF0660-notes_right-L009}{}{sio Arabica illa videtur quam, priore IX saec. Chr.,}
\DLine{AB01-PDF0660-notes_right-L010}{}{in usum suum conscribendam \Name{al-Kindī} curaverat}
\DLine{AB01-PDF0660-notes_right-L011}{}{(\textit{Fihrist} l. c.).}
\end{minipage}
\DSource{0661}{212}
\DLine{AB01-PDF0661-body-L001}{}{amussim Ptolemaicis respondent, excepta parva differentia in notis ad nr. 15 et 16 indicata;}
\DLine{AB01-PDF0661-body-L002}{}{deest \textit{Africae} \textgreek{Μερόη} (quae tamen in conspectu generali IV libri Ptolemaei pro \textgreek{ἐπαρχία} ab}
\DLine{AB01-PDF0661-body-L003}{}{aliis distincta non habetur), sed Numidia, ut in IV libri conspectu, ab Africa provincia distin-}
\DLine{AB01-PDF0661-body-L004}{}{guitur, qua re regionum numerus semper est 12. Discrepantiae et errores quoque in provin-}
\DLine{AB01-PDF0661-body-L005}{5}{ciis \textit{Asiaticis} occurrunt. Phrygia et Paphlagonia, quae apud Ptolemaeum pertinent altera ad}
\DLine{AB01-PDF0661-body-L006}{}{Asiam Propriam altera ad Galatiam, duae peculiares provinciae in tabulis Albatenii fiunt;}
\DLine{AB01-PDF0661-body-L007}{}{earumque additio compensat Arabiam Desertam et Mesopotamiam miro modo omissas. Duo}
\DLine{AB01-PDF0661-body-L008}{}{errores parvi momenti in nr. 51 et 66; longe maior in nr. 92, quia Ḥimyaritarum regio iam}
\DLine{AB01-PDF0661-body-L009}{}{nr. 75 continetur, et hoc modo provinciae Asiaticae 47 loco 46 fiunt. — Alia quoque ani-}
\DLine{AB01-PDF0661-body-L010}{10}{madvertenda sunt. Provinciae \textit{Europae} ac \textit{Africae,} item \textit{Asiae} usque ad nr. 68, Graecis no-}
\DLine{AB01-PDF0661-body-L011}{}{minibus immutatis aut summum Arabice versis designantur; at inde a nr. 69 loco Graecorum}
\DLine{AB01-PDF0661-body-L012}{}{nomina provinciarum ab Arabibus usurpata ponuntur. Non semper recte, nam Ṭabaristān in}
\DLine{AB01-PDF0661-body-L013}{}{extremo Oriente collocatur ad \textgreek{Σηρική} repraesentandam (nr. 83), et regionis Farghānah latitudo}
\DLine{AB01-PDF0661-body-L014}{}{septem gradibus minuitur ut Paropanisadis respondeat. — In latitudinibus longitudinibusque}
\DLine{AB01-PDF0661-body-L015}{15}{provinciarum a 21 ad 93 codex minuta omittit; qua de causa nescio. — Situs provinciarum}
\DLine{AB01-PDF0661-body-L016}{}{omnium sunt mediae quantitates inter extremas longitudines ac latitudines singulis provinciis}
\DLine{AB01-PDF0661-body-L017}{}{a \Name{Ptolemaeo} assignatas.}
\hypertarget{AB01-PDF0661-P01}{}
\DLine{AB01-PDF0661-body-L018}{}{\Indent Longe maioris ponderis quam provinciarum tabula est urbium series, in quibus Graeca}
\DLine{AB01-PDF0661-body-L019}{}{traditio novis Arabum studiis geographicis certat.}
\hypertarget{AB01-PDF0661-P02}{}
\DLine{AB01-PDF0661-body-L020}{20}{\Indent In \textit{Africa} urbes paucae memorantur, eaeque, praeter Aegyptias, nominibus ac situbus}
\DLine{AB01-PDF0661-body-L021}{}{omnino Ptolemaicis (1); Aegyptus vero nomina Arabica situsque Khuwārizmicos complectitur,}
\DLine{AB01-PDF0661-body-L022}{}{licet ambae Oases, Diospolis Magna et Alexandria originem Graecam redoleant. \textit{Europa,} seu}
\DLine{AB01-PDF0661-body-L023}{}{quatuor urbes Roma, Ṭrāqiyah (?), Athenae, Constantinopolis, omnino Ptolemaica, sicut \textit{Asia}}
\DLine{AB01-PDF0661-body-L024}{}{\textit{Minor} (excepta ‘Ammūriyah) et extrema \textit{Asia Orientalis} (Thinae ac Sera). \textit{India, Turke-}}
\DLine{AB01-PDF0661-body-L025}{25}{\textit{stān, Persia} ex al-Khuwārizmī sumptae, 10 gradibus semper ad longitudines adiectis (2); si-}
\DLine{AB01-PDF0661-body-L026}{}{tuum Ptolemaicorum nonnisi Herāt (nr. 191) et Nisibis in Herāt (nr. 232) supersunt, a reli-}
\DLine{AB01-PDF0661-body-L027}{}{quis Persiae urbibus ita multum discrepantes. In \textit{Armenia} et \textit{Caucasicis} regionibus auctor}
\DLine{AB01-PDF0661-body-L028}{}{al-Khuwārizmī ducem habet, sed eius coordinatas aliquanto mutat ut situs illos cum Ptole-}
\DLine{AB01-PDF0661-body-L029}{}{maicis Asiae Minoris et Khuwārizmicis Persiae aliquo modo componere possit. Sed miro modo,}
\DLine{AB01-PDF0661-body-L030}{30}{simul cum Dabīl, Ǵorzān, Taflis, Bardha‘ah aliisque apparet Madīnat al-Abwāb (nr. 200) seu}
\DLine{AB01-PDF0661-body-L031}{}{Albaniae Portae eodem situ quo apud Ptolemaeum memoratur!}
\hypertarget{AB01-PDF0661-P03}{}
\DLine{AB01-PDF0661-body-L032}{}{\Indent In \textit{Syria} potissimum, \textit{Mesopotamia} et \textit{Arabia} arduus auctoris labor apparet Ptolemaei}
\DLine{AB01-PDF0661-body-L033}{}{cum al-Khuwārizmī conciliandi. In \textit{Syria} et \textit{Mesopotamia} amplius dimidia pars situum Pto-}
\DLine{AB01-PDF0661-body-L034}{}{lemaici sunt (3); alii multi Arabicae originis sed Graecis accommodati (4); 11 tantum ex al-}
\DLine{AB01-PDF0661-body-L035}{35}{Khuwārizmī absque mutationibus (praeter 10° longitudinibus adiectos) sumpti, nempe ar-Ram-}
\DLine{AB01-PDF0661-body-L036}{}{leh, Āmid, Ḥarrān, Baghdād, Surra man ra’à, al-Kūfah, al-Baṣrah, Wāsiṭ, Ḥulwān, Beyrūt}
\DLine{AB01-PDF0661-body-L037}{}{ac Hīt. Itaque verum monstrum est situs Beyrūt (nr. 227), si cum aliis Syriae urbibus com-}
\DLine{AB01-PDF0661-body-L038}{}{paretur. — In \textit{Arabia} 17 situs ex al-Khuwārizmī sumpti; sed aṭ-Ṭā’if, al-Ǵār, al-Yamāmah}
\DLine{AB01-PDF0661-body-L039}{}{locusque non probatus urbis Mekkah aliam originem ostendunt, nec possunt cum Khuwāriz-}
\par\vspace{6pt}\hrule height.25pt\vspace{5pt}
\noindent\begin{minipage}[t]{.48\textwidth}\vspace{0pt}\fontsize{9}{11.5}\selectfont\setlength{\GutterOffset}{0pt}
\hypertarget{AB01-PDF0661-N01}{}
\DLine{AB01-PDF0661-notes_left-L001}{}{\Indent (1) \textit{Madīnat aṭ-Ṭīb} (nr. 238) excipienda, Grae-}
\DLine{AB01-PDF0661-notes_left-L002}{}{cae originis quidem sed ex al-Khuwārizmī sumpta,}
\DLine{AB01-PDF0661-notes_left-L003}{}{e quo etiam \textit{Donqolah} (nr. 257) provenit.}
\hypertarget{AB01-PDF0661-N02}{}
\DLine{AB01-PDF0661-notes_left-L004}{}{\Indent (2) Cuius rei infra causam exponam. — Miro}
\DLine{AB01-PDF0661-notes_left-L005}{}{modo incrementum longitudinis non acceperunt \textit{Tub-}}
\DLine{AB01-PDF0661-notes_left-L006}{}{\textit{bat} (nr. 186) et \textit{al-Mayd} (nr. 239); fortasse ex li-}
\DLine{AB01-PDF0661-notes_left-L007}{}{brariorum errore.}
\end{minipage}
\hfill\begin{minipage}[t]{.48\textwidth}\vspace{0pt}\fontsize{9}{11.5}\selectfont\setlength{\GutterOffset}{0pt}
\hypertarget{AB01-PDF0661-N03}{}
\DLine{AB01-PDF0661-notes_right-L001}{}{\Indent (3) Nr. 110–114, 119-128, 130, 132–136, 139,}
\DLine{AB01-PDF0661-notes_right-L002}{}{140, 142, 144, 145, 148, 149, 153, 156, 158, 163 (?),}
\DLine{AB01-PDF0661-notes_right-L003}{}{171, 173, 199, 202-205, 207, 211-213, 219, 222 (?),}
\DLine{AB01-PDF0661-notes_right-L004}{}{223-225, 272, 273.}
\hypertarget{AB01-PDF0661-N04}{}
\DLine{AB01-PDF0661-notes_right-L005}{}{\Indent (4) Nr. 129, 131, 137, 138, 141, 150, 151, 154,}
\DLine{AB01-PDF0661-notes_right-L006}{}{155, 157, 159-162, 180, 201, 206, 208, 209, 220, 221,}
\DLine{AB01-PDF0661-notes_right-L007}{}{226, 228, 229 (?).}
\end{minipage}
\DSource{0662}{213}
\DLine{AB01-PDF0662-body-L001}{}{micis componi. Mappae Ptolemaicae nonnisi Sāffārā (nr. 97), insula Serapidis (nr. 98) et for-}
\DLine{AB01-PDF0662-body-L002}{}{tasse nr. 230 superfuerunt.}
\hypertarget{AB01-PDF0662-P01}{}
\DLine{AB01-PDF0662-body-L003}{}{\Indent Satis patet auctorem \textit{Libri Figurae Terrae,} traditionem Graecam omnino respuere non}
\DLine{AB01-PDF0662-body-L004}{}{ausum, fontes suos non semper feliciter commiscuisse. Cuius rei vestigia certa etiam apud}
\DLine{AB01-PDF0662-body-L005}{5}{nonnullos seriores astronomos inveniuntur, qui tabulas geographicas vetustiorum auctorum fere}
\DLine{AB01-PDF0662-body-L006}{}{exscribunt, incerta ac inutilia interdum servantes sicut in notis ad nr. 234, 235, 238, 239,}
\DLine{AB01-PDF0662-body-L007}{}{240 ostendi.}
\hypertarget{AB01-PDF0662-P02}{}
\DLine{AB01-PDF0662-body-L008}{}{\Indent Tribus locis, idest Mekkah (nr. 103), ar-Raqqah (nr. 150), Baghdād (nr. 168), additur}
\DLine{AB01-PDF0662-body-L009}{}{\textit{mumtaḥan} « [situs] probatus, verificatus ». Procul dubio observationibus astronomicis suis}
\DLine{AB01-PDF0662-body-L010}{10}{quoad ar-Raqqah; aliorum quoad Baghdād, nam huius urbis coordinatae eaedem sunt ac}
\DLine{AB01-PDF0662-body-L011}{}{Khuwārizmicae, et longitudo ita erat determinata iam ab astronomis khalifae al-Ma’mūn in}
\DLine{AB01-PDF0662-body-L012}{}{\textit{az-zīǵ al-mumtaḥan} « opere astronomico probato » sicut \Name{al-Mas‘ūdī}, \textit{at-Tanbīh} p. 45 te-}
\DLine{AB01-PDF0662-body-L013}{}{statur. Situs Mekkah ab ipso Albatenio constitutus videtur, sed longitudinem in codice errore}
\DLine{AB01-PDF0662-body-L014}{}{affectam puto. — Cur autem situm Antiochiae iuxta suas observationes al-Battānī non emen-}
\DLine{AB01-PDF0662-body-L015}{15}{daverit nescio.}
\hypertarget{AB01-PDF0662-P03}{}
\DLine{AB01-PDF0662-body-L016}{}{\Indent Restat ut de \textit{meridiano primo} nonnulla dicam. Dissertatione mea \textit{Al-Ḫuwârizmî e il}}
\DLine{AB01-PDF0662-body-L017}{}{\textit{suo rifacimento} etc., p. 24-25, ostendi primum meridianum Khuwārizmicum eundem esse ac}
\DLine{AB01-PDF0662-body-L018}{}{Ptolemaicum \textit{Insularum Fortunatarum;} sed, cum felicissime Arabs 9° vel 10° minuisset}
\DLine{AB01-PDF0662-body-L019}{}{longitudinem Ptolemaicam Maris Mediterranei et 10° circiter etiam situs ab ortu huius ma-}
\DLine{AB01-PDF0662-body-L020}{20}{ris, false seriores Geographos Arabicos putavisse \textit{primum meridianum Khuwārizmicum a}}
\DLine{AB01-PDF0662-body-L021}{}{\textit{Ptolemaico differre} et per littora \textit{Africae Occidentalis,} 10° ab ortu Insularum Fortunata-}
\DLine{AB01-PDF0662-body-L022}{}{rum, transire. — Dicit \Name{al-Battānī} t. I, p. 19 et 20, longitudines « ab Insulis Habitatis}
\DLine{AB01-PDF0662-body-L023}{}{« (\textit{al-ǵazā’ir al-‘āmirah}), quae sunt mari Oceani Occidentalis (\textit{ūqiyānus al-gharbī}), » nume-}
\DLine{AB01-PDF0662-body-L024}{}{rari; sed his tabulis longitudines omnium locorum ex al-Khuwārizmī sumptorum semper 10°}
\DLine{AB01-PDF0662-body-L025}{25}{\textit{augentur} (1). Unde patet iam auctorem \textit{Libri Figurae Terrae} primum meridianum Khuwā-}
\DLine{AB01-PDF0662-body-L026}{}{rizmicum a Ptolemaico diversum perperam existimavisse, et hoc modo novum eumque falsum}
\DLine{AB01-PDF0662-body-L027}{}{meridianum Insularum Fortunatarum instituisse, quo irrita facta est optima emendatio ab al-}
\DLine{AB01-PDF0662-body-L028}{}{Khuwārizmī longitudini Mediterranei illata. Potuit ita quidem auctor \textit{Libri Figurae Terrae}}
\DLine{AB01-PDF0662-body-L029}{}{satis bene longitudines Khuwārizmicas cum Ptolemaicis conciliare; sed eandem rem assecutus}
\DLine{AB01-PDF0662-body-L030}{30}{esset, et longe melius, si, longitudinibus al-Khuwārizmī servatis, Ptolemaicas a Mediterranei}
\DLine{AB01-PDF0662-body-L031}{}{ortu 10 gradibus fere minuisset.}
\hypertarget{AB01-PDF0662-P04}{}
\DLine{AB01-PDF0662-body-L032}{}{\Indent Cum his tabulis nullo certo ordine urbes memorentur, earum, per regiones distributarum,}
\DLine{AB01-PDF0662-body-L033}{}{conspectum addere visum est.}
\par\vspace{4mm}
\par\vspace{2pt}\noindent{\lineskiplimit=-\maxdimen\begin{minipage}[t]{51.4mm}\vspace{0pt}
\DLine{AB01-PDF0662-body-L034}{}{\hypertarget{AB01-PDF0662-H34}{}\CenterLine{\fontsize{8.5}{10}\selectfont\textbf{Europa.}}}
\par\vspace{1\baselineskip}
\DLine{AB01-PDF0662-body-L035}{}{Athenae 217.}
\DLine{AB01-PDF0662-body-L036}{}{Constantinopolis 182.}
\DLine{AB01-PDF0662-body-L037}{}{Roma 181.}
\DLine{AB01-PDF0662-body-L038}{}{Ṭrāqiyah (?) 218.}
\par\vspace{0.5\baselineskip}
\DLine{AB01-PDF0662-body-L039}{}{\hypertarget{AB01-PDF0662-H39}{}\CenterLine{\fontsize{8.5}{10}\selectfont\textbf{Africa.}}}
\par\vspace{0.5\baselineskip}
\DLine{AB01-PDF0662-body-L040}{}{Akhmīm 250.}
\DLine{AB01-PDF0662-body-L041}{}{Alexandria 109.}
\end{minipage}\begin{minipage}[t]{51.4mm}\vspace{0pt}
\DLine{AB01-PDF0662-body-L042}{}{Auxume 215.}
\DLine{AB01-PDF0662-body-L043}{}{Catabathmus magnus 107.}
\DLine{AB01-PDF0662-body-L044}{}{Carthago 105.}
\DLine{AB01-PDF0662-body-L045}{}{Dimyāṭ 177.}
\DLine{AB01-PDF0662-body-L046}{}{Diospolis magna 100.}
\DLine{AB01-PDF0662-body-L047}{}{Donqolah 257.}
\DLine{AB01-PDF0662-body-L048}{}{al-Faramā 195.}
\DLine{AB01-PDF0662-body-L049}{}{al-Fusṭāṭ 178.}
\DLine{AB01-PDF0662-body-L050}{}{Gira metropolis 94.}
\DLine{AB01-PDF0662-body-L051}{}{Iarzitha (? \textarabic{سونا ثليا}) 96.}
\end{minipage}\begin{minipage}[t]{51.4mm}\vspace{0pt}
\DLine{AB01-PDF0662-body-L052}{}{Leptis magna 106.}
\DLine{AB01-PDF0662-body-L053}{}{Madīnat aṭ-Ṭīb 238.}
\DLine{AB01-PDF0662-body-L054}{}{Nigira 95.}
\DLine{AB01-PDF0662-body-L055}{}{Oasis maior 108.}
\DLine{AB01-PDF0662-body-L056}{}{Oasis minor 101.}
\DLine{AB01-PDF0662-body-L057}{}{al-Qolzum 252.}
\DLine{AB01-PDF0662-body-L058}{}{Qūṣ 251.}
\DLine{AB01-PDF0662-body-L059}{}{Syene 188.}
\DLine{AB01-PDF0662-body-L060}{}{Tinnīs 194.}
\end{minipage}}\par
\par\vspace{6pt}\hrule height.25pt\vspace{5pt}{\fontsize{9}{11.5}\selectfont
\hypertarget{AB01-PDF0662-N01}{}
\DLine{AB01-PDF0662-notes-L001}{}{\Indent (1) Cfr. p. 21, adn. 2.}
}
\DSource{0663}{214}
\par\vspace{2pt}\noindent{\lineskiplimit=-\maxdimen\begin{minipage}[t]{51.4mm}\vspace{0pt}
\par\vspace{0.5\baselineskip}
\DLine{AB01-PDF0663-body-L001}{}{\hypertarget{AB01-PDF0663-H01}{}\CenterLine{\fontsize{8.5}{10}\selectfont\textbf{Asia Minor.}}}
\par\vspace{0.5\baselineskip}
\DLine{AB01-PDF0663-body-L002}{}{Adhanah 120.}
\DLine{AB01-PDF0663-body-L003}{}{‘Ammūriyah 183.}
\DLine{AB01-PDF0663-body-L004}{}{‘Ayn Zarbah 179.}
\DLine{AB01-PDF0663-body-L005}{}{Laodicea Phrygia 116.}
\DLine{AB01-PDF0663-body-L006}{}{Malaṭyah 143.}
\DLine{AB01-PDF0663-body-L007}{}{al-Maṣṣīṣah 121.}
\DLine{AB01-PDF0663-body-L008}{}{Rhodus 117.}
\DLine{AB01-PDF0663-body-L009}{}{Salamis Cypri (?) 118.}
\DLine{AB01-PDF0663-body-L010}{}{Ṭarasūs 119.}
\DLine{AB01-PDF0663-body-L011}{}{Trapezus 265.}
\DLine{AB01-PDF0663-body-L012}{}{Zibaṭrah 214.}
\par\vspace{0.5\baselineskip}
\DLine{AB01-PDF0663-body-L013}{}{\hypertarget{AB01-PDF0663-H13}{}\CenterLine{\fontsize{8.5}{10}\selectfont\textbf{Armenia.}}}
\par\vspace{0.5\baselineskip}
\DLine{AB01-PDF0663-body-L014}{}{Bardha‘ah 167.}
\DLine{AB01-PDF0663-body-L015}{}{Dabīl 165.}
\DLine{AB01-PDF0663-body-L016}{}{Ǵorzān 187.}
\DLine{AB01-PDF0663-body-L017}{}{Khilāṭ 164.}
\DLine{AB01-PDF0663-body-L018}{}{Madīnat al-abwāb 200.}
\DLine{AB01-PDF0663-body-L019}{}{Taflīs 166.}
\par\vspace{0.5\baselineskip}
\DLine{AB01-PDF0663-body-L020}{}{\hypertarget{AB01-PDF0663-H20}{}\CenterLine{\fontsize{8.5}{10}\selectfont\textbf{Arabia.}}}
\par\vspace{0.5\baselineskip}
\DLine{AB01-PDF0663-body-L021}{}{‘Aden 185.}
\DLine{AB01-PDF0663-body-L022}{}{Alqis (Ocelis?) 234.}
\DLine{AB01-PDF0663-body-L023}{}{al-Baḥrayn 246.}
\DLine{AB01-PDF0663-body-L024}{}{\textarabic{بِرهور} 236.}
\DLine{AB01-PDF0663-body-L025}{}{al-Ǵār 253.}
\DLine{AB01-PDF0663-body-L026}{}{Ǵorash 243.}
\DLine{AB01-PDF0663-body-L027}{}{Ḥaḍramūt 237.}
\DLine{AB01-PDF0663-body-L028}{}{Haǵar 254.}
\DLine{AB01-PDF0663-body-L029}{}{Mārā 235.}
\DLine{AB01-PDF0663-body-L030}{}{Mahrah 244.}
\DLine{AB01-PDF0663-body-L031}{}{Mekkah 103.}
\DLine{AB01-PDF0663-body-L032}{}{\textarabic{مصره} 249.}
\DLine{AB01-PDF0663-body-L033}{}{\textarabic{المعلى} 240.}
\DLine{AB01-PDF0663-body-L034}{}{‘Omān 247.}
\DLine{AB01-PDF0663-body-L035}{}{Saba’ 242.}
\DLine{AB01-PDF0663-body-L036}{}{Ṣan‘ā’ 184.}
\DLine{AB01-PDF0663-body-L037}{}{Sapphara 97.}
\DLine{AB01-PDF0663-body-L038}{}{Serapidis insula 98.}
\DLine{AB01-PDF0663-body-L039}{}{Tabālah 245.}
\DLine{AB01-PDF0663-body-L040}{}{aṭ-Ṭā’if 193.}
\DLine{AB01-PDF0663-body-L041}{}{al-Yamāmah 192.}
\DLine{AB01-PDF0663-body-L042}{}{Yathrib (Medina) 104.}
\DLine{AB01-PDF0663-body-L043}{}{Ẓafār 241.}
\end{minipage}\begin{minipage}[t]{51.4mm}\vspace{0pt}
\par\vspace{0.5\baselineskip}
\DLine{AB01-PDF0663-body-L044}{}{\hypertarget{AB01-PDF0663-H44}{}\CenterLine{\fontsize{8.5}{10}\selectfont\textbf{Syria et Palaestina.}}}
\par\vspace{0.5\baselineskip}
\DLine{AB01-PDF0663-body-L045}{}{‘Akkā 127.}
\DLine{AB01-PDF0663-body-L046}{}{‘Amawās 211.}
\DLine{AB01-PDF0663-body-L047}{}{Antiochia 142.}
\DLine{AB01-PDF0663-body-L048}{}{Aradus 272.}
\DLine{AB01-PDF0663-body-L049}{}{Asdūd 213.}
\DLine{AB01-PDF0663-body-L050}{}{‘Asqalān 113.}
\DLine{AB01-PDF0663-body-L051}{}{Ba‘labakk 134.}
\DLine{AB01-PDF0663-body-L052}{}{Bayt al-Maqdis (Ierusalem) 273.}
\DLine{AB01-PDF0663-body-L053}{}{Bayt Ǵibrīn 228.}
\DLine{AB01-PDF0663-body-L054}{}{Beyrūt 227.}
\DLine{AB01-PDF0663-body-L055}{}{Caesarea Panias 110.}
\DLine{AB01-PDF0663-body-L056}{}{Damascus 133.}
\DLine{AB01-PDF0663-body-L057}{}{Dolūk 140.}
\DLine{AB01-PDF0663-body-L058}{}{Fāmiyah 132.}
\DLine{AB01-PDF0663-body-L059}{}{Filasṭīn 111.}
\DLine{AB01-PDF0663-body-L060}{}{Ǵabalah 224.}
\DLine{AB01-PDF0663-body-L061}{}{Ǵindāros 220.}
\DLine{AB01-PDF0663-body-L062}{}{Ǵisr Anṭākiyah 180.}
\DLine{AB01-PDF0663-body-L063}{}{Ǵubayl 202.}
\DLine{AB01-PDF0663-body-L064}{}{Ḥaleb (Aleppo) 136.}
\DLine{AB01-PDF0663-body-L065}{}{Ḥamāh 130.}
\DLine{AB01-PDF0663-body-L066}{}{Ḥimṣ 128.}
\DLine{AB01-PDF0663-body-L067}{}{Hīt 271.}
\DLine{AB01-PDF0663-body-L068}{}{Ḥuwwārīn 208.}
\DLine{AB01-PDF0663-body-L069}{}{‘Irqah (‘Arqah) 124.}
\DLine{AB01-PDF0663-body-L070}{}{al-Iskandarūnah 219.}
\DLine{AB01-PDF0663-body-L071}{}{Laodicea 122.}
\DLine{AB01-PDF0663-body-L072}{}{Ma‘arrat an-Nu‘mān 138.}
\DLine{AB01-PDF0663-body-L073}{}{Orthosia 222.}
\DLine{AB01-PDF0663-body-L074}{}{Qārā 226.}
\DLine{AB01-PDF0663-body-L075}{}{Qinnasrīn 137.}
\DLine{AB01-PDF0663-body-L076}{}{Qūros 139.}
\DLine{AB01-PDF0663-body-L077}{}{Ra‘bān 141.}
\DLine{AB01-PDF0663-body-L078}{}{Rāfiyah 212.}
\DLine{AB01-PDF0663-body-L079}{}{ar-Ramleh 115.}
\DLine{AB01-PDF0663-body-L080}{}{ar-Rastan 129.}
\DLine{AB01-PDF0663-body-L081}{}{ar-Ruṣāfah 201.}
\DLine{AB01-PDF0663-body-L082}{}{Rūsis 225.}
\DLine{AB01-PDF0663-body-L083}{}{Sabasṭiyah 114.}
\DLine{AB01-PDF0663-body-L084}{}{Salamiyyah 131.}
\DLine{AB01-PDF0663-body-L085}{}{Sanǵah 223.}
\DLine{AB01-PDF0663-body-L086}{}{Ṣaydā’ (Sidon) 126.}
\DLine{AB01-PDF0663-body-L087}{}{Shayzar 206.}
\DLine{AB01-PDF0663-body-L088}{}{Ṣūr (Tyrus) 125.}
\end{minipage}\begin{minipage}[t]{54mm}\vspace{0pt}
\DLine{AB01-PDF0663-body-L089}{}{Tadmur (Palmyra) 135.}
\DLine{AB01-PDF0663-body-L090}{}{Tell Mannas 207.}
\DLine{AB01-PDF0663-body-L091}{}{Tripolis 123.}
\DLine{AB01-PDF0663-body-L092}{}{Ūrim 204.}
\DLine{AB01-PDF0663-body-L093}{}{Waǵh al-ḥaǵar 221.}
\DLine{AB01-PDF0663-body-L094}{}{Zoghmah 205.}
\par\vspace{0.5\baselineskip}
\DLine{AB01-PDF0663-body-L095}{}{\hypertarget{AB01-PDF0663-H95}{}\makebox[51.4mm][c]{\fontsize{8.5}{10}\selectfont\textbf{Mesopotamia.}}}
\par\vspace{0.5\baselineskip}
\DLine{AB01-PDF0663-body-L096}{}{‘Abbādān 268.}
\DLine{AB01-PDF0663-body-L097}{}{Āmid 146.}
\DLine{AB01-PDF0663-body-L098}{}{al-‘Āqūl 209.}
\DLine{AB01-PDF0663-body-L099}{}{Arzan 147.}
\DLine{AB01-PDF0663-body-L100}{}{Babel 171.}
\DLine{AB01-PDF0663-body-L101}{}{Baghdād 168.}
\DLine{AB01-PDF0663-body-L102}{}{Balad 161.}
\DLine{AB01-PDF0663-body-L103}{}{Bālis 149.}
\DLine{AB01-PDF0663-body-L104}{}{al-Baṣrah 174.}
\DLine{AB01-PDF0663-body-L105}{}{Ctesiphon (al-Madā’in) 199.}
\DLine{AB01-PDF0663-body-L106}{}{Dārā 159.}
\DLine{AB01-PDF0663-body-L107}{}{Ḥarrān 152.}
\DLine{AB01-PDF0663-body-L108}{}{Kafar Tūthā 157.}
\DLine{AB01-PDF0663-body-L109}{}{al-Kūfah 170.}
\DLine{AB01-PDF0663-body-L110}{}{Manbiǵ 154.}
\DLine{AB01-PDF0663-body-L111}{}{Māridīn 160.}
\DLine{AB01-PDF0663-body-L112}{}{al-Mawṣil (Mossul) 162.}
\DLine{AB01-PDF0663-body-L113}{}{Mayyāfāriqīn 145.}
\DLine{AB01-PDF0663-body-L114}{}{Naṣībīn 158.}
\DLine{AB01-PDF0663-body-L115}{}{Ninive 173.}
\DLine{AB01-PDF0663-body-L116}{}{Qirqīsiyā’ 151.}
\DLine{AB01-PDF0663-body-L117}{}{ar-Raqqah 150.}
\DLine{AB01-PDF0663-body-L118}{}{Ra’s al-‘ayn 156.}
\DLine{AB01-PDF0663-body-L119}{}{ar-Rohā’ (Edessa) 153.}
\DLine{AB01-PDF0663-body-L120}{}{Shimshāṭ 144.}
\DLine{AB01-PDF0663-body-L121}{}{Sinǵār 163.}
\DLine{AB01-PDF0663-body-L122}{}{Sumaysāṭ 148.}
\DLine{AB01-PDF0663-body-L123}{}{Sūrā 229.}
\DLine{AB01-PDF0663-body-L124}{}{Surra man ra’à (Sāmarrā) 169.}
\DLine{AB01-PDF0663-body-L125}{}{Tell Mawzan 155.}
\DLine{AB01-PDF0663-body-L126}{}{Wāsiṭ 176.}
\par\vspace{0.5\baselineskip}
\DLine{AB01-PDF0663-body-L127}{}{\hypertarget{AB01-PDF0663-H127}{}\makebox[51.4mm][c]{\fontsize{8.5}{10}\selectfont\textbf{Persia.}}}
\par\vspace{0.5\baselineskip}
\DLine{AB01-PDF0663-body-L128}{}{Āmul 263.}
\DLine{AB01-PDF0663-body-L129}{}{Donbāwend 262.}
\DLine{AB01-PDF0663-body-L130}{}{Ǵiruft 255.}
\DLine{AB01-PDF0663-body-L131}{}{Hamadhān 210.}
\DLine{AB01-PDF0663-body-L132}{}{Herāt 191.}
\end{minipage}}\par
\DSource{0664}{215}
\par\vspace{2pt}\noindent{\lineskiplimit=-\maxdimen\begin{minipage}[t]{51.4mm}\vspace{0pt}
\DLine{AB01-PDF0664-body-L001}{}{Ḥulwān 198.}
\DLine{AB01-PDF0664-body-L002}{}{Kābul 256.}
\DLine{AB01-PDF0664-body-L003}{}{Khuwayy 266.}
\DLine{AB01-PDF0664-body-L004}{}{al-Muḥammadiyyah 259.}
\DLine{AB01-PDF0664-body-L005}{}{Nisibis Ariae 232.}
\DLine{AB01-PDF0664-body-L006}{}{Qaṣr al-milḥ 260.}
\DLine{AB01-PDF0664-body-L007}{}{Qazwīn 190.}
\DLine{AB01-PDF0664-body-L008}{}{Qumm 197.}
\DLine{AB01-PDF0664-body-L009}{}{ar-Rayy 172.}
\DLine{AB01-PDF0664-body-L010}{}{ar-Rūyān 258.}
\DLine{AB01-PDF0664-body-L011}{}{Sarakhs 270.}
\end{minipage}\begin{minipage}[t]{51.4mm}\vspace{0pt}
\DLine{AB01-PDF0664-body-L012}{}{Sāriyah 264.}
\DLine{AB01-PDF0664-body-L013}{}{Sīrāf 175.}
\DLine{AB01-PDF0664-body-L014}{}{as-Sīraǵān 261.}
\DLine{AB01-PDF0664-body-L015}{}{Ṭūs 269.}
\par\vspace{0.5\baselineskip}
\DLine{AB01-PDF0664-body-L016}{}{\hypertarget{AB01-PDF0664-H16}{}\CenterLine{\fontsize{8.5}{10}\selectfont\textbf{India et Asia Orientalis.}}}
\par\vspace{0.5\baselineskip}
\DLine{AB01-PDF0664-body-L017}{}{ad-Daybul 189.}
\DLine{AB01-PDF0664-body-L018}{}{al-Mayd 239.}
\DLine{AB01-PDF0664-body-L019}{}{an-Nīrūn 248.}
\DLine{AB01-PDF0664-body-L020}{}{Osrūshanah 267.}
\DLine{AB01-PDF0664-body-L021}{}{Sīrās 231.}
\end{minipage}\begin{minipage}[t]{51.4mm}\vspace{0pt}
\DLine{AB01-PDF0664-body-L022}{}{Thinae 99.}
\DLine{AB01-PDF0664-body-L023}{}{Tubbat 186.}
\DLine{AB01-PDF0664-body-L024}{}{aṭ-Ṭurāraband 196.}
\par\vspace{0.5\baselineskip}
\DLine{AB01-PDF0664-body-L025}{}{\hypertarget{AB01-PDF0664-H25}{}\CenterLine{\fontsize{8.5}{10}\selectfont\textbf{Incertae.}}}
\par\vspace{0.5\baselineskip}
\DLine{AB01-PDF0664-body-L026}{}{\textarabic{بجلزا} seu \textarabic{بجلنرا} 102.}
\DLine{AB01-PDF0664-body-L027}{}{\textarabic{سقراطاوس} 112.}
\DLine{AB01-PDF0664-body-L028}{}{\textarabic{ارام} 230.}
\DLine{AB01-PDF0664-body-L029}{}{\textarabic{جزيل} 203.}
\DLine{AB01-PDF0664-body-L030}{}{\textarabic{بلد اور} 233.}
\DLine{AB01-PDF0664-body-L031}{}{\textarabic{داڡا مدينة الفرس} 216.}
\end{minipage}}\par
\par\vspace{4mm}
\hypertarget{AB01-PDF0664-P01}{}
\DLine{AB01-PDF0664-body-L032}{}{\Indent Qua ratione codicis numeros emendaverim ex ipsis notis patet; consulantur quoque \textit{Ad-}}
\DLine{AB01-PDF0664-body-L033}{}{\textit{denda et Emendanda.} Quae iam \Name{Lelewel} feliciter coniecerat semper suis locis memoravi;}
\DLine{AB01-PDF0664-body-L034}{}{illius praeclari viri mendas etiam notare inutile duxi. \Name{Ptolemaei} \textit{Geographia} usus sum}
\DLine{AB01-PDF0664-body-L035}{35}{iuxta editionem C. Fr. A. Nobbe, ed. stereotypa, Lipsiae 1881 (3 tomi in-18°); \Name{Ibn Yūnus}}
\DLine{AB01-PDF0664-body-L036}{}{ad fidem codicis Lugduno-Batavi (p. 133–36) adhibui. Nomina Graeca nunquam apud Ara-}
\DLine{AB01-PDF0664-body-L037}{}{bicos scriptores vulgata semper stellula * instruxi.}
\hypertarget{AB01-PDF0664-P02}{}
\DLine{AB01-PDF0664-body-L038}{}{\Indent Versio Hispanica (f. 42,v.), quam hic referre iuvat, 129 urbes ex tabula Albateniana}
\DLine{AB01-PDF0664-body-L039}{}{excerpsit, tribus aliis Hispaniae urbibus praemissis quae apud al-Battānī desunt. Urbibus addo}
\DLine{AB01-PDF0664-body-L040}{40}{ipsarum numeros ordinis iuxta textum Arabicum; longitudines et latitudines falsas discrepan-}
\DLine{AB01-PDF0664-body-L041}{}{tesque a codice Escurialensi stellula * instruo, falsas sed codici Escurialensi communes litteris}
\DLine{AB01-PDF0664-body-L042}{}{inclinatis seu Italicis indico.}
\par\vspace{14mm}
\par\noindent\rule{157mm}{1.2pt}\par\nointerlineskip\vspace{1pt}\noindent\rule{157mm}{.4pt}\par
\par\nointerlineskip\noindent\hspace*{0.0mm}\hspace{36.6mm}\hspace{3mm}\rule{.4pt}{4mm}\hspace{2mm}\hspace{8.0mm}\hspace{7.0mm}\hspace{3mm}\rule{.4pt}{4mm}\hspace{2mm}\hspace{6.0mm}\hspace{7.5mm}\hspace{3mm}\rule{.4pt}{4mm}\hspace{1.2pt}\rule{.4pt}{4mm}\hspace{2mm}\hspace{36.2mm}\hspace{3mm}\rule{.4pt}{4mm}\hspace{2mm}\hspace{8.0mm}\hspace{7.0mm}\hspace{3mm}\rule{.4pt}{4mm}\hspace{2mm}\hspace{6.0mm}\hspace{6.9mm}\par\nointerlineskip
\DLine{AB01-PDF0664-body-L043}{}{\hspace*{0.0mm}\makebox[36.6mm][l]{\smash{Cordoua despanna}}\hspace{3mm}\rule[-\dp\strutbox]{.4pt}{\baselineskip}\hspace{2mm}\makebox[8.0mm][r]{9\rlap{°}}\makebox[7.0mm][r]{9\rlap{′}}\hspace{3mm}\rule[-\dp\strutbox]{.4pt}{\baselineskip}\hspace{2mm}\makebox[6.0mm][r]{38\rlap{°}}\makebox[7.5mm][r]{20\rlap{′}}\hspace{3mm}\rule[-\dp\strutbox]{.4pt}{\baselineskip}\hspace{1.2pt}\rule[-\dp\strutbox]{.4pt}{\baselineskip}\hspace{2mm}\makebox[36.2mm][l]{\smash{56. Tarçoç}}\hspace{3mm}\rule[-\dp\strutbox]{.4pt}{\baselineskip}\hspace{2mm}\makebox[8.0mm][r]{68\rlap{°}}\makebox[7.0mm][r]{\smash{}}\hspace{3mm}\rule[-\dp\strutbox]{.4pt}{\baselineskip}\hspace{2mm}\makebox[6.0mm][r]{37\rlap{°}}\makebox[6.9mm][r]{\smash{}}}
\DLine{AB01-PDF0664-body-L044}{}{\hspace*{0.0mm}\makebox[36.6mm][l]{\smash{Toledo}}\hspace{3mm}\rule[-\dp\strutbox]{.4pt}{\baselineskip}\hspace{2mm}\makebox[8.0mm][r]{\smash{8}}\makebox[7.0mm][r]{\smash{}}\hspace{3mm}\rule[-\dp\strutbox]{.4pt}{\baselineskip}\hspace{2mm}\makebox[6.0mm][r]{\smash{39}}\makebox[7.5mm][r]{\smash{20}}\hspace{3mm}\rule[-\dp\strutbox]{.4pt}{\baselineskip}\hspace{1.2pt}\rule[-\dp\strutbox]{.4pt}{\baselineskip}\hspace{2mm}\makebox[36.2mm][l]{\smash{62. ysla de cabroz}}\hspace{3mm}\rule[-\dp\strutbox]{.4pt}{\baselineskip}\hspace{2mm}\makebox[8.0mm][r]{\smash{66}}\makebox[7.0mm][r]{\smash{}}\hspace{3mm}\rule[-\dp\strutbox]{.4pt}{\baselineskip}\hspace{2mm}\makebox[6.0mm][r]{\smash{35}}\makebox[6.9mm][r]{\smash{\textit{30}}}}
\DLine{AB01-PDF0664-body-L045}{45}{\hspace*{0.0mm}\makebox[36.6mm][l]{\smash{çaragoça despanna}}\hspace{3mm}\rule[-\dp\strutbox]{.4pt}{\baselineskip}\hspace{2mm}\makebox[8.0mm][r]{\smash{11}}\makebox[7.0mm][r]{\smash{}}\hspace{3mm}\rule[-\dp\strutbox]{.4pt}{\baselineskip}\hspace{2mm}\makebox[6.0mm][r]{\smash{42}}\makebox[7.5mm][r]{\smash{}}\hspace{3mm}\rule[-\dp\strutbox]{.4pt}{\baselineskip}\hspace{1.2pt}\rule[-\dp\strutbox]{.4pt}{\baselineskip}\hspace{2mm}\makebox[36.2mm][l]{\smash{63. Çorya}}\hspace{3mm}\rule[-\dp\strutbox]{.4pt}{\baselineskip}\hspace{2mm}\makebox[8.0mm][r]{\smash{71}}\makebox[7.0mm][r]{\smash{}}\hspace{3mm}\rule[-\dp\strutbox]{.4pt}{\baselineskip}\hspace{2mm}\makebox[6.0mm][r]{\smash{36}}\makebox[6.9mm][r]{\smash{}}}
\DLine{AB01-PDF0664-body-L046}{}{\hspace*{0.0mm}\makebox[36.6mm][l]{\smash{16. Dalmathya}}\hspace{3mm}\rule[-\dp\strutbox]{.4pt}{\baselineskip}\hspace{2mm}\makebox[8.0mm][r]{\smash{46}}\makebox[7.0mm][r]{\smash{}}\hspace{3mm}\rule[-\dp\strutbox]{.4pt}{\baselineskip}\hspace{2mm}\makebox[6.0mm][r]{\smash{42}}\makebox[7.5mm][r]{\smash{}}\hspace{3mm}\rule[-\dp\strutbox]{.4pt}{\baselineskip}\hspace{1.2pt}\rule[-\dp\strutbox]{.4pt}{\baselineskip}\hspace{2mm}\makebox[36.2mm][l]{\smash{66. Tierra de arabigos}}\hspace{3mm}\rule[-\dp\strutbox]{.4pt}{\baselineskip}\hspace{2mm}\makebox[8.0mm][r]{\smash{*60}}\makebox[7.0mm][r]{\smash{}}\hspace{3mm}\rule[-\dp\strutbox]{.4pt}{\baselineskip}\hspace{2mm}\makebox[6.0mm][r]{\smash{29}}\makebox[6.9mm][r]{\smash{*31}}}
\DLine{AB01-PDF0664-body-L047}{}{\hspace*{0.0mm}\makebox[36.6mm][l]{\smash{17. ysla de ythalia}}\hspace{3mm}\rule[-\dp\strutbox]{.4pt}{\baselineskip}\hspace{2mm}\makebox[8.0mm][r]{\smash{36}}\makebox[7.0mm][r]{\smash{40}}\hspace{3mm}\rule[-\dp\strutbox]{.4pt}{\baselineskip}\hspace{2mm}\makebox[6.0mm][r]{\smash{41}}\makebox[7.5mm][r]{\smash{40}}\hspace{3mm}\rule[-\dp\strutbox]{.4pt}{\baselineskip}\hspace{1.2pt}\rule[-\dp\strutbox]{.4pt}{\baselineskip}\hspace{2mm}\makebox[36.2mm][l]{\smash{68. Assur}}\hspace{3mm}\rule[-\dp\strutbox]{.4pt}{\baselineskip}\hspace{2mm}\makebox[8.0mm][r]{\smash{80}}\makebox[7.0mm][r]{\smash{}}\hspace{3mm}\rule[-\dp\strutbox]{.4pt}{\baselineskip}\hspace{2mm}\makebox[6.0mm][r]{\smash{37}}\makebox[6.9mm][r]{\smash{}}}
\DLine{AB01-PDF0664-body-L048}{}{\hspace*{0.0mm}\makebox[36.6mm][l]{\smash{19. ysla de sardenna}}\hspace{3mm}\rule[-\dp\strutbox]{.4pt}{\baselineskip}\hspace{2mm}\makebox[8.0mm][r]{\smash{31}}\makebox[7.0mm][r]{\smash{*}}\hspace{3mm}\rule[-\dp\strutbox]{.4pt}{\baselineskip}\hspace{2mm}\makebox[6.0mm][r]{\smash{38}}\makebox[7.5mm][r]{\smash{*}}\hspace{3mm}\rule[-\dp\strutbox]{.4pt}{\baselineskip}\hspace{1.2pt}\rule[-\dp\strutbox]{.4pt}{\baselineskip}\hspace{2mm}\makebox[36.2mm][l]{\smash{70. açuç}}\hspace{3mm}\rule[-\dp\strutbox]{.4pt}{\baselineskip}\hspace{2mm}\makebox[8.0mm][r]{\smash{83}}\makebox[7.0mm][r]{\smash{}}\hspace{3mm}\rule[-\dp\strutbox]{.4pt}{\baselineskip}\hspace{2mm}\makebox[6.0mm][r]{\smash{34}}\makebox[6.9mm][r]{\smash{}}}
\DLine{AB01-PDF0664-body-L049}{}{\hspace*{0.0mm}\makebox[36.6mm][l]{\smash{20. cezilia}}\hspace{3mm}\rule[-\dp\strutbox]{.4pt}{\baselineskip}\hspace{2mm}\makebox[8.0mm][r]{\smash{39}}\makebox[7.0mm][r]{\smash{*}}\hspace{3mm}\rule[-\dp\strutbox]{.4pt}{\baselineskip}\hspace{2mm}\makebox[6.0mm][r]{\smash{37}}\makebox[7.5mm][r]{\smash{*}}\hspace{3mm}\rule[-\dp\strutbox]{.4pt}{\baselineskip}\hspace{1.2pt}\rule[-\dp\strutbox]{.4pt}{\baselineskip}\hspace{2mm}\makebox[36.2mm][l]{\smash{71. perssia}}\hspace{3mm}\rule[-\dp\strutbox]{.4pt}{\baselineskip}\hspace{2mm}\makebox[8.0mm][r]{\smash{90}}\makebox[7.0mm][r]{\smash{}}\hspace{3mm}\rule[-\dp\strutbox]{.4pt}{\baselineskip}\hspace{2mm}\makebox[6.0mm][r]{\smash{32}}\makebox[6.9mm][r]{\smash{}}}
\DLine{AB01-PDF0664-body-L050}{50}{\hspace*{0.0mm}\makebox[36.6mm][l]{\smash{35. Taniar}}\hspace{3mm}\rule[-\dp\strutbox]{.4pt}{\baselineskip}\hspace{2mm}\makebox[8.0mm][r]{\smash{8}}\makebox[7.0mm][r]{\smash{}}\hspace{3mm}\rule[-\dp\strutbox]{.4pt}{\baselineskip}\hspace{2mm}\makebox[6.0mm][r]{\smash{32}}\makebox[7.5mm][r]{\smash{}}\hspace{3mm}\rule[-\dp\strutbox]{.4pt}{\baselineskip}\hspace{1.2pt}\rule[-\dp\strutbox]{.4pt}{\baselineskip}\hspace{2mm}\makebox[36.2mm][l]{\smash{72. azpahen}}\hspace{3mm}\rule[-\dp\strutbox]{.4pt}{\baselineskip}\hspace{2mm}\makebox[8.0mm][r]{\smash{96}}\makebox[7.0mm][r]{\smash{}}\hspace{3mm}\rule[-\dp\strutbox]{.4pt}{\baselineskip}\hspace{2mm}\makebox[6.0mm][r]{\smash{*34}}\makebox[6.9mm][r]{\smash{}}}
\DLine{AB01-PDF0664-body-L051}{}{\hspace*{0.0mm}\makebox[36.6mm][l]{\smash{37. affrica}}\hspace{3mm}\rule[-\dp\strutbox]{.4pt}{\baselineskip}\hspace{2mm}\makebox[8.0mm][r]{\smash{36}}\makebox[7.0mm][r]{\smash{}}\hspace{3mm}\rule[-\dp\strutbox]{.4pt}{\baselineskip}\hspace{2mm}\makebox[6.0mm][r]{\smash{31}}\makebox[7.5mm][r]{\smash{}}\hspace{3mm}\rule[-\dp\strutbox]{.4pt}{\baselineskip}\hspace{1.2pt}\rule[-\dp\strutbox]{.4pt}{\baselineskip}\hspace{2mm}\makebox[36.2mm][l]{\smash{74. poblado de carmen}}\hspace{3mm}\rule[-\dp\strutbox]{.4pt}{\baselineskip}\hspace{2mm}\makebox[8.0mm][r]{\smash{99}}\makebox[7.0mm][r]{\smash{}}\hspace{3mm}\rule[-\dp\strutbox]{.4pt}{\baselineskip}\hspace{2mm}\makebox[6.0mm][r]{\smash{25}}\makebox[6.9mm][r]{\smash{}}}
\DLine{AB01-PDF0664-body-L052}{}{\hspace*{0.0mm}\makebox[36.6mm][l]{\smash{45. cuç desuso degipto}}\hspace{3mm}\rule[-\dp\strutbox]{.4pt}{\baselineskip}\hspace{2mm}\makebox[8.0mm][r]{\smash{62}}\makebox[7.0mm][r]{\smash{}}\hspace{3mm}\rule[-\dp\strutbox]{.4pt}{\baselineskip}\hspace{2mm}\makebox[6.0mm][r]{\smash{\textit{36}}}\makebox[7.5mm][r]{\smash{}}\hspace{3mm}\rule[-\dp\strutbox]{.4pt}{\baselineskip}\hspace{1.2pt}\rule[-\dp\strutbox]{.4pt}{\baselineskip}\hspace{2mm}\makebox[36.2mm][l]{\smash{76. Iorgen}}\hspace{3mm}\rule[-\dp\strutbox]{.4pt}{\baselineskip}\hspace{2mm}\makebox[8.0mm][r]{\smash{95}}\makebox[7.0mm][r]{\smash{}}\hspace{3mm}\rule[-\dp\strutbox]{.4pt}{\baselineskip}\hspace{2mm}\makebox[6.0mm][r]{\smash{*48}}\makebox[6.9mm][r]{\smash{}}}
\DLine{AB01-PDF0664-body-L053}{}{\hspace*{0.0mm}\makebox[36.6mm][l]{\smash{46. Cuç la de dentro}}\hspace{3mm}\rule[-\dp\strutbox]{.4pt}{\baselineskip}\hspace{2mm}\makebox[8.0mm][r]{\smash{50}}\makebox[7.0mm][r]{\smash{}}\hspace{3mm}\rule[-\dp\strutbox]{.4pt}{\baselineskip}\hspace{2mm}\makebox[6.0mm][r]{\smash{*42}}\makebox[7.5mm][r]{\smash{}}\hspace{3mm}\rule[-\dp\strutbox]{.4pt}{\baselineskip}\hspace{1.2pt}\rule[-\dp\strutbox]{.4pt}{\baselineskip}\hspace{2mm}\makebox[36.2mm][l]{\smash{78. Balach}}\hspace{3mm}\rule[-\dp\strutbox]{.4pt}{\baselineskip}\hspace{2mm}\makebox[8.0mm][r]{\smash{116}}\makebox[7.0mm][r]{\smash{}}\hspace{3mm}\rule[-\dp\strutbox]{.4pt}{\baselineskip}\hspace{2mm}\makebox[6.0mm][r]{\smash{41}}\makebox[6.9mm][r]{\smash{}}}
\DLine{AB01-PDF0664-body-L054}{}{\hspace*{0.0mm}\makebox[36.6mm][l]{\smash{55. arminna la menor}}\hspace{3mm}\rule[-\dp\strutbox]{.4pt}{\baselineskip}\hspace{2mm}\makebox[8.0mm][r]{\smash{71}}\makebox[7.0mm][r]{\smash{}}\hspace{3mm}\rule[-\dp\strutbox]{.4pt}{\baselineskip}\hspace{2mm}\makebox[6.0mm][r]{\smash{39}}\makebox[7.5mm][r]{\smash{}}\hspace{3mm}\rule[-\dp\strutbox]{.4pt}{\baselineskip}\hspace{1.2pt}\rule[-\dp\strutbox]{.4pt}{\baselineskip}\hspace{2mm}\makebox[36.2mm][l]{\smash{79. açooth}}\hspace{3mm}\rule[-\dp\strutbox]{.4pt}{\baselineskip}\hspace{2mm}\makebox[8.0mm][r]{\smash{114}}\makebox[7.0mm][r]{\smash{}}\hspace{3mm}\rule[-\dp\strutbox]{.4pt}{\baselineskip}\hspace{2mm}\makebox[6.0mm][r]{\smash{45}}\makebox[6.9mm][r]{\smash{}}}
\par\vspace{4mm}
\par\fontsize{9}{12}\selectfont
\DLine{AB01-PDF0664-body-L055}{55}{17. \textit{Ysla,} quia Arabice \textit{ǵazīrah} insulam ac paeninsulam significat.}
\DLine{AB01-PDF0664-body-L056}{}{66. Minuta latit. 31′ forte quia translator Hispanus in exemplari suo invenerat hoc loco, ad 0′ designan-}
\DLine{AB01-PDF0664-body-L057}{}{\hspace*{2.5em}dum, negationem Arabicam \textarabic{لا} (sicut ter vel quater etiam in cod. Escur. fit), quam iuxta suum valo-}
\DLine{AB01-PDF0664-body-L058}{}{\hspace*{2.5em}rem numeralem accipiendam putavit.}
\DLine{AB01-PDF0664-body-L059}{}{72. Lat. 34° (ergo \textarabic{لد}), in cod. Esc. 37° \textarabic{لز}. \Name{Ptolemaeus} VI, 5 extremas latitudines 33° ac 39° statuit;}
\DLine{AB01-PDF0664-body-L060}{60}{\hspace*{2.5em}vera lectio igitur 36° \textarabic{لو}, quae p. 35 restituenda.}
\DLine{AB01-PDF0664-body-L061}{}{79. Nomen legit \textarabic{الصعد} pro \textarabic{الصغد}.}
\DSource{0665}{216}
\par\vspace{3mm}
\par\noindent\rule{157mm}{1.2pt}\par\nointerlineskip\vspace{1pt}\noindent\rule{157mm}{.4pt}\par
\par\nointerlineskip\noindent\hspace*{0.0mm}\hspace{36.6mm}\hspace{3mm}\rule{.4pt}{4mm}\hspace{2mm}\hspace{8.0mm}\hspace{7.0mm}\hspace{3mm}\rule{.4pt}{4mm}\hspace{2mm}\hspace{6.0mm}\hspace{7.5mm}\hspace{3mm}\rule{.4pt}{4mm}\hspace{1.2pt}\rule{.4pt}{4mm}\hspace{2mm}\hspace{36.2mm}\hspace{3mm}\rule{.4pt}{4mm}\hspace{2mm}\hspace{8.0mm}\hspace{7.0mm}\hspace{3mm}\rule{.4pt}{4mm}\hspace{2mm}\hspace{6.0mm}\hspace{6.9mm}\par\nointerlineskip
\DLine{AB01-PDF0665-body-L001}{}{\hspace*{0.0mm}\makebox[36.6mm][l]{\smash{81. Turcos de dentro}}\hspace{3mm}\rule[-\dp\strutbox]{.4pt}{\baselineskip}\hspace{2mm}\makebox[8.0mm][r]{120\rlap{°}}\makebox[7.0mm][r]{\smash{}}\hspace{3mm}\rule[-\dp\strutbox]{.4pt}{\baselineskip}\hspace{2mm}\makebox[6.0mm][r]{*16\rlap{°}}\makebox[7.5mm][r]{\smash{}}\hspace{3mm}\rule[-\dp\strutbox]{.4pt}{\baselineskip}\hspace{1.2pt}\rule[-\dp\strutbox]{.4pt}{\baselineskip}\hspace{2mm}\makebox[36.2mm][l]{\smash{127. acre}}\hspace{3mm}\rule[-\dp\strutbox]{.4pt}{\baselineskip}\hspace{2mm}\makebox[8.0mm][r]{66\rlap{°}}\makebox[7.0mm][r]{50\rlap{′}}\hspace{3mm}\rule[-\dp\strutbox]{.4pt}{\baselineskip}\hspace{2mm}\makebox[6.0mm][r]{33\rlap{°}}\makebox[6.9mm][r]{\smash{}}}
\DLine{AB01-PDF0665-body-L002}{}{\hspace*{0.0mm}\makebox[36.6mm][l]{\smash{82. los turcos de fuera}}\hspace{3mm}\rule[-\dp\strutbox]{.4pt}{\baselineskip}\hspace{2mm}\makebox[8.0mm][r]{\smash{*107}}\makebox[7.0mm][r]{\smash{}}\hspace{3mm}\rule[-\dp\strutbox]{.4pt}{\baselineskip}\hspace{2mm}\makebox[6.0mm][r]{\smash{*43}}\makebox[7.5mm][r]{\smash{}}\hspace{3mm}\rule[-\dp\strutbox]{.4pt}{\baselineskip}\hspace{1.2pt}\rule[-\dp\strutbox]{.4pt}{\baselineskip}\hspace{2mm}\makebox[36.2mm][l]{\smash{128. hemz}}\hspace{3mm}\rule[-\dp\strutbox]{.4pt}{\baselineskip}\hspace{2mm}\makebox[8.0mm][r]{\smash{69}}\makebox[7.0mm][r]{\smash{5}}\hspace{3mm}\rule[-\dp\strutbox]{.4pt}{\baselineskip}\hspace{2mm}\makebox[6.0mm][r]{\smash{34}}\makebox[6.9mm][r]{\smash{}}}
\DLine{AB01-PDF0665-body-L003}{}{\hspace*{0.0mm}\makebox[36.6mm][l]{\smash{83. Tabarçathen}}\hspace{3mm}\rule[-\dp\strutbox]{.4pt}{\baselineskip}\hspace{2mm}\makebox[8.0mm][r]{\smash{*162}}\makebox[7.0mm][r]{\smash{}}\hspace{3mm}\rule[-\dp\strutbox]{.4pt}{\baselineskip}\hspace{2mm}\makebox[6.0mm][r]{\smash{45}}\makebox[7.5mm][r]{\smash{}}\hspace{3mm}\rule[-\dp\strutbox]{.4pt}{\baselineskip}\hspace{1.2pt}\rule[-\dp\strutbox]{.4pt}{\baselineskip}\hspace{2mm}\makebox[36.2mm][l]{\smash{130. hama}}\hspace{3mm}\rule[-\dp\strutbox]{.4pt}{\baselineskip}\hspace{2mm}\makebox[8.0mm][r]{\smash{69}}\makebox[7.0mm][r]{\smash{30}}\hspace{3mm}\rule[-\dp\strutbox]{.4pt}{\baselineskip}\hspace{2mm}\makebox[6.0mm][r]{\smash{\textit{35}}}\makebox[6.9mm][r]{\smash{20}}}
\DLine{AB01-PDF0665-body-L004}{}{\hspace*{0.0mm}\makebox[36.6mm][l]{\smash{85. fragana}}\hspace{3mm}\rule[-\dp\strutbox]{.4pt}{\baselineskip}\hspace{2mm}\makebox[8.0mm][r]{\smash{116}}\makebox[7.0mm][r]{\smash{}}\hspace{3mm}\rule[-\dp\strutbox]{.4pt}{\baselineskip}\hspace{2mm}\makebox[6.0mm][r]{\smash{35}}\makebox[7.5mm][r]{\smash{}}\hspace{3mm}\rule[-\dp\strutbox]{.4pt}{\baselineskip}\hspace{1.2pt}\rule[-\dp\strutbox]{.4pt}{\baselineskip}\hspace{2mm}\makebox[36.2mm][l]{\smash{133. Damascho}}\hspace{3mm}\rule[-\dp\strutbox]{.4pt}{\baselineskip}\hspace{2mm}\makebox[8.0mm][r]{\smash{69}}\makebox[7.0mm][r]{\smash{}}\hspace{3mm}\rule[-\dp\strutbox]{.4pt}{\baselineskip}\hspace{2mm}\makebox[6.0mm][r]{\smash{33}}\makebox[6.9mm][r]{\smash{}}}
\DLine{AB01-PDF0665-body-L005}{5}{\hspace*{0.0mm}\makebox[36.6mm][l]{\smash{86. çarazthen}}\hspace{3mm}\rule[-\dp\strutbox]{.4pt}{\baselineskip}\hspace{2mm}\makebox[8.0mm][r]{\smash{108}}\makebox[7.0mm][r]{\smash{}}\hspace{3mm}\rule[-\dp\strutbox]{.4pt}{\baselineskip}\hspace{2mm}\makebox[6.0mm][r]{\smash{29}}\makebox[7.5mm][r]{\smash{}}\hspace{3mm}\rule[-\dp\strutbox]{.4pt}{\baselineskip}\hspace{1.2pt}\rule[-\dp\strutbox]{.4pt}{\baselineskip}\hspace{2mm}\makebox[36.2mm][l]{\smash{134. balbach}}\hspace{3mm}\rule[-\dp\strutbox]{.4pt}{\baselineskip}\hspace{2mm}\makebox[8.0mm][r]{\smash{68}}\makebox[7.0mm][r]{\smash{20}}\hspace{3mm}\rule[-\dp\strutbox]{.4pt}{\baselineskip}\hspace{2mm}\makebox[6.0mm][r]{\smash{33}}\makebox[6.9mm][r]{\smash{15}}}
\DLine{AB01-PDF0665-body-L006}{}{\hspace*{0.0mm}\makebox[36.6mm][l]{\smash{87. Tierras dazinge}}\hspace{3mm}\rule[-\dp\strutbox]{.4pt}{\baselineskip}\hspace{2mm}\makebox[8.0mm][r]{\smash{115}}\makebox[7.0mm][r]{\smash{}}\hspace{3mm}\rule[-\dp\strutbox]{.4pt}{\baselineskip}\hspace{2mm}\makebox[6.0mm][r]{\smash{29}}\makebox[7.5mm][r]{\smash{}}\hspace{3mm}\rule[-\dp\strutbox]{.4pt}{\baselineskip}\hspace{1.2pt}\rule[-\dp\strutbox]{.4pt}{\baselineskip}\hspace{2mm}\makebox[36.2mm][l]{\smash{136. halab}}\hspace{3mm}\rule[-\dp\strutbox]{.4pt}{\baselineskip}\hspace{2mm}\makebox[8.0mm][r]{\smash{71}}\makebox[7.0mm][r]{\smash{}}\hspace{3mm}\rule[-\dp\strutbox]{.4pt}{\baselineskip}\hspace{2mm}\makebox[6.0mm][r]{\smash{*35}}\makebox[6.9mm][r]{\smash{50}}}
\DLine{AB01-PDF0665-body-L007}{}{\hspace*{0.0mm}\makebox[36.6mm][l]{\smash{88. Tierras daçindi}}\hspace{3mm}\rule[-\dp\strutbox]{.4pt}{\baselineskip}\hspace{2mm}\makebox[8.0mm][r]{\smash{110}}\makebox[7.0mm][r]{\smash{}}\hspace{3mm}\rule[-\dp\strutbox]{.4pt}{\baselineskip}\hspace{2mm}\makebox[6.0mm][r]{\smash{23}}\makebox[7.5mm][r]{\smash{}}\hspace{3mm}\rule[-\dp\strutbox]{.4pt}{\baselineskip}\hspace{1.2pt}\rule[-\dp\strutbox]{.4pt}{\baselineskip}\hspace{2mm}\makebox[36.2mm][l]{\smash{138. magarat anaomen}}\hspace{3mm}\rule[-\dp\strutbox]{.4pt}{\baselineskip}\hspace{2mm}\makebox[8.0mm][r]{\smash{69}}\makebox[7.0mm][r]{\smash{55}}\hspace{3mm}\rule[-\dp\strutbox]{.4pt}{\baselineskip}\hspace{2mm}\makebox[6.0mm][r]{\smash{34}}\makebox[6.9mm][r]{\smash{}}}
\DLine{AB01-PDF0665-body-L008}{}{\hspace*{0.0mm}\makebox[36.6mm][l]{\smash{89. India dalle del rio}}\hspace{3mm}\rule[-\dp\strutbox]{.4pt}{\baselineskip}\hspace{2mm}\makebox[8.0mm][r]{\smash{132}}\makebox[7.0mm][r]{\smash{}}\hspace{3mm}\rule[-\dp\strutbox]{.4pt}{\baselineskip}\hspace{2mm}\makebox[6.0mm][r]{\smash{*26}}\makebox[7.5mm][r]{\smash{}}\hspace{3mm}\rule[-\dp\strutbox]{.4pt}{\baselineskip}\hspace{1.2pt}\rule[-\dp\strutbox]{.4pt}{\baselineskip}\hspace{2mm}\makebox[36.2mm][l]{\smash{139. coriz}}\hspace{3mm}\rule[-\dp\strutbox]{.4pt}{\baselineskip}\hspace{2mm}\makebox[8.0mm][r]{\smash{70}}\makebox[7.0mm][r]{\smash{\textit{50}}}\hspace{3mm}\rule[-\dp\strutbox]{.4pt}{\baselineskip}\hspace{2mm}\makebox[6.0mm][r]{\smash{*34}}\makebox[6.9mm][r]{\smash{20}}}
\DLine{AB01-PDF0665-body-L009}{}{\hspace*{0.0mm}\makebox[36.6mm][l]{\smash{90. India daquen del}}\hspace{3mm}\rule[-\dp\strutbox]{.4pt}{\baselineskip}\hspace{2mm}\makebox[8.0mm][r]{\smash{}}\makebox[7.0mm][r]{\smash{}}\hspace{3mm}\rule[-\dp\strutbox]{.4pt}{\baselineskip}\hspace{2mm}\makebox[6.0mm][r]{\smash{}}\makebox[7.5mm][r]{\smash{}}\hspace{3mm}\rule[-\dp\strutbox]{.4pt}{\baselineskip}\hspace{1.2pt}\rule[-\dp\strutbox]{.4pt}{\baselineskip}\hspace{2mm}\makebox[36.2mm][l]{\smash{142. Anthiochia}}\hspace{3mm}\rule[-\dp\strutbox]{.4pt}{\baselineskip}\hspace{2mm}\makebox[8.0mm][r]{\smash{69}}\makebox[7.0mm][r]{\smash{}}\hspace{3mm}\rule[-\dp\strutbox]{.4pt}{\baselineskip}\hspace{2mm}\makebox[6.0mm][r]{\smash{35}}\makebox[6.9mm][r]{\smash{30}}}
\DLine{AB01-PDF0665-body-L010}{10}{\hspace*{0.0mm}\makebox[36.6mm][l]{\smash{\hspace*{6mm}Ryo}}\hspace{3mm}\rule[-\dp\strutbox]{.4pt}{\baselineskip}\hspace{2mm}\makebox[8.0mm][r]{\smash{*180}}\makebox[7.0mm][r]{\smash{}}\hspace{3mm}\rule[-\dp\strutbox]{.4pt}{\baselineskip}\hspace{2mm}\makebox[6.0mm][r]{\smash{22}}\makebox[7.5mm][r]{\smash{}}\hspace{3mm}\rule[-\dp\strutbox]{.4pt}{\baselineskip}\hspace{1.2pt}\rule[-\dp\strutbox]{.4pt}{\baselineskip}\hspace{2mm}\makebox[36.2mm][l]{\smash{143. malthya}}\hspace{3mm}\rule[-\dp\strutbox]{.4pt}{\baselineskip}\hspace{2mm}\makebox[8.0mm][r]{\smash{71}}\makebox[7.0mm][r]{\smash{}}\hspace{3mm}\rule[-\dp\strutbox]{.4pt}{\baselineskip}\hspace{2mm}\makebox[6.0mm][r]{\smash{39}}\makebox[6.9mm][r]{\smash{}}}
\DLine{AB01-PDF0665-body-L011}{}{\hspace*{0.0mm}\makebox[36.6mm][l]{\smash{91. Isla de çarandib}}\hspace{3mm}\rule[-\dp\strutbox]{.4pt}{\baselineskip}\hspace{2mm}\makebox[8.0mm][r]{\smash{124}}\makebox[7.0mm][r]{\smash{}}\hspace{3mm}\rule[-\dp\strutbox]{.4pt}{\baselineskip}\hspace{2mm}\makebox[6.0mm][r]{\smash{3}}\makebox[7.5mm][r]{\smash{}}\hspace{3mm}\rule[-\dp\strutbox]{.4pt}{\baselineskip}\hspace{1.2pt}\rule[-\dp\strutbox]{.4pt}{\baselineskip}\hspace{2mm}\makebox[36.2mm][l]{\smash{144. Simxath}}\hspace{3mm}\rule[-\dp\strutbox]{.4pt}{\baselineskip}\hspace{2mm}\makebox[8.0mm][r]{\smash{73}}\makebox[7.0mm][r]{\smash{20}}\hspace{3mm}\rule[-\dp\strutbox]{.4pt}{\baselineskip}\hspace{2mm}\makebox[6.0mm][r]{\smash{38}}\makebox[6.9mm][r]{\smash{40}}}
\DLine{AB01-PDF0665-body-L012}{}{\hspace*{0.0mm}\makebox[36.6mm][l]{\smash{92. comedio de he-}}\hspace{3mm}\rule[-\dp\strutbox]{.4pt}{\baselineskip}\hspace{2mm}\makebox[8.0mm][r]{\smash{}}\makebox[7.0mm][r]{\smash{}}\hspace{3mm}\rule[-\dp\strutbox]{.4pt}{\baselineskip}\hspace{2mm}\makebox[6.0mm][r]{\smash{}}\makebox[7.5mm][r]{\smash{}}\hspace{3mm}\rule[-\dp\strutbox]{.4pt}{\baselineskip}\hspace{1.2pt}\rule[-\dp\strutbox]{.4pt}{\baselineskip}\hspace{2mm}\makebox[36.2mm][l]{\smash{150. Arraca}}\hspace{3mm}\rule[-\dp\strutbox]{.4pt}{\baselineskip}\hspace{2mm}\makebox[8.0mm][r]{\smash{73}}\makebox[7.0mm][r]{\smash{15}}\hspace{3mm}\rule[-\dp\strutbox]{.4pt}{\baselineskip}\hspace{2mm}\makebox[6.0mm][r]{\smash{36}}\makebox[6.9mm][r]{\smash{*47}}}
\DLine{AB01-PDF0665-body-L013}{}{\hspace*{0.0mm}\makebox[36.6mm][l]{\smash{\hspace*{6mm}myar}}\hspace{3mm}\rule[-\dp\strutbox]{.4pt}{\baselineskip}\hspace{2mm}\makebox[8.0mm][r]{\smash{\textit{102}}}\makebox[7.0mm][r]{\smash{}}\hspace{3mm}\rule[-\dp\strutbox]{.4pt}{\baselineskip}\hspace{2mm}\makebox[6.0mm][r]{\smash{12}}\makebox[7.5mm][r]{\smash{30}}\hspace{3mm}\rule[-\dp\strutbox]{.4pt}{\baselineskip}\hspace{1.2pt}\rule[-\dp\strutbox]{.4pt}{\baselineskip}\hspace{2mm}\makebox[36.2mm][l]{\smash{152. harran}}\hspace{3mm}\rule[-\dp\strutbox]{.4pt}{\baselineskip}\hspace{2mm}\makebox[8.0mm][r]{\smash{73}}\makebox[7.0mm][r]{\smash{}}\hspace{3mm}\rule[-\dp\strutbox]{.4pt}{\baselineskip}\hspace{2mm}\makebox[6.0mm][r]{\smash{36}}\makebox[6.9mm][r]{\smash{40}}}
\DLine{AB01-PDF0665-body-L014}{}{\hspace*{0.0mm}\makebox[36.6mm][l]{\smash{93. Tierras daçin}}\hspace{3mm}\rule[-\dp\strutbox]{.4pt}{\baselineskip}\hspace{2mm}\makebox[8.0mm][r]{\smash{177}}\makebox[7.0mm][r]{\smash{}}\hspace{3mm}\rule[-\dp\strutbox]{.4pt}{\baselineskip}\hspace{2mm}\makebox[6.0mm][r]{\smash{22}}\makebox[7.5mm][r]{\smash{}}\hspace{3mm}\rule[-\dp\strutbox]{.4pt}{\baselineskip}\hspace{1.2pt}\rule[-\dp\strutbox]{.4pt}{\baselineskip}\hspace{2mm}\makebox[36.2mm][l]{\smash{153. azuhe}}\hspace{3mm}\rule[-\dp\strutbox]{.4pt}{\baselineskip}\hspace{2mm}\makebox[8.0mm][r]{\smash{72}}\makebox[7.0mm][r]{\smash{*}}\hspace{3mm}\rule[-\dp\strutbox]{.4pt}{\baselineskip}\hspace{2mm}\makebox[6.0mm][r]{\smash{*50}}\makebox[6.9mm][r]{\smash{*37}}}
\DLine{AB01-PDF0665-body-L015}{15}{\hspace*{0.0mm}\makebox[36.6mm][l]{\smash{102. Hara}}\hspace{3mm}\rule[-\dp\strutbox]{.4pt}{\baselineskip}\hspace{2mm}\makebox[8.0mm][r]{\smash{43}}\makebox[7.0mm][r]{\smash{}}\hspace{3mm}\rule[-\dp\strutbox]{.4pt}{\baselineskip}\hspace{2mm}\makebox[6.0mm][r]{\smash{23}}\makebox[7.5mm][r]{\smash{30}}\hspace{3mm}\rule[-\dp\strutbox]{.4pt}{\baselineskip}\hspace{1.2pt}\rule[-\dp\strutbox]{.4pt}{\baselineskip}\hspace{2mm}\makebox[36.2mm][l]{\smash{156. Raçaalahyn}}\hspace{3mm}\rule[-\dp\strutbox]{.4pt}{\baselineskip}\hspace{2mm}\makebox[8.0mm][r]{\smash{74}}\makebox[7.0mm][r]{\smash{}}\hspace{3mm}\rule[-\dp\strutbox]{.4pt}{\baselineskip}\hspace{2mm}\makebox[6.0mm][r]{\smash{36}}\makebox[6.9mm][r]{\smash{50}}}
\DLine{AB01-PDF0665-body-L016}{}{\hspace*{0.0mm}\makebox[36.6mm][l]{\smash{103. mecha}}\hspace{3mm}\rule[-\dp\strutbox]{.4pt}{\baselineskip}\hspace{2mm}\makebox[8.0mm][r]{\smash{71}}\makebox[7.0mm][r]{\smash{}}\hspace{3mm}\rule[-\dp\strutbox]{.4pt}{\baselineskip}\hspace{2mm}\makebox[6.0mm][r]{\smash{21}}\makebox[7.5mm][r]{\smash{40}}\hspace{3mm}\rule[-\dp\strutbox]{.4pt}{\baselineskip}\hspace{1.2pt}\rule[-\dp\strutbox]{.4pt}{\baselineskip}\hspace{2mm}\makebox[36.2mm][l]{\smash{157. cafartota}}\hspace{3mm}\rule[-\dp\strutbox]{.4pt}{\baselineskip}\hspace{2mm}\makebox[8.0mm][r]{\smash{74}}\makebox[7.0mm][r]{\smash{35}}\hspace{3mm}\rule[-\dp\strutbox]{.4pt}{\baselineskip}\hspace{2mm}\makebox[6.0mm][r]{\smash{37}}\makebox[6.9mm][r]{\smash{5}}}
\DLine{AB01-PDF0665-body-L017}{}{\hspace*{0.0mm}\makebox[36.6mm][l]{\smash{109. alexandria degipto}}\hspace{3mm}\rule[-\dp\strutbox]{.4pt}{\baselineskip}\hspace{2mm}\makebox[8.0mm][r]{\smash{60}}\makebox[7.0mm][r]{\smash{30}}\hspace{3mm}\rule[-\dp\strutbox]{.4pt}{\baselineskip}\hspace{2mm}\makebox[6.0mm][r]{\smash{30}}\makebox[7.5mm][r]{\smash{58}}\hspace{3mm}\rule[-\dp\strutbox]{.4pt}{\baselineskip}\hspace{1.2pt}\rule[-\dp\strutbox]{.4pt}{\baselineskip}\hspace{2mm}\makebox[36.2mm][l]{\smash{159. Dara}}\hspace{3mm}\rule[-\dp\strutbox]{.4pt}{\baselineskip}\hspace{2mm}\makebox[8.0mm][r]{\smash{75}}\makebox[7.0mm][r]{\smash{15}}\hspace{3mm}\rule[-\dp\strutbox]{.4pt}{\baselineskip}\hspace{2mm}\makebox[6.0mm][r]{\smash{37}}\makebox[6.9mm][r]{\smash{10}}}
\DLine{AB01-PDF0665-body-L018}{}{\hspace*{0.0mm}\makebox[36.6mm][l]{\smash{104. yaçrib de moho-}}\hspace{3mm}\rule[-\dp\strutbox]{.4pt}{\baselineskip}\hspace{2mm}\makebox[8.0mm][r]{\smash{}}\makebox[7.0mm][r]{\smash{}}\hspace{3mm}\rule[-\dp\strutbox]{.4pt}{\baselineskip}\hspace{2mm}\makebox[6.0mm][r]{\smash{}}\makebox[7.5mm][r]{\smash{}}\hspace{3mm}\rule[-\dp\strutbox]{.4pt}{\baselineskip}\hspace{1.2pt}\rule[-\dp\strutbox]{.4pt}{\baselineskip}\hspace{2mm}\makebox[36.2mm][l]{\smash{160. maraditzh}}\hspace{3mm}\rule[-\dp\strutbox]{.4pt}{\baselineskip}\hspace{2mm}\makebox[8.0mm][r]{\smash{75}}\makebox[7.0mm][r]{\smash{}}\hspace{3mm}\rule[-\dp\strutbox]{.4pt}{\baselineskip}\hspace{2mm}\makebox[6.0mm][r]{\smash{37}}\makebox[6.9mm][r]{\smash{15}}}
\DLine{AB01-PDF0665-body-L019}{}{\hspace*{0.0mm}\makebox[36.6mm][l]{\smash{\hspace*{6mm}mat}}\hspace{3mm}\rule[-\dp\strutbox]{.4pt}{\baselineskip}\hspace{2mm}\makebox[8.0mm][r]{\smash{75}}\makebox[7.0mm][r]{\smash{}}\hspace{3mm}\rule[-\dp\strutbox]{.4pt}{\baselineskip}\hspace{2mm}\makebox[6.0mm][r]{\smash{25}}\makebox[7.5mm][r]{\smash{}}\hspace{3mm}\rule[-\dp\strutbox]{.4pt}{\baselineskip}\hspace{1.2pt}\rule[-\dp\strutbox]{.4pt}{\baselineskip}\hspace{2mm}\makebox[36.2mm][l]{\smash{162. almauçil}}\hspace{3mm}\rule[-\dp\strutbox]{.4pt}{\baselineskip}\hspace{2mm}\makebox[8.0mm][r]{\smash{78}}\makebox[7.0mm][r]{\smash{10}}\hspace{3mm}\rule[-\dp\strutbox]{.4pt}{\baselineskip}\hspace{2mm}\makebox[6.0mm][r]{\smash{36}}\makebox[6.9mm][r]{\smash{30}}}
\DLine{AB01-PDF0665-body-L020}{20}{\hspace*{0.0mm}\makebox[36.6mm][l]{\smash{111. falaztyn}}\hspace{3mm}\rule[-\dp\strutbox]{.4pt}{\baselineskip}\hspace{2mm}\makebox[8.0mm][r]{\smash{66}}\makebox[7.0mm][r]{\smash{15}}\hspace{3mm}\rule[-\dp\strutbox]{.4pt}{\baselineskip}\hspace{2mm}\makebox[6.0mm][r]{\smash{32}}\makebox[7.5mm][r]{\smash{30}}\hspace{3mm}\rule[-\dp\strutbox]{.4pt}{\baselineskip}\hspace{1.2pt}\rule[-\dp\strutbox]{.4pt}{\baselineskip}\hspace{2mm}\makebox[36.2mm][l]{\smash{163. çiniar}}\hspace{3mm}\rule[-\dp\strutbox]{.4pt}{\baselineskip}\hspace{2mm}\makebox[8.0mm][r]{\smash{77}}\makebox[7.0mm][r]{\smash{30}}\hspace{3mm}\rule[-\dp\strutbox]{.4pt}{\baselineskip}\hspace{2mm}\makebox[6.0mm][r]{\smash{36}}\makebox[6.9mm][r]{\smash{}}}
\DLine{AB01-PDF0665-body-L021}{}{\hspace*{0.0mm}\makebox[36.6mm][l]{\smash{113. azcalen}}\hspace{3mm}\rule[-\dp\strutbox]{.4pt}{\baselineskip}\hspace{2mm}\makebox[8.0mm][r]{\smash{65}}\makebox[7.0mm][r]{\smash{}}\hspace{3mm}\rule[-\dp\strutbox]{.4pt}{\baselineskip}\hspace{2mm}\makebox[6.0mm][r]{\smash{31}}\makebox[7.5mm][r]{\smash{50}}\hspace{3mm}\rule[-\dp\strutbox]{.4pt}{\baselineskip}\hspace{1.2pt}\rule[-\dp\strutbox]{.4pt}{\baselineskip}\hspace{2mm}\makebox[36.2mm][l]{\smash{164. hylach}}\hspace{3mm}\rule[-\dp\strutbox]{.4pt}{\baselineskip}\hspace{2mm}\makebox[8.0mm][r]{\smash{78}}\makebox[7.0mm][r]{\smash{}}\hspace{3mm}\rule[-\dp\strutbox]{.4pt}{\baselineskip}\hspace{2mm}\makebox[6.0mm][r]{\smash{39}}\makebox[6.9mm][r]{\smash{20}}}
\DLine{AB01-PDF0665-body-L022}{}{\hspace*{0.0mm}\makebox[36.6mm][l]{\smash{115. arrambla}}\hspace{3mm}\rule[-\dp\strutbox]{.4pt}{\baselineskip}\hspace{2mm}\makebox[8.0mm][r]{\smash{65}}\makebox[7.0mm][r]{\smash{*7}}\hspace{3mm}\rule[-\dp\strutbox]{.4pt}{\baselineskip}\hspace{2mm}\makebox[6.0mm][r]{\smash{31}}\makebox[7.5mm][r]{\smash{\textit{35}}}\hspace{3mm}\rule[-\dp\strutbox]{.4pt}{\baselineskip}\hspace{1.2pt}\rule[-\dp\strutbox]{.4pt}{\baselineskip}\hspace{2mm}\makebox[36.2mm][l]{\smash{167. bardaha}}\hspace{3mm}\rule[-\dp\strutbox]{.4pt}{\baselineskip}\hspace{2mm}\makebox[8.0mm][r]{\smash{*104}}\makebox[7.0mm][r]{\smash{}}\hspace{3mm}\rule[-\dp\strutbox]{.4pt}{\baselineskip}\hspace{2mm}\makebox[6.0mm][r]{\smash{42}}\makebox[6.9mm][r]{\smash{}}}
\DLine{AB01-PDF0665-body-L023}{}{\hspace*{0.0mm}\makebox[36.6mm][l]{\smash{117. ysla de Rodyz}}\hspace{3mm}\rule[-\dp\strutbox]{.4pt}{\baselineskip}\hspace{2mm}\makebox[8.0mm][r]{\smash{58}}\makebox[7.0mm][r]{\smash{40}}\hspace{3mm}\rule[-\dp\strutbox]{.4pt}{\baselineskip}\hspace{2mm}\makebox[6.0mm][r]{\smash{36}}\makebox[7.5mm][r]{\smash{}}\hspace{3mm}\rule[-\dp\strutbox]{.4pt}{\baselineskip}\hspace{1.2pt}\rule[-\dp\strutbox]{.4pt}{\baselineskip}\hspace{2mm}\makebox[36.2mm][l]{\smash{168. baldach}}\hspace{3mm}\rule[-\dp\strutbox]{.4pt}{\baselineskip}\hspace{2mm}\makebox[8.0mm][r]{\smash{80}}\makebox[7.0mm][r]{\smash{}}\hspace{3mm}\rule[-\dp\strutbox]{.4pt}{\baselineskip}\hspace{2mm}\makebox[6.0mm][r]{\smash{33}}\makebox[6.9mm][r]{\smash{9}}}
\DLine{AB01-PDF0665-body-L024}{}{\hspace*{0.0mm}\makebox[36.6mm][l]{\smash{118. Çalmaoz}}\hspace{3mm}\rule[-\dp\strutbox]{.4pt}{\baselineskip}\hspace{2mm}\makebox[8.0mm][r]{\smash{*46}}\makebox[7.0mm][r]{\smash{45}}\hspace{3mm}\rule[-\dp\strutbox]{.4pt}{\baselineskip}\hspace{2mm}\makebox[6.0mm][r]{\smash{35}}\makebox[7.5mm][r]{\smash{*50}}\hspace{3mm}\rule[-\dp\strutbox]{.4pt}{\baselineskip}\hspace{1.2pt}\rule[-\dp\strutbox]{.4pt}{\baselineskip}\hspace{2mm}\makebox[36.2mm][l]{\smash{169. curramanrae}}\hspace{3mm}\rule[-\dp\strutbox]{.4pt}{\baselineskip}\hspace{2mm}\makebox[8.0mm][r]{\smash{79}}\makebox[7.0mm][r]{\smash{45}}\hspace{3mm}\rule[-\dp\strutbox]{.4pt}{\baselineskip}\hspace{2mm}\makebox[6.0mm][r]{\smash{34}}\makebox[6.9mm][r]{\smash{}}}
\DLine{AB01-PDF0665-body-L025}{25}{\hspace*{0.0mm}\makebox[36.6mm][l]{\smash{119. Tarçoz}}\hspace{3mm}\rule[-\dp\strutbox]{.4pt}{\baselineskip}\hspace{2mm}\makebox[8.0mm][r]{\smash{*57}}\makebox[7.0mm][r]{\smash{40}}\hspace{3mm}\rule[-\dp\strutbox]{.4pt}{\baselineskip}\hspace{2mm}\makebox[6.0mm][r]{\smash{36}}\makebox[7.5mm][r]{\smash{*50}}\hspace{3mm}\rule[-\dp\strutbox]{.4pt}{\baselineskip}\hspace{1.2pt}\rule[-\dp\strutbox]{.4pt}{\baselineskip}\hspace{2mm}\makebox[36.2mm][l]{\smash{170. Alcufa}}\hspace{3mm}\rule[-\dp\strutbox]{.4pt}{\baselineskip}\hspace{2mm}\makebox[8.0mm][r]{\smash{79}}\makebox[7.0mm][r]{\smash{30}}\hspace{3mm}\rule[-\dp\strutbox]{.4pt}{\baselineskip}\hspace{2mm}\makebox[6.0mm][r]{\smash{31}}\makebox[6.9mm][r]{\smash{30}}}
\DLine{AB01-PDF0665-body-L026}{}{\hspace*{0.0mm}\makebox[36.6mm][l]{\smash{120. Edane}}\hspace{3mm}\rule[-\dp\strutbox]{.4pt}{\baselineskip}\hspace{2mm}\makebox[8.0mm][r]{\smash{68}}\makebox[7.0mm][r]{\smash{15}}\hspace{3mm}\rule[-\dp\strutbox]{.4pt}{\baselineskip}\hspace{2mm}\makebox[6.0mm][r]{\smash{36}}\makebox[7.5mm][r]{\smash{50}}\hspace{3mm}\rule[-\dp\strutbox]{.4pt}{\baselineskip}\hspace{1.2pt}\rule[-\dp\strutbox]{.4pt}{\baselineskip}\hspace{2mm}\makebox[36.2mm][l]{\smash{171. babillonia la anti-}}\hspace{3mm}\rule[-\dp\strutbox]{.4pt}{\baselineskip}\hspace{2mm}\makebox[8.0mm][r]{\smash{}}\makebox[7.0mm][r]{\smash{}}\hspace{3mm}\rule[-\dp\strutbox]{.4pt}{\baselineskip}\hspace{2mm}\makebox[6.0mm][r]{\smash{}}\makebox[6.9mm][r]{\smash{}}}
\DLine{AB01-PDF0665-body-L027}{}{\hspace*{0.0mm}\makebox[36.6mm][l]{\smash{121. Almaçeyça}}\hspace{3mm}\rule[-\dp\strutbox]{.4pt}{\baselineskip}\hspace{2mm}\makebox[8.0mm][r]{\smash{\textit{67}}}\makebox[7.0mm][r]{\smash{50}}\hspace{3mm}\rule[-\dp\strutbox]{.4pt}{\baselineskip}\hspace{2mm}\makebox[6.0mm][r]{\smash{36}}\makebox[7.5mm][r]{\smash{45}}\hspace{3mm}\rule[-\dp\strutbox]{.4pt}{\baselineskip}\hspace{1.2pt}\rule[-\dp\strutbox]{.4pt}{\baselineskip}\hspace{2mm}\makebox[36.2mm][l]{\smash{\hspace*{6mm}gua}}\hspace{3mm}\rule[-\dp\strutbox]{.4pt}{\baselineskip}\hspace{2mm}\makebox[8.0mm][r]{\smash{79}}\makebox[7.0mm][r]{\smash{}}\hspace{3mm}\rule[-\dp\strutbox]{.4pt}{\baselineskip}\hspace{2mm}\makebox[6.0mm][r]{\smash{35}}\makebox[6.9mm][r]{\smash{}}}
\DLine{AB01-PDF0665-body-L028}{}{\hspace*{0.0mm}\makebox[36.6mm][l]{\smash{122. aladaquia}}\hspace{3mm}\rule[-\dp\strutbox]{.4pt}{\baselineskip}\hspace{2mm}\makebox[8.0mm][r]{\smash{68}}\makebox[7.0mm][r]{\smash{30}}\hspace{3mm}\rule[-\dp\strutbox]{.4pt}{\baselineskip}\hspace{2mm}\makebox[6.0mm][r]{\smash{35}}\makebox[7.5mm][r]{\smash{5}}\hspace{3mm}\rule[-\dp\strutbox]{.4pt}{\baselineskip}\hspace{1.2pt}\rule[-\dp\strutbox]{.4pt}{\baselineskip}\hspace{2mm}\makebox[36.2mm][l]{\smash{172. Arryo}}\hspace{3mm}\rule[-\dp\strutbox]{.4pt}{\baselineskip}\hspace{2mm}\makebox[8.0mm][r]{\smash{*96}}\makebox[7.0mm][r]{\smash{}}\hspace{3mm}\rule[-\dp\strutbox]{.4pt}{\baselineskip}\hspace{2mm}\makebox[6.0mm][r]{\smash{*37}}\makebox[6.9mm][r]{\smash{30}}}
\DLine{AB01-PDF0665-body-L029}{}{\hspace*{0.0mm}\makebox[36.6mm][l]{\smash{123. Tripolo}}\hspace{3mm}\rule[-\dp\strutbox]{.4pt}{\baselineskip}\hspace{2mm}\makebox[8.0mm][r]{\smash{67}}\makebox[7.0mm][r]{\smash{30}}\hspace{3mm}\rule[-\dp\strutbox]{.4pt}{\baselineskip}\hspace{2mm}\makebox[6.0mm][r]{\smash{34}}\makebox[7.5mm][r]{\smash{20}}\hspace{3mm}\rule[-\dp\strutbox]{.4pt}{\baselineskip}\hspace{1.2pt}\rule[-\dp\strutbox]{.4pt}{\baselineskip}\hspace{2mm}\makebox[36.2mm][l]{\smash{173. Niniue}}\hspace{3mm}\rule[-\dp\strutbox]{.4pt}{\baselineskip}\hspace{2mm}\makebox[8.0mm][r]{\smash{78}}\makebox[7.0mm][r]{\smash{}}\hspace{3mm}\rule[-\dp\strutbox]{.4pt}{\baselineskip}\hspace{2mm}\makebox[6.0mm][r]{\smash{36}}\makebox[6.9mm][r]{\smash{*49}}}
\DLine{AB01-PDF0665-body-L030}{30}{\hspace*{0.0mm}\makebox[36.6mm][l]{\smash{124. Tarcha}}\hspace{3mm}\rule[-\dp\strutbox]{.4pt}{\baselineskip}\hspace{2mm}\makebox[8.0mm][r]{\smash{68}}\makebox[7.0mm][r]{\smash{30}}\hspace{3mm}\rule[-\dp\strutbox]{.4pt}{\baselineskip}\hspace{2mm}\makebox[6.0mm][r]{\smash{34}}\makebox[7.5mm][r]{\smash{}}\hspace{3mm}\rule[-\dp\strutbox]{.4pt}{\baselineskip}\hspace{1.2pt}\rule[-\dp\strutbox]{.4pt}{\baselineskip}\hspace{2mm}\makebox[36.2mm][l]{\smash{174. Albaçra}}\hspace{3mm}\rule[-\dp\strutbox]{.4pt}{\baselineskip}\hspace{2mm}\makebox[8.0mm][r]{\smash{*80}}\makebox[7.0mm][r]{\smash{*10}}\hspace{3mm}\rule[-\dp\strutbox]{.4pt}{\baselineskip}\hspace{2mm}\makebox[6.0mm][r]{\smash{31}}\makebox[6.9mm][r]{\smash{}}}
\DLine{AB01-PDF0665-body-L031}{}{\hspace*{0.0mm}\makebox[36.6mm][l]{\smash{125. Çor}}\hspace{3mm}\rule[-\dp\strutbox]{.4pt}{\baselineskip}\hspace{2mm}\makebox[8.0mm][r]{\smash{67}}\makebox[7.0mm][r]{\smash{}}\hspace{3mm}\rule[-\dp\strutbox]{.4pt}{\baselineskip}\hspace{2mm}\makebox[6.0mm][r]{\smash{33}}\makebox[7.5mm][r]{\smash{20}}\hspace{3mm}\rule[-\dp\strutbox]{.4pt}{\baselineskip}\hspace{1.2pt}\rule[-\dp\strutbox]{.4pt}{\baselineskip}\hspace{2mm}\makebox[36.2mm][l]{\smash{175. cayraf}}\hspace{3mm}\rule[-\dp\strutbox]{.4pt}{\baselineskip}\hspace{2mm}\makebox[8.0mm][r]{\smash{*109}}\makebox[7.0mm][r]{\smash{30}}\hspace{3mm}\rule[-\dp\strutbox]{.4pt}{\baselineskip}\hspace{2mm}\makebox[6.0mm][r]{\smash{29}}\makebox[6.9mm][r]{\smash{30}}}
\DLine{AB01-PDF0665-body-L032}{}{\hspace*{0.0mm}\makebox[36.6mm][l]{\smash{126. Çayde}}\hspace{3mm}\rule[-\dp\strutbox]{.4pt}{\baselineskip}\hspace{2mm}\makebox[8.0mm][r]{\smash{*66}}\makebox[7.0mm][r]{\smash{20}}\hspace{3mm}\rule[-\dp\strutbox]{.4pt}{\baselineskip}\hspace{2mm}\makebox[6.0mm][r]{\smash{33}}\makebox[7.5mm][r]{\smash{30}}\hspace{3mm}\rule[-\dp\strutbox]{.4pt}{\baselineskip}\hspace{1.2pt}\rule[-\dp\strutbox]{.4pt}{\baselineskip}\hspace{2mm}\makebox[36.2mm][l]{\smash{176. Vaçith}}\hspace{3mm}\rule[-\dp\strutbox]{.4pt}{\baselineskip}\hspace{2mm}\makebox[8.0mm][r]{\smash{*101}}\makebox[7.0mm][r]{\smash{30}}\hspace{3mm}\rule[-\dp\strutbox]{.4pt}{\baselineskip}\hspace{2mm}\makebox[6.0mm][r]{\smash{*34}}\makebox[6.9mm][r]{\smash{30}}}
\par\vspace{6mm}
\par\fontsize{9}{12}\selectfont
\DLine{AB01-PDF0665-body-L033}{}{87. Nomen legit \textarabic{الزنج} \textit{az-Zinǵ} pro \textarabic{الرخج} \textit{ar-Rukhkhaǵ.}}
\DLine{AB01-PDF0665-body-L034}{}{102. Nomen incertum videtur \textarabic{حارا} legisse.}
\DLine{AB01-PDF0665-body-L035}{35}{103. Omittit longitudinem urbis probatam.}
\DLine{AB01-PDF0665-body-L036}{}{109. Minuta latitudinis 58′ \textarabic{نح} longe meliora quam 18′ \textarabic{يح} codicis Escurialensis, a me p. 38 haud recte servata.}
\DLine{AB01-PDF0665-body-L037}{}{118. Nomen sicut in cod. Escur. scriptum est invenit.}
\DLine{AB01-PDF0665-body-L038}{}{124. Nominis corruptela amanuensi Hispano tribuenda videtur.}
\DLine{AB01-PDF0665-body-L039}{}{153. Nomen legit \textarabic{الزها} pro \textarabic{الرها}; numeri false ab amanuensi quodam collocati pro 72° 50′, 37° 0′.}
\DLine{AB01-PDF0665-body-L040}{40}{164. Nominis corruptela (pro \textit{hylath}) amanuensi Hispano tribuenda.}
\DLine{AB01-PDF0665-body-L041}{}{169. Minuta long. ut apud \Name{al-Khuwārizmī}; minus recte cod. Escur. 10′, a me pag. 42 servata.}
\DLine{AB01-PDF0665-body-L042}{}{175. Minuta latit. iustiora quam 29′ cod. Escur., a me pag. 43 servata et ex numero graduum false re-}
\DLine{AB01-PDF0665-body-L043}{}{\hspace*{2.5em}petito orientia.}
\DLine{AB01-PDF0665-body-L044}{}{176. Gradus latit. emendationem meam pag. 43 confirmant; cod. Escur. 30° \textarabic{ل}, interpres Hispanus 34° \textarabic{لد},}
\DLine{AB01-PDF0665-body-L045}{45}{\hspace*{2.5em}vera lectio 32° \textarabic{لٮ} (= \textarabic{لب}).}
\end{document}
